% ==================================================================
%  徐森林、金亚东、薛春华《数学分析》（第三册）
%  第 15 章  含参变量积分
%  源：原书印刷页 177-243（PDF 物理页 = 印刷页 + 6）
% ==================================================================

\chapter{含参变量积分}\label{ch:15}

在第 13 章与第 14 章已经看到, 无穷级数是构造新函数的一种重要工具, 用它构造了处处连续处处不可导的函数, 还用它构造了能填满正方形的连续曲线. 和无穷级数具有同样重要意义的另一种构造新函数的工具是含参变量的积分.

考虑闭区间 $[a,b]\times[\alpha,\beta]$ 上的二元连续函数 $f(x,u)$, 则对固定的 $u\in[\alpha,\beta]$, 函数 $f(x,u)$ 对变量 $x$ 在 $[a,b]$ 上 Riemann 可积, 称积分
\[
  \int_{a}^{b}f(x,u)\mathrm{d}x
\]
是含参变量 $u$ 的正常 (常义) 积分. 如果对于固定的 $u$, $f(x,u)$ 是变量 $x$ 在 $[a,b]$ 中的无界函数, 或者 $[a,b]$ 是无穷区间, 则称相应的积分是含参变量 $u$ 的反常 (广义) 积分.

本章就是要讨论由含参变量的常义积分与广义积分所确定的函数的分析性质, 即它们的连续性、可导性, 以及如何计算它们的导数与积分. 其起关键作用的是要引进参变量广义积分的一致收敛性. 最后, 我们还研究了两个重要的特殊函数——$\Gamma$ 函数与 $B$ 函数, 附带还给出了余元公式与 Stirling 公式.

% ==================================================================
%  徐森林、金亚东、薛春华《数学分析》（第三册）
%  S15.1  含参变量的正常积分
%  源：原书印刷页 177-183（PDF 物理页 183-189）
%  底本：PaddleOCR-VL-1.6，见 chapters/_ch15_15_1_raw.md
%  约定：全西文标点；行间公式用一对方括号；多行公式用 align/aligned；
%        跨册引用一律 \extref；数学模式内的中文必须包 \text{}
%  已订正的 OCR 错误与原书照录清单见 tools/_make_ch15_1.py 头部说明
% ==================================================================
\section{含参变量的正常积分}\label{sec:15.1}

应用函数项级数 $\sum_{n=1}^{\infty}u_{n}(x)$ 可以构造新函数, 它就是这个函数项级数的和函数. 收敛的幂级数是构造新函数的最简单而且又是非常重要的一种方法. 第 16 章将用三角级数, 或者更确切地用 Fourier 级数
\[
  \frac{a_{0}}{2}+\sum_{n=1}^{\infty}[a_{n}\cos nx+b_{n}\sin nx]
\]
来构造新函数.

现在我们用含参变量 (的正常或反常 (广义)) 积分来构造新函数.

\begin{definition}\label{def:15.1.1}
设二元函数 $f(x,u)$ 在闭区间 $[a,b]\times[\alpha,\beta]$ 上连续, 我们要研究由含参变量 $u$ 的正常积分所确定的关于 $u$ 的函数
\[
  \varphi(u)=\int_{a}^{b}f(x,u)\mathrm{d}x
\]
的分析性质，并称它为含参变量 $u$ 的正常积分.
\end{definition}

\begin{theorem}[含参变量正常积分保持连续性]\label{thm:15.1.1}
如果 $f(x,u)$ 在闭区间 $[a,b]\times[\alpha,\beta]$ 上连续, 则
\[
  \varphi(u)=\int_{a}^{b}f(x,u)\mathrm{d}x
\]
为区间 $[\alpha,\beta]$ 上的连续函数，即 $\forall u_{0}\in[\alpha,\beta]$，有
\[
  \lim_{u\to u_{0}}\int_{a}^{b}f(x,u)\mathrm{d}x=\lim_{u\to u_{0}}\varphi(u)=\varphi(u_{0})=\int_{a}^{b}f(x,u_{0})\mathrm{d}x=\int_{a}^{b}\lim_{u\to u_{0}}f(x,u)\mathrm{d}x
\]
(极限号与积分号可交换或可积分号下求极限).
\end{theorem}

\begin{proof}
设 $u_{0}\in[\alpha,\beta]$. 因为 $f(x,u)$ 在 $[a,b]\times[\alpha,\beta]$ 上连续, 所以它一致连续, 故 $\forall\varepsilon>0$, $\exists\delta>0$, 当 $(x_{1},u_{1})$, $(x_{2},u_{2})\in[a,b]\times[\alpha,\beta]$ 时, 只要
\[
  \rho_{0}\big((x_{1},u_{1}),(x_{2},u_{2})\big)=\sqrt{(x_{1}-x_{2})^{2}+(u_{1}-u_{2})^{2}}<\delta,
\]
就有
\[
  |f(x_{1},u_{1})-f(x_{2},u_{2})|<\frac{\varepsilon}{b-a}.
\]
因此, 当 $|u-u_{0}|<\delta$ 时, $|f(x,u)-f(x,u_{0})|<\frac{\varepsilon}{b-a}$. 于是
\[
  \begin{aligned}
    |\varphi(u)-\varphi(u_{0})|&=\left|\int_{a}^{b}f(x,u)\mathrm{d}x-\int_{a}^{b}f(x,u_{0})\mathrm{d}x\right|\\
    &\leqslant\int_{a}^{b}|f(x,u)-f(x,u_{0})|\mathrm{d}x<\frac{\varepsilon}{b-a}\int_{a}^{b}\mathrm{d}x=\varepsilon,
  \end{aligned}
\]
故 $\varphi(u)$ 在 $u_{0}$ 处连续. 由 $u_{0}\in[\alpha,\beta]$ 是任取的, 故 $\varphi(u)$ 在 $[\alpha,\beta]$ 上连续.
\end{proof}

\begin{theorem}[积分号可交换]\label{thm:15.1.2}
如果 $f(x,u)$ 在闭区间 $[a,b]\times[\alpha,\beta]$ 上连续, 则
\[
  \varphi(u)=\int_{a}^{b}f(x,u)\mathrm{d}x
\]
在 $[\alpha,\beta]$ 上连续 (当然 Riemann 可积), 且
\[
  \begin{aligned}
    \int_{\alpha}^{\beta}\varphi(u)\mathrm{d}u&=\int_{\alpha}^{\beta}\left[\int_{a}^{b}f(x,u)\mathrm{d}x\right]\mathrm{d}u\\
    &=\int_{a}^{b}\left[\int_{\alpha}^{\beta}f(x,u)\mathrm{d}u\right]\mathrm{d}x,
  \end{aligned}
\]
即积分号可交换.
\end{theorem}

\begin{proof}
因为 $f(x,u)$ 在 $[a,b]\times[\alpha,\beta]$ 上连续, 故\extref{定理 10.3.3} 表明二重积分与累次积分都存在, 且
\[
  \int_{\alpha}^{\beta}\left[\int_{a}^{b}f(x,u)\mathrm{d}x\right]\mathrm{d}u=\iint\limits_{[a,b]\times[a,b]}f(x,u)\mathrm{d}x\mathrm{d}u=\int_{a}^{b}\left[\int_{\alpha}^{\beta}f(x,u)\mathrm{d}u\right]\mathrm{d}x.
\]
\end{proof}

\begin{theorem}[可在积分号下求导]\label{thm:15.1.3}
如果 $f(x,u)$ 与偏导函数 $\dfrac{\partial f}{\partial u}(x,u)$ 都在闭区间 $[a,b]\times[\alpha,\beta]$ 上连续, 则
\[
  \varphi(u)=\int_{a}^{b}f(x,u)\mathrm{d}x
\]
在 $[\alpha,\beta]$ 上可导, 且
\[
  \varphi'(u)=\frac{\mathrm{d}\varphi(u)}{\mathrm{d}u}=\int_{a}^{b}\frac{\partial f(x,u)}{\partial u}\mathrm{d}x=\int_{a}^{b}f_{u}'(x,u)\mathrm{d}x
\]
(求导与积分号可交换, 或可积分号下求导).
\end{theorem}

\begin{proof}
令 $g(u)=\displaystyle\int_{a}^{b}\frac{\partial f(x,u)}{\partial u}\mathrm{d}x$, $u\in[\alpha,\beta]$, 则 $\forall v\in[\alpha,\beta]$ 有
\[
  \begin{aligned}
    \int_{\alpha}^{v}g(u)\mathrm{d}u&=\int_{\alpha}^{v}\left[\int_{a}^{b}\frac{\partial f(x,u)}{\partial u}\mathrm{d}x\right]\mathrm{d}u\xlongequal{\text{定理~\ref{thm:15.1.2}}}\int_{a}^{b}\left[\int_{\alpha}^{v}\frac{\partial f(x,u)}{\partial u}\mathrm{d}u\right]\mathrm{d}x\\
    &\xlongequal{\text{Newton-Leibniz 公式}}\int_{a}^{b}[f(x,v)-f(x,\alpha)]\mathrm{d}x=\varphi(v)-\varphi(\alpha),
  \end{aligned}
\]
由定理~\ref{thm:15.1.1} 知, $g(u)$ 在 $[\alpha,\beta]$ 上连续, 根据\extref{定理 6.3.2} (微积分基本定理) 得到
\[
  \begin{aligned}
    \int_{a}^{b}f_{u}'(x,v)\mathrm{d}x&=g(v)=\left(\int_{\alpha}^{v}g(u)\mathrm{d}u\right)'=\left(\varphi(v)-\varphi(\alpha)\right)'\\
    &=\varphi'(v)=\left(\int_{a}^{b}f(x,v)\mathrm{d}x\right)'.
  \end{aligned}
\]
\end{proof}

\begin{theorem}\label{thm:15.1.4}
如果 $f(x,u)$ 在闭区间 $[a,b]\times[\alpha,\beta]$ 上连续, $a(u)$, $b(u)$ 都在闭区间 $[\alpha,\beta]$ 上连续, 且 $a\leqslant a(u)$, $b(u)\leqslant b$, $u\in[\alpha,\beta]$, 则
\[
  \psi(u)=\int_{a(u)}^{b(u)}f(x,u)\mathrm{d}x
\]
在 $[\alpha,\beta]$ 上也连续.
\end{theorem}

\begin{proof}
证法 1 设 $M=\sup\limits_{(x,u)\in[a,b]\times[\alpha,\beta]}|f(x,u)|$, $u_{0}$, $u\in[\alpha,\beta]$, 由 $a(u)$ 与 $b(u)$ 在 $[\alpha,\beta]$ 都连续知
\[
  \begin{aligned}
    \left|\int_{a(u)}^{a(u_{0})}f(x,u)\mathrm{d}x\right|&\leqslant\int_{a(u)}^{a(u_{0})}|f(x,u)|\mathrm{d}x\leqslant M|a(u)-a(u_{0})|\to0\quad(u\to u_{0}),\\
    \left|\int_{b(u_{0})}^{b(u)}f(x,u)\mathrm{d}x\right|&\leqslant\int_{b(u_{0})}^{b(u)}|f(x,u)|\mathrm{d}x\leqslant M|b(u)-b(u_{0})|\to0\quad(u\to u_{0}).
  \end{aligned}
\]
因此
\[
  \lim_{u\to u_{0}}\int_{a(u)}^{a(u_{0})}f(x,u)\mathrm{d}x=0=\lim_{u\to u_{0}}\int_{b(u_{0})}^{b(u)}f(x,u)\mathrm{d}x.
\]
于是
\[
  \begin{aligned}
    \lim_{u\to u_{0}}\psi(u)&=\lim_{u\to u_{0}}\int_{a(u)}^{b(u)}f(x,u)\mathrm{d}x\\
    &=\lim_{u\to u_{0}}\left[\int_{a(u)}^{a(u_{0})}f(x,u)\mathrm{d}x+\int_{a(u_{0})}^{b(u_{0})}f(x,u)\mathrm{d}x+\int_{b(u_{0})}^{b(u)}f(x,u)\mathrm{d}x\right]\\
    &\xlongequal{\text{定理~\ref{thm:15.1.1}}}0+\int_{a(u_{0})}^{b(u_{0})}\lim_{u\to u_{0}}f(x,u)\mathrm{d}x+0=\int_{a(u_{0})}^{b(u_{0})}f(x,u_{0})\mathrm{d}x=\psi(u_{0}),
  \end{aligned}
\]
所以, $\psi(u)$ 在 $u_{0}$ 处连续, 由 $u_{0}\in[\alpha,\beta]$ 的任取性, $\psi(u)$ 在 $[\alpha,\beta]$ 上连续.

证法 2 设 $M=\sup\limits_{(x,u)\in[a,b]\times[\alpha,\beta]}|f(x,u)|$, $u_{0}\in[\alpha,\beta]$. $\forall\varepsilon>0$, 由于 $a(u)$, $b(u)$ 都在闭区间 $[\alpha,\beta]$ 上连续, 故 $\exists\delta_{1}>0$, 使得当 $u\in[\alpha,\beta]$, $|u-u_{0}|<\delta_{1}$ 时, 有 $|a(u)-a(u_{0})|<\dfrac{\varepsilon}{3(M+1)}$, $|b(u)-b(u_{0})|<\dfrac{\varepsilon}{3(M+1)}$. 再由 $f(x,u)$ 在 $[a,b]\times[\alpha,\beta]$ 上一致连续, 故 $\exists\delta_{2}>0$, 当 $\rho_{0}^{2}\big((x',u'),(x'',u'')\big)<\delta_{2}$ 时, $|f(x',u')-f(x'',u'')|<\dfrac{\varepsilon}{3(|b(u_{0})-a(u_{0})|+1)}$. 于是, 令 $\delta=\min\{\delta_{1},\delta_{2}\}$, 当 $|u-u_{0}|<\delta$ 时, 有
\[
  \begin{aligned}
    |\psi(u)-\psi(u_{0})|&=\left|\int_{a(u)}^{b(u)}f(x,u)\mathrm{d}x-\int_{a(u_{0})}^{b(u_{0})}f(x,u_{0})\mathrm{d}x\right|\\
    &\leqslant\left|\int_{a(u)}^{a(u_{0})}f(x,u)\mathrm{d}x\right|+\left|\int_{a(u_{0})}^{b(u_{0})}f(x,u)\mathrm{d}x-\int_{b(u_{0})}^{b(u)}f(x,u_{0})\mathrm{d}x\right|\\
    &\quad+\left|\int_{b(u_{0})}^{b(u)}f(x,u)\mathrm{d}x\right|\\
    &\leqslant\int_{a(u)}^{a(u_{0})}|f(x,u)|\mathrm{d}x+\int_{a(u_{0})}^{b(u_{0})}|f(x,u)-f(x,u_{0})|\mathrm{d}x\\
    &\quad+\int_{b(u_{0})}^{b(u)}|f(x,u)|\mathrm{d}x\\
    &\leqslant M|a(u)-a(u_{0})|+\frac{\varepsilon}{3(|b(u_{0})-a(u_{0})|+1)}|b(u_{0})-a(u_{0})|\\
    &\quad+M|b(u)-b(u_{0})|\\
    &\leqslant\frac{\varepsilon}{3}+\frac{\varepsilon}{3}+\frac{\varepsilon}{3}=\varepsilon.
  \end{aligned}
\]
这就证明了 $\lim\limits_{u\to u_{0}}\psi(u)=\psi(u_{0})$, 即 $\psi(u)$ 在 $u_{0}$ 连续. 再由 $u_{0}\in[\alpha,\beta]$ 任取知, $\psi(u)$ 在 $[\alpha,\beta]$ 上连续.
\end{proof}

\begin{theorem}[一般的参变量积分的求导公式]\label{thm:15.1.5}
设 $f(x,u)$ 在闭区间 $[a,b]\times[\alpha,\beta]$ 上连续, $\dfrac{\partial f}{\partial u}(x,u)$ 在 $[a,b]\times[\alpha,\beta]$ 也连续, $a(u)$, $b(u)$ 在 $[\alpha,\beta]$ 上可导, 且
\[
  a\leqslant a(u),\quad b(u)\leqslant b,\quad u\in[\alpha,\beta],
\]
则 $\psi(u)=\displaystyle\int_{a(u)}^{b(u)}f(x,u)\mathrm{d}x$ 在 $[\alpha,\beta]$ 上可导, 并且
\[
  \psi'(u)=\int_{a(u)}^{b(u)}f_{u}'(x,u)\mathrm{d}x+f(b(u),u)b'(u)-f(a(u),u)a'(u).
\]
\end{theorem}

\begin{proof}
证法 1 设 $\psi(u)=\displaystyle\int_{a(u)}^{b(u)}f(x,u)\mathrm{d}x$ 是由函数
\[
  F(u,a,b)=\int_{a}^{b}f(x,u)\mathrm{d}x
\]
与 $a=a(u)$, $b=b(u)$ 复合而成的. 根据复合函数求导的链规则得到
\[
  \begin{aligned}
    \psi'(u)&=F_{u}'+F_{b}'b'(u)+F_{a}'a'(u)\\
    &=\int_{a(u)}^{b(u)}f_{u}'(x,u)\mathrm{d}x+f(b(u),u)b'(u)-f(a(u),u)a'(u).
  \end{aligned}
\]

证法 2 $\forall u_{0}\in[\alpha,\beta]$, 有
\[
  \begin{aligned}
    \psi'(u_{0})&=\lim_{u\to u_{0}}\frac{\psi(u)-\psi(u_{0})}{u-u_{0}}=\lim_{u\to u_{0}}\frac{1}{u-u_{0}}\left[\int_{a(u)}^{b(u)}f(x,u)\mathrm{d}x-\int_{a(u_{0})}^{b(u_{0})}f(x,u_{0})\mathrm{d}x\right]\\
    &=\lim_{u\to u_{0}}\frac{1}{u-u_{0}}\left[\int_{a(u)}^{a(u_{0})}f(x,u)\mathrm{d}x+\int_{a(u_{0})}^{b(u_{0})}[f(x,u)-f(x,u_{0})]\mathrm{d}x+\int_{b(u_{0})}^{b(u)}f(x,u)\mathrm{d}x\right]\\
    &\xlongequal{\text{积分中值定理}}\lim_{u\to u_{0}}\left[-\frac{a(u)-a(u_{0})}{u-u_{0}}f(\overline{x},u)\right.\\
    &\left.\quad+\frac{1}{u-u_{0}}\int_{a(u_{0})}^{b(u_{0})}f_{u}'\big(x,u_{0}+\theta(u-u_{0})\big)(u-u_{0})\mathrm{d}x+\frac{b(u)-b(u_{0})}{u-u_{0}}f(\overline{\overline{x}},u)\right]\\
    &=-a'(u_{0})f(a(u_{0}),u_{0})+\int_{a(u_{0})}^{b(u_{0})}f_{u}'(x,u_{0})\mathrm{d}x+b'(u_{0})f(b(u_{0}),u_{0}).
  \end{aligned}
\]
由此, $\forall u\in[\alpha,\beta]$ 有
\[
  \psi'(u)=\int_{a(u)}^{b(u)}f_{u}'(x,u)\mathrm{d}x+f(b(u),u)b'(u)-f(a(u),u)a'(u).
\]
\end{proof}

\begin{example}\label{ex:15.1.1}
计算积分
\[
  \int_{0}^{1}\frac{x^{b}-x^{a}}{\ln x}\mathrm{d}x,\quad0<a<b.
\]
\end{example}

\begin{proof}
解法 1 由于 $x^{n}$ 在 $[0,1]\times[a,b]$ 上连续, 故
\[
  \begin{aligned}
    \int_{0}^{1}\frac{x^{b}-x^{a}}{\ln x}\mathrm{d}x&=\int_{0}^{1}\left[\int_{a}^{b}x^{u}\mathrm{d}u\right]\mathrm{d}x\xlongequal{\text{定理~\ref{thm:15.1.2}}}\int_{a}^{b}\left[\int_{0}^{1}x^{u}\mathrm{d}x\right]\mathrm{d}u\\
    &=\int_{a}^{b}\left.\frac{x^{u+1}}{u+1}\right|_{0}^{1}\mathrm{d}u=\int_{a}^{b}\frac{\mathrm{d}u}{u+1}=\left.\ln(u+1)\right|_{a}^{b}=\ln\frac{b+1}{a+1}.
  \end{aligned}
\]
解法 2 视 $I(b)=\displaystyle\int_{0}^{1}\frac{x^{b}-x^{a}}{\ln x}\mathrm{d}x$ 中的 $a$ 为固定的常数, 而 $b$ 视为参变量, 则
\[
  \frac{\mathrm{d}I}{\mathrm{d}b}\xlongequal{\text{定理~\ref{thm:15.1.3}}}\int_{0}^{1}\frac{\mathrm{d}}{\mathrm{d}b}\left(\frac{x^{b}-x^{a}}{\ln x}\right)\mathrm{d}x=\int_{0}^{1}x^{b}\mathrm{d}x=\left.\frac{x^{b+1}}{b+1}\right|_{0}^{1}=\frac{1}{b+1},
\]
所以, 重新得到
\[
  I(b)=\int_{a}^{b}\frac{\mathrm{d}I}{\mathrm{d}u}\mathrm{d}u+I(a)=\int_{a}^{b}\frac{\mathrm{d}u}{u+1}+0=\left.\ln(u+1)\right|_{a}^{b}=\ln\frac{b+1}{a+1}.
\]
\end{proof}

\begin{remark}\label{rem:15.1.1}
例~\ref{ex:15.1.1} 的解法中, $\dfrac{x^{b}-x^{a}}{\ln x}$ 当 $x=0,1$ 时无定义, 正确理解应为
\[
  \int_{a}^{b}x^{u}\mathrm{d}u=\begin{cases}
    0, & x=0,\\
    \left.\dfrac{x^{u}}{\ln x}\right|_{a}^{b}=\dfrac{x^{b}-x^{a}}{\ln x}, & 0<x<1,\\
    b-a, & x=1.
\end{cases}
\]
\[
  \begin{aligned}
    &\lim_{x\to0^{+}}\frac{x^{b}-x^{a}}{\ln x}\xlongequal{\text{L'Hospital 法则}}\lim_{x\to0^{+}}\frac{bx^{b-1}-ax^{a-1}}{\frac{1}{x}}=\lim_{x\to0^{+}}(bx^{b}-ax^{a})=0-0=0=\int_{a}^{b}0^{u}\mathrm{d}u,\\
    &\lim_{x\to1^{-}}\frac{x^{b}-x^{a}}{\ln x}\xlongequal{\text{L'Hospital 法则}}\lim_{x\to1^{-}}\frac{bx^{b-1}-ax^{a-1}}{\frac{1}{x}}=\lim_{x\to1^{-}}(bx^{b}-ax^{a})=b-a=\int_{a}^{b}1^{u}\mathrm{d}u.
  \end{aligned}
\]
例~\ref{ex:15.1.1} 的解法 2 中, 如果令
\[
  f(x,b)=\begin{cases}
    0, & x=0,\\
    \dfrac{x^{b}-x^{a}}{\ln x}, & 0<x<1,\\
    b-a, & x=1,
\end{cases}
\]
则
\[
  \begin{aligned}
    \frac{\mathrm{d}}{\mathrm{d}b}f(x,b)&=\begin{cases}
      0, & x=0,\\
      x^{b}, & 0<x<1,\\
      1, & x=1,
\end{cases}
  \end{aligned}
\]
在 $[0,1]\times[a,+\infty)$ 上连续. 因此, 上面应用定理~\ref{thm:15.1.3} 是合理的.

我们给出这样的注解是完全必要的, 大部分读者都会忽略这些! 但是, 千万不能因为这样枝节的问题, 束缚自己的手脚, 这也不敢做, 那也不敢做. 您只管大胆地用定理~\ref{thm:15.1.1}, 定理~\ref{thm:15.1.2} 与定理~\ref{thm:15.1.3}, 计算出结果后仔细验证是否满足定理的条件.
\end{remark}

\begin{example}\label{ex:15.1.2}
计算积分
\[
  I(r)=\int_{0}^{\pi}\ln(1-2r\cos x+r^{2})\mathrm{d}x,\quad|r|<1.
\]
\end{example}

\begin{proof}
将 $r$ 视作参变量, 应用定理~\ref{thm:15.1.3}, 对 $r$ 求导得
\[
  \begin{aligned}
    I'(r)&=\int_{0}^{\pi}\frac{\partial}{\partial r}\ln\left(1-2r\cos x+r^{2}\right)\mathrm{d}x\\
    &=\int_{0}^{\pi}\frac{-2\cos x+2r}{1-2r\cos x+r^{2}}\mathrm{d}x,
  \end{aligned}
\]
\[
  I'(0)=-\int_{0}^{\pi}2\cos x\mathrm{d}x=-\left.2\sin x\right|_{0}^{\pi}=0.
\]
当 $r\neq0$ 时, 有
\[
  \begin{aligned}
    I'(r)&=\frac{1}{r}\int_{0}^{\pi}\left(1-\frac{1-r^{2}}{1-2r\cos x+r^{2}}\right)\mathrm{d}x\\
    &\xlongequal{t=\tan\frac{x}{2}}\frac{1}{r}\left[\pi-(1-r^{2})\int_{0}^{+\infty}\frac{1}{1-2r\frac{1-t^{2}}{1+t^{2}}+r^{2}}\cdot\frac{2\mathrm{d}t}{1+t^{2}}\right]
  \end{aligned}
\]
\[
  \begin{aligned}
    &=\frac{1}{r}\left[\pi-2(1-r^{2})\int_{0}^{+\infty}\frac{\mathrm{d}t}{(1-r)^{2}+(1+r)^{2}t^{2}}\right]\\
    &=\frac{1}{r}\left[\pi-\frac{1+r}{1-r}\int_{0}^{+\infty}\frac{2\mathrm{d}t}{1+\left(\frac{1+r}{1-r}\right)^{2}}\right]\\
    &=\frac{1}{r}\left[\pi-2\arctan\frac{1+r}{1-r}t\bigg|_{0}^{+\infty}\right]=\frac{1}{r}\left[\pi-2\left(\frac{\pi}{2}-0\right)\right]=0.
  \end{aligned}
\]
由上得到 $I'(r)=0$, $|r|<1$. 所以 $I(r)=c$ (常值). 从而
\[
  I(r)=\int_{0}^{\pi}\ln(1-2r\cos x+r^{2})\mathrm{d}x=I(0)=\int_{0}^{\pi}0\mathrm{d}x=0,\quad|r|<1.
\]
\end{proof}

\begin{exblock}{练习题 15.1}
\begin{enumerate}
  \item 求极限:
  \[
    (1)\ \lim_{a\to0}\int_{-1}^{1}\sqrt{x^{2}+a^{2}}\mathrm{d}x;\qquad
    (2)\ \lim_{t\to0}\int_{0}^{2}x^{2}\cos tx\,\mathrm{d}x.
\]
  \item 设 $f(x)$ 为可导函数, 令 $F(x)=\displaystyle\int_{0}^{u}(x+u)f(x)\mathrm{d}x$, 计算 $F''(u)$.
  \item 计算下列函数的导数:

  (1) $f(x)=\displaystyle\int_{\sin x}^{\cos x}\mathrm{e}^{(1+t)^{2}}\mathrm{d}t$;

  (2) $f(x)=\displaystyle\int_{x}^{x^{2}}\mathrm{e}^{-x^{2}u^{2}}\mathrm{d}u$;

  (3) $f(x)=\displaystyle\int_{a+x}^{b+x}\frac{\sin xt}{t}\mathrm{d}t$;

  (4) $f(u)=\displaystyle\int_{0}^{u}g(x+u,x-u)\,\mathrm{d}x$, 其中 $g$ 为具有连续偏导数的二元函数.
  \item 设 $\varphi,\psi$ 分别可以微分两次与一次. 证明:
  \[
    u(x,t)=\frac{1}{2}\big[\varphi(x-at)+\varphi(x+at)\big]+\frac{1}{2a}\int_{x-at}^{x+at}\psi(s)\mathrm{d}s
\]
  \[
    \frac{\partial^{2}u}{\partial t^{2}}=a^{2}\frac{\partial^{2}u}{\partial x^{2}}
\]
  满足弦振动方程.
  \item 设 $f(t)$ 在闭区间 $[0,a]$ 上连续, 且当 $t\in[0,a]$ 时, $(x-t)^{2}+y^{2}+z^{2}\neq0$. 证明: 函数
  \[
    u(x,y,z)=\int_{0}^{a}\frac{f(t)}{\sqrt{(x-t)^{2}+y^{2}+z^{2}}}\mathrm{d}t
\]
  满足 Laplace 方程
  \[
    \Delta u=\frac{\partial^{2}u}{\partial x^{2}}+\frac{\partial^{2}u}{\partial y^{2}}+\frac{\partial^{2}u}{\partial z^{2}}=0.
\]
  \item 设 $f(t)$ 在 $[a,A]$ 上连续, 证明:
  \[
    \lim_{h\to0}\frac{1}{h}\int_{a}^{x}[f(t+h)-f(t)]\mathrm{d}t=f(x)-f(a),\quad a<x<A.
\]
  \item 设 $a<b$, $f(x)$ 为可导函数, 令
  \[
    \varphi(u)=\int_{a}^{b}f(x)|x-u|\mathrm{d}x,
\]
  计算 $\varphi''(u)$.
  \item 在区间 $[1,3]$ 上用线性函数 $a+bx$ 近似代替 $f(x)=x^{2}$, 试选 $a$, $b$, 使得
  \[
    \int_{1}^{3}(a+bx-x^{2})^{2}\mathrm{d}x
\]
  达到最小值.
\end{enumerate}
\end{exblock}

\begin{exblock}{思考题 15.1}
\begin{enumerate}
  \item 证明: $n$ 阶 Bessel 函数
  \[
    \mathrm{J}_{n}(x)=\frac{1}{\pi}\int_{0}^{\pi}\cos(n\varphi-x\sin\varphi)\mathrm{d}\varphi
\]
  满足 Bessel 方程
  \[
    x^{2}\mathrm{J}_{n}''(x)+x\mathrm{J}_{n}'(x)+(x^{2}-n^{2})\mathrm{J}_{n}(x)=0.
\]
  \item 利用对参数的求导法, 计算下列积分:

  (1) $\displaystyle\int_{0}^{\frac{\pi}{2}}\ln(a^{2}\sin^{2}x+b^{2}\cos^{2}x)\mathrm{d}x$;

  (2) $\displaystyle\int_{0}^{\frac{\pi}{2}}\frac{\arctan(a\tan x)}{\tan x}\mathrm{d}x$.
  \item 用下面两种方法证明
  \[
    \int_{0}^{\frac{\pi}{2}}\ln\frac{1+a\cos x}{1-a\cos x}\frac{\mathrm{d}x}{\cos x}=\pi\operatorname{arcsin}a,\quad|a|<1.
\]
  (1) 对参数 $a$ 求导;

  (2) 由 $\dfrac{1}{\cos x}\ln\dfrac{1+a\cos x}{1-a\cos x}=2a\displaystyle\int_{0}^{1}\frac{\mathrm{d}y}{1-(a^{2}\cos^{2}x)y^{2}}$.
  \item 用下面两种方法证明
  \[
    \int_{0}^{1}\frac{\ln(1+x)}{1+x^{2}}\mathrm{d}x=\frac{\pi}{8}\ln2.
\]
  (1) $I(\alpha)=\displaystyle\int_{0}^{1}\frac{\ln(1+\alpha x)}{1+x^{2}}\mathrm{d}x$ 关于参数 $\alpha$ 求导;

  (2) $\displaystyle\int_{0}^{1}\frac{\ln(1+x)}{1+x^{2}}\mathrm{d}x=\int_{0}^{\frac{\pi}{2}}\ln(1+\tan\theta)\mathrm{d}\theta$.
\end{enumerate}
\end{exblock}

% ==================================================================
%  徐森林、金亚东、薛春华《数学分析》（第三册）
%  S15.2  含参变量广义积分的一致收敛
%  源：原书印刷页 185-200（PDF 物理页 191-206）
%  底本：PaddleOCR-VL-1.6，见 chapters/_ch15_15_2_raw.md
%  约定：全西文标点；行间公式用一对方括号；多行公式用 align/aligned；
%        跨册引用一律 \extref；数学模式内的中文必须包 \text{}
%  带撇编号：定理 15.2.6'（函数项级数的逐项无穷积分）用 theoremp
%           + \setcounter{theorem}{6}，见 _make_ch15_2b.py
%  已订正的 OCR 错误与原书照录清单见 tools/_make_ch15_2a.py 头部说明
% ==================================================================
\section{含参变量广义积分的一致收敛}\label{sec:15.2}

在 6.5 节研究了广义积分的敛散性, 而在 15.1 节中引进了参变量的正常积分 $\displaystyle\int_{a}^{b}f(x,u)\mathrm{d}x$, 并给出了关于 $u$ 的连续、导数及积分的重要定理.

现考虑参变量广义积分, 它既是参变量积分, 又是广义积分, 它是两者的结合. 设 $f(x,u)$ 在 $[a,+\infty)\times[\alpha,\beta]$ 上连续. 如果对 $[\alpha,\beta]$ 中的任意 $u$, 无穷积分
\[
  \int_{a}^{+\infty}f(x,u)\mathrm{d}x
\]
都收敛, 则它就确定了 $[\alpha,\beta]$ 上的一个函数
\[
  \varphi(u)=\int_{a}^{+\infty}f(x,u)\mathrm{d}x.
\]
如何研究 $\varphi(u)$ 的连续、可导、积分等分析性质? 在第 13 章中已经看到, 研究由函数项级数所确定的函数的分析性质时, 函数项级数 $\sum_{n=1}^{\infty}u_{n}(x)$ 关于 $x\in I$ 的一致收敛概念起了关键的作用. 这里 $x$ 起到参量的作用, 它相当于 $\displaystyle\int_{1}^{+\infty}f(x,u)\mathrm{d}x$ 中的参变量 $u$.

对 $\displaystyle\int_{a}^{+\infty}f(x,u)\mathrm{d}x$, 类似的概念也起着同样重要的作用. 设积分 $\displaystyle\int_{a}^{+\infty}f(x,u)\mathrm{d}x$ 对任意 $u\in[\alpha,\beta]$ 都收敛, 即 $\forall\varepsilon>0$, $\exists A_{0}=A_{0}(\varepsilon,u)>a$, 当 $A>A_{0}$ 时, 有
\[
  \left|\int_{a}^{A}f(x,u)\mathrm{d}x-\int_{a}^{+\infty}f(x,u)\mathrm{d}x\right|=\left|\int_{A}^{+\infty}f(x,u)\mathrm{d}x\right|<\varepsilon.
\]
概括这样的考察, 就引出了下面类似于函数项级数的一致收敛概念.

\begin{definition}\label{def:15.2.1}
如果 $\forall\varepsilon>0$, 总 $\exists A_{0}=A_{0}(\varepsilon)$ (仅与 $\varepsilon$ 有关, 而与 $x\in I$ 无关!) $>a$, 当 $A>A_{0}$ 时, 有
\[
  \left|\int_{A}^{+\infty}f(x,u)\mathrm{d}x\right|<\varepsilon,\quad\forall u\in[\alpha,\beta],
\]
则称参变量无穷积分 $\displaystyle\int_{a}^{+\infty}f(x,u)\mathrm{d}x$ 关于 $u$ 在 $[\alpha,\beta]$ 上一致收敛.
\end{definition}

对于瑕积分也有类似的定义.

\begin{definition}\label{def:15.2.2}
设 $a$ 为瑕点, $\forall u\in I$, $\displaystyle\int_{a}f(x,u)\mathrm{d}x$ 收敛. 如果 $\forall\varepsilon>0$, $\exists\delta_{0}=\delta_{0}(\varepsilon)$ (仅与 $\varepsilon$ 有关, 而与 $x\in I$ 无关!) $>0$, 当 $\delta\in(0,\delta_{0})$, 有
\[
  \left|\int_{a}^{a+\delta}f(x,u)\mathrm{d}x\right|=\left|\int_{a+\delta}^{b}f(x,u)\mathrm{d}x-\int_{a}^{b}f(x,u)\mathrm{d}x\right|<\varepsilon,\quad\forall u\in[\alpha,\beta],
\]
则称瑕积分 $\displaystyle\int_{a}^{b}f(x,u)\mathrm{d}x$ 关于 $u\in[\alpha,\beta]$ 一致收敛.
\end{definition}

上面两个定义中的 $[\alpha,\beta]$ 可以是开区间也可以是无穷区间.

与广义积分收敛判别法一样, 这里也有一系列和函数项级数类似的一致收敛判别法.

下面只对无穷积分来讨论这些判别法. 而对于瑕积分也有类似的结果, 就不再一一说明了.

\begin{theorem}\label{thm:15.2.1}
设 $\eta(A)=\sup\limits_{u\in[\alpha,\beta]}\left|\int_{A}^{+\infty}f(x,u)\mathrm{d}x\right|$, 则无穷积分
\[
  \int_{a}^{+\infty}f(x,u)\mathrm{d}x\ \text{在}\ [\alpha,\beta]\ \text{上一致收敛}\ \Leftrightarrow\ \lim_{A\to+\infty}\eta(A)=0.
\]
\end{theorem}

\begin{proof}
($\Leftarrow$) 因为 $\lim\limits_{A\to+\infty}\eta(A)=0$, 所以 $\forall\varepsilon>0$, $\exists A_{0}=A_{0}(\varepsilon)>a$, 当 $A>A_{0}$ 时, 有
\[
  \eta(A)=|\eta(A)-0|<\varepsilon.
\]
于是, 当 $A>A_{0}$ 时, 有
\[
  \begin{aligned}
    \left|\int_{a}^{A}f(x,u)\mathrm{d}x-\int_{a}^{+\infty}f(x,u)\mathrm{d}x\right|&=\left|\int_{A}^{+\infty}f(x,u)\mathrm{d}x\right|<\sup_{v\in[\alpha,\beta]}\left|\int_{A}^{+\infty}f(x,v)\mathrm{d}x\right|\\
    &=\eta(A)<\varepsilon,
  \end{aligned}
\]
这就证明无穷积分 $\displaystyle\int_{a}^{+\infty}f(x,u)\mathrm{d}x$ 在 $[\alpha,\beta]$ 上一致收敛.

($\Rightarrow$) 设无穷积分 $\displaystyle\int_{a}^{+\infty}f(x,u)\mathrm{d}x$ 在 $[\alpha,\beta]$ 上一致收敛, 故 $\forall\varepsilon>0$, $\exists A_{0}=A_{0}(\varepsilon)>a$, 当 $A>A_{0}$ 时, 有
\[
  \left|\int_{A}^{+\infty}f(x,u)\mathrm{d}x\right|=\left|\int_{a}^{A}f(x,u)\mathrm{d}x-\int_{a}^{+\infty}f(x,u)\mathrm{d}x\right|<\frac{\varepsilon}{2},\quad\forall u\in[\alpha,\beta].
\]
于是, 当 $A>A_{0}$ 时, 有
\[
  \eta(A)=\sup_{u\in[\alpha,\beta]}\left|\int_{A}^{+\infty}f(x,u)\mathrm{d}x\right|\leqslant\frac{\varepsilon}{2}<\varepsilon,
\]
即
\[
  \lim_{A\to+\infty}\eta(A)=0.
\]
\end{proof}

\begin{example}\label{ex:15.2.1}
讨论无穷积分 $\displaystyle\int_{0}^{+\infty}ue^{-xu}\mathrm{d}x$ 分别在 (1) $u\in[\delta,+\infty)$, $\delta>0$, (2) $u\in(0,+\infty)$ 上关于 $u$ 的一致收敛性.
\end{example}

\begin{proof}
解法 1

(1) 因 $\delta>0$, 故
\[
  \begin{aligned}
    \eta(A)&=\sup_{u\in[\delta,+\infty)}\left|\int_{A}^{+\infty}u\mathrm{e}^{-xu}\mathrm{d}x\right|=\sup_{u\in[\delta,+\infty)}\left|-\mathrm{e}^{-xu}\right|_{A}^{+\infty}\\
    &=\sup_{u\in[\delta,+\infty)}\mathrm{e}^{-Au}=\mathrm{e}^{-A\delta}\to0\quad(A\to+\infty).
  \end{aligned}
\]
根据定理~\ref{thm:15.2.1} 知, 无穷积分 $\displaystyle\int_{0}^{+\infty}ue^{-xu}\mathrm{d}x$ 在 $[\delta,+\infty)$ ($\delta>0$) 上关于 $u$ 一致收敛.

(2) 因 $\eta(A)=\sup\limits_{u\in(0,+\infty)}\left|\int_{A}^{+\infty}ue^{-xu}\mathrm{d}x\right|=\sup\limits_{u\in(0,+\infty)}\left|-\mathrm{e}^{-xu}\right|_{A}^{+\infty}=\sup\limits_{u\in(0,+\infty)}\mathrm{e}^{-Au}=1$,
根据定理~\ref{thm:15.2.1} 知, 无穷积分 $\displaystyle\int_{0}^{+\infty}ue^{-xu}\mathrm{d}x$ 在 $(0,+\infty)$ 上关于 $u$ 不一致收敛.

解法 2 显然, $\forall u\geqslant0$, $\displaystyle\int_{0}^{+\infty}ue^{-xu}\mathrm{d}x$ 收敛.

(1) 当 $u\geqslant\delta>0$ 时, $\forall\varepsilon>0$, 取 $A_{0}=A_{0}(\varepsilon)>\dfrac{-\ln\varepsilon}{\delta}$, 当 $A>A_{0}$ 时, 有
\[
  \left|\int_{A}^{+\infty}u\mathrm{e}^{-xu}\mathrm{d}x\right|=\left|-\mathrm{e}^{-xu}\right|_{A}^{+\infty}=\mathrm{e}^{-Au}\leqslant\mathrm{e}^{-A_{0}\delta}<\varepsilon,\quad\forall u\in[\delta,+\infty).
\]
所以, $\displaystyle\int_{0}^{+\infty}ue^{-xu}\mathrm{d}x$ 在 $u\in[\delta,+\infty)$ 中一致收敛.

(2) (反证) 假设 $\displaystyle\int_{0}^{+\infty}ue^{-xu}\mathrm{d}x$ 在 $[0,+\infty)$ 中一致收敛, 则对 $\varepsilon_{0}=\frac{1}{4}>0$, $\exists A_{0}>0$, 当 $A>A_{0}$ 时, 有
\[
  \left|\int_{A}^{+\infty}u\mathrm{e}^{-xu}\mathrm{d}x\right|=\left|-\mathrm{e}^{-xu}\right|_{A}^{+\infty}=\mathrm{e}^{-Au}<\varepsilon_{0}=\frac{1}{4},\quad\forall u\in(0,+\infty).
\]
取定 $A>A_{0}$, 令 $u=\dfrac{1}{A}$, 则
\[
  \mathrm{e}^{-A\cdot\frac{1}{A}}=\mathrm{e}^{-1}=\frac{1}{\mathrm{e}}<\frac{1}{4},
\]
或令 $u\to0^{+}$ 得
\[
  1=\lim_{u\to0^{+}}\mathrm{e}^{-Au}\leqslant\frac{1}{4},
\]
矛盾.
\end{proof}

\begin{theorem}[参变量无穷积分的 Cauchy 收敛准则]\label{thm:15.2.2}
无穷积分 $\displaystyle\int_{a}^{+\infty}f(x,u)\mathrm{d}x$ 在 $[\alpha,\beta]$ 上一致收敛 $\Leftrightarrow$ $\forall\varepsilon>0$, $\exists A_{0}=A_{0}(\varepsilon)$ (仅与 $\varepsilon$ 有关, 而与 $x\in[\alpha,\beta]$ 无关!) $>a$, 当 $A'$, $A''>A_{0}$ 时, 有
\[
  \left|\int_{a}^{A''}f(x,u)\mathrm{d}x-\int_{a}^{A'}f(x,u)\mathrm{d}x\right|=\left|\int_{A'}^{A''}f(x,u)\mathrm{d}x\right|<\varepsilon.
\]
\end{theorem}

\begin{proof}
($\Rightarrow$) 设 $\displaystyle\int_{a}^{+\infty}f(x,u)\mathrm{d}x$ 在 $[\alpha,\beta]$ 上一致收敛, 则 $\forall\varepsilon>0$, $\exists A_{0}=A(\varepsilon)>a$, 当 $A>A_{0}$ 时, 有
\[
  \left|\int_{A}^{+\infty}f(x,u)\mathrm{d}x\right|=\left|\int_{a}^{A}f(x,u)\mathrm{d}x-\int_{a}^{+\infty}f(x,u)\mathrm{d}x\right|<\frac{\varepsilon}{2},\quad\forall u\in[\alpha,\beta].
\]
所以, 当 $A'$, $A''>A_{0}$ 时, 有
\[
  \begin{aligned}
    \left|\int_{A'}^{A''}f(x,u)\mathrm{d}x\right|&=\left|\int_{A'}^{+\infty}f(x,u)\mathrm{d}x-\int_{A''}^{+\infty}f(x,u)\mathrm{d}x\right|\\
    &\leqslant\left|\int_{A'}^{+\infty}f(x,u)\mathrm{d}x\right|+\left|\int_{A''}^{+\infty}f(x,u)\mathrm{d}x\right|<\frac{\varepsilon}{2}+\frac{\varepsilon}{2}=\varepsilon.
  \end{aligned}
\]

($\Leftarrow$) 由右边条件, $\forall\varepsilon>0$, $\exists A_{0}=A_{0}(\varepsilon)>a$, 当 $A$, $A'>A_{0}$ 时, 有
\[
  \left|\int_{A}^{A'}f(x,u)\mathrm{d}x\right|<\frac{\varepsilon}{2},\quad\forall u\in[\alpha,\beta].
\]
固定 $A$, 令 $A'\to+\infty$ 得到
\[
  \left|\int_{A}^{+\infty}f(x,u)\mathrm{d}x\right|\leqslant\frac{\varepsilon}{2}<\varepsilon,\quad\forall u\in[\alpha,\beta].
\]
这就证明了无穷积分 $\displaystyle\int_{a}^{+\infty}f(x,u)\mathrm{d}x$ 在 $[\alpha,\beta]$ 上一致收敛.
\end{proof}

\begin{theorem}[参变量无穷积分的 Weierstrass 判别法——优函数判别法]\label{thm:15.2.3}
设 $f(x,u)$ 对 $x$ 在 $[a,+\infty)$ 上连续. 如果存在 $[a,+\infty)$ 上的连续函数 $F$, 使得 $\displaystyle\int_{a}^{+\infty}F(x)\mathrm{d}x$ 收敛, 而且对一切充分大的 $x$ 及 $[\alpha,\beta]$ 上的一切 $u$, 都有
\[
  |f(x,u)|\leqslant F(x),
\]
则无穷积分 $\displaystyle\int_{a}^{+\infty}f(x,u)\mathrm{d}x$ 在 $[\alpha,\beta]$ 上一致收敛.
\end{theorem}

\begin{proof}
因为 $\displaystyle\int_{a}^{+\infty}F(x)\mathrm{d}x$ 收敛, 故由无穷积分的 Cauchy 收敛准则 $\forall\varepsilon>0$, $\exists A_{0}=A_{0}(\varepsilon)>a$, 只要 $A'$, $A''>A_{0}$, 便有
\[
  \left|\int_{A'}^{A''}F(x)\mathrm{d}x\right|<\varepsilon.
\]
由此便知
\[
  \left|\int_{A'}^{A''}f(x,u)\mathrm{d}x\right|\leqslant\int_{A'}^{A''}|f(x,u)|\mathrm{d}x\leqslant\left|\int_{A'}^{A''}F(x)\mathrm{d}x\right|<\varepsilon,\quad\forall u\in[\alpha,\beta].
\]
根据参变量无穷积分一致收敛的 Cauchy 收敛准则知, $\displaystyle\int_{a}^{+\infty}f(x,u)\mathrm{d}x$ 在 $[\alpha,\beta]$ 上一致收敛.
\end{proof}

\begin{example}\label{ex:15.2.2}
证明: (1) $\displaystyle\int_{1}^{+\infty}\frac{\sin ux}{x^{2}}\mathrm{d}x$; (2) $\displaystyle\int_{0}^{+\infty}\frac{\cos ux}{1+x^{2}}\mathrm{d}x$ 在 $u\in(-\infty,+\infty)$ 中一致收敛.
\end{example}

\begin{proof}
(1) 因为
\[
  \left|\frac{\sin ux}{x^{2}}\right|\leqslant\frac{1}{x^{2}},\quad\forall u\in(-\infty,+\infty),
\]
而 $\displaystyle\int_{1}^{+\infty}\frac{\mathrm{d}x}{x^{2}}$ 收敛, 根据 Weierstrass 判别法 (定理~\ref{thm:15.2.3}) 知, $\displaystyle\int_{1}^{+\infty}\frac{\sin ux}{x^{2}}\mathrm{d}x$ 在 $(-\infty,+\infty)$ 上一致收敛.

(2) 因为
\[
  \left|\frac{\cos ux}{1+x^{2}}\right|\leqslant\frac{1}{1+x^{2}},\quad\forall u\in(-\infty,+\infty),
\]
而 $\displaystyle\int_{0}^{+\infty}\frac{\mathrm{d}x}{1+x^{2}}$ 收敛, 根据 Weierstrass 判别法 (定理~\ref{thm:15.2.3}) 知, $\displaystyle\int_{0}^{+\infty}\frac{\cos ux}{1+x^{2}}\mathrm{d}x$ 在 $(-\infty,+\infty)$ 上一致收敛.
\end{proof}

\begin{example}\label{ex:15.2.3}
证明: 无穷积分
\[
  \int_{0}^{+\infty}\mathrm{e}^{-(u^{2}+\alpha)t}\sin t\,\mathrm{d}t\quad(\alpha>0\ \text{为固定常数})
\]
在 $u\in[0,+\infty)$ 中一致收敛.
\end{example}

\begin{proof}
因为
\[
  \left|\mathrm{e}^{-(u^{2}+\alpha)t}\sin t\right|\leqslant\mathrm{e}^{-\alpha t},\quad t\in[0,+\infty),\ u\in[0,+\infty),
\]
而 $\displaystyle\int_{0}^{+\infty}\mathrm{e}^{-\alpha t}\mathrm{d}t=-\frac{1}{\alpha}\left.\mathrm{e}^{-\alpha t}\right|_{0}^{+\infty}=\frac{1}{\alpha}$ 收敛, 故由 Weierstrass 判别法 (定理~\ref{thm:15.2.3}) 知, $\displaystyle\int_{0}^{+\infty}\mathrm{e}^{-(u^{2}+\alpha)t}\sin t\,\mathrm{d}t$ 在 $u\in[0,+\infty)$ 中一致收敛.
\end{proof}

\begin{example}\label{ex:15.2.4}
证明: 无穷积分
\[
  \int_{0}^{+\infty}\mathrm{e}^{-(u^{2}+\alpha)t}\sin t\,\mathrm{d}u\quad(\alpha\geqslant0\ \text{为固定常数})
\]
在 $t\in[0,+\infty)$ 中一致收敛.
\end{example}

\begin{proof}
因为
\[
  \lim_{t\to0^{+}}\frac{\mathrm{e}^{-\alpha t}\sin t}{\sqrt{t}}=\lim_{t\to0^{+}}\frac{\sin t}{t}\mathrm{e}^{-\alpha t}\sqrt{t}=1\cdot0=0,
\]
所以, $\forall\varepsilon>0$, $\exists\delta>0$, 当 $t\in(0,\delta)$ 时, 有
\[
  \left|\frac{\mathrm{e}^{-\alpha t}\sin t}{\sqrt{t}}\right|<\frac{2\varepsilon}{\sqrt{\pi}}.
\]
于是, $\forall A>0$, 有
\[
  \begin{aligned}
    \left|\int_{A}^{+\infty}\mathrm{e}^{-(u^{2}+\alpha)t}\sin t\,\mathrm{d}u\right|&=\left|\mathrm{e}^{-\alpha t}\sin t\int_{A}^{+\infty}\mathrm{e}^{-tu^{2}}\mathrm{d}u\right|\\
    &\xlongequal{x=\sqrt{t}u}\left|\frac{\mathrm{e}^{-\alpha t}\sin t}{\sqrt{t}}\int_{\sqrt{t}A}^{+\infty}\mathrm{e}^{-x^{2}}\mathrm{d}x\right|<\frac{2\varepsilon}{\sqrt{\pi}}\int_{0}^{+\infty}\mathrm{e}^{-x^{2}}\mathrm{d}x=\varepsilon,\quad\forall t\in(0,\delta).
  \end{aligned}
\]
由于 $t=0$ 时, 上式显然也成立, 因此, $\forall t\in[0,\delta)$, $\forall A>0$, 有
\[
  \left|\int_{A}^{+\infty}\mathrm{e}^{-(u^{2}+\alpha)t}\sin t\,\mathrm{d}u\right|<\varepsilon.
\]
当 $t\in[\delta,+\infty)$ 时, 由于
\[
  \left|\mathrm{e}^{-(u^{2}+\alpha)t}\sin t\right|\leqslant\mathrm{e}^{-\delta(\alpha+u^{2})}\leqslant\mathrm{e}^{-\delta u^{2}},
\]
而 $\displaystyle\int_{0}^{+\infty}\mathrm{e}^{-\delta u^{2}}\mathrm{d}u$ 收敛, 故由 Weierstrass 判别法 (定理~\ref{thm:15.2.3}) 知, 无穷积分 $\displaystyle\int_{0}^{+\infty}\mathrm{e}^{-(u^{2}+\alpha)t}\sin t\,\mathrm{d}u$ 在 $t\in[\delta,+\infty)$ 中一致收敛. 因而, $\exists A_{0}>0$, 当 $A>A_{0}$ 时, 有
\[
  \left|\int_{A}^{+\infty}\mathrm{e}^{-(u^{2}+\alpha)t}\sin t\,\mathrm{d}u\right|<\varepsilon
\]
对任意的 $t\in[\delta,+\infty)$ 成立. 综合上述知, 当 $A>A_{0}$ 时, 有
\[
  \left|\int_{A}^{+\infty}\mathrm{e}^{-(u^{2}+\alpha)t}\sin t\,\mathrm{d}u\right|<\varepsilon
\]
对任意 $t\in[0,+\infty)$ 成立. 这就证明了无穷积分
\[
  \int_{0}^{+\infty}\mathrm{e}^{-(u^{2}+\alpha)t}\sin t\,\mathrm{d}u
\]
在 $t\in[0,+\infty)$ 中一致收敛.
\end{proof}

更细致的判别法有下面的定理.

\begin{theorem}[参变量无穷积分的 Dirichlet 判别法]\label{thm:15.2.4}
如果函数 $f(x,u)$, $g(x,u)$ 满足:

(1) $g(x,u)$ 为 $x$ 的单调函数, 且当 $x\to+\infty$ 时关于 $u\in[\alpha,\beta]$ 一致趋于 $0$, 即 $\forall\varepsilon>0$, $\exists A_{0}=A_{0}(\varepsilon)>a$, 当 $x>A_{0}$ 时, $|g(x,u)|<\varepsilon$;

(2) $\forall A\geqslant a$, $\displaystyle\int_{a}^{A}f(x,u)\mathrm{d}x$ 对 $u\in[\alpha,\beta]$ 一致有界, 即 $\exists M>0$ ($M$ 为常数), 使得
\[
  \left|\int_{a}^{A}f(x,u)\mathrm{d}x\right|\leqslant M,\quad\forall u\in[\alpha,\beta],
\]
则 $\displaystyle\int_{a}^{+\infty}f(x,u)g(x,u)\mathrm{d}x$ 在 $u\in[\alpha,\beta]$ 上一致收敛.
\end{theorem}

\begin{proof}
$\forall\varepsilon>0$, 由条件 (1), $\exists A_{0}=A_{0}(\varepsilon)>0$, 当 $x>A_{0}$ 时, 有
\[
  |g(x,u)|<\frac{\varepsilon}{4M},\quad\forall u\in[\alpha,\beta].
\]
于是, 当 $A''>A'>A_{0}$ 时, 由一般的第二积分中值定理, $\exists\xi\in[A',A'']$ 使
\[
  \begin{aligned}
    \left|\int_{A'}^{A''}f(x,u)g(x,u)\mathrm{d}x\right|
    &=\left|g(A',u)\int_{A'}^{\xi}f(x,u)\mathrm{d}x+g(A'',u)\int_{\xi}^{A''}f(x,u)\mathrm{d}x\right|\\
    &\leqslant|g(A',u)|\left|\int_{A'}^{\xi}f(x,u)\mathrm{d}x\right|+|g(A',u)|\left|\int_{\xi}^{A''}f(x,u)\mathrm{d}x\right|\\
    &<\frac{\varepsilon}{4M}\left|\int_{a}^{\xi}f(x,u)\mathrm{d}x-\int_{a}^{A'}f(x,u)\mathrm{d}x\right|\\
    &\quad+\frac{\varepsilon}{4M}\left|\int_{a}^{A''}f(x,u)\mathrm{d}x-\int_{a}^{\xi}f(x,u)\mathrm{d}x\right|\\
    &\leqslant\frac{\varepsilon}{4M}(2M+2M)=\varepsilon,\quad\forall u\in[\alpha,\beta].
  \end{aligned}
\]
再由参变量无穷积分的 Cauchy 收敛准则知, $\displaystyle\int_{a}^{+\infty}f(x,u)g(x,u)\mathrm{d}x$ 在 $u\in[\alpha,\beta]$ 上一致收敛.
\end{proof}

\begin{theorem}[参变量无穷积分的 Abel 判别法]\label{thm:15.2.5}
如果 $f(x,u)$, $g(x,u)$ 满足:

(1) $g(x,u)$ 对 $x$ 单调, 且关于 $u\in[\alpha,\beta]$ 一致有界, 即 $\exists M>0$ (常数), 使得
\[
  |g(x,u)|\leqslant M,\quad(x,u)\in[a,+\infty)\times[\alpha,\beta];
\]

(2) 无穷积分 $\displaystyle\int_{a}^{+\infty}f(x,u)\mathrm{d}x$ 关于 $u\in[\alpha,\beta]$ 一致收敛.

则 $\displaystyle\int_{a}^{+\infty}f(x,u)g(x,u)\mathrm{d}x$ 在 $u\in[\alpha,\beta]$ 上一致收敛.
\end{theorem}

\begin{proof}
由条件 (1), $|g(x,u)|\leqslant M$, $(x,u)\in[a,+\infty)\times[\alpha,\beta]$.

由条件 (2), $\forall\varepsilon>0$, $\exists A_{0}=A_{0}(\varepsilon)>a$, 当 $A''>A'>A_{0}$ 时, 有
\[
  \left|\int_{A'}^{A}f(x,u)\mathrm{d}x\right|<\frac{\varepsilon}{2M},\quad u\in[\alpha,\beta].
\]
于是, 由一般的第二积分中值定理, $\exists\xi\in[A',A'']$ 使
\[
  \begin{aligned}
    \left|\int_{A'}^{A''}f(x,u)g(x,u)\mathrm{d}x\right|
    &=\left|g(A',u)\int_{A'}^{\xi}f(x,u)\mathrm{d}x+g(A'',u)\int_{\xi}^{A''}f(x,u)\mathrm{d}x\right|\\
    &\leqslant|g(A',u)|\left|\int_{A'}^{\xi}f(x,u)\mathrm{d}x\right|+|g(A'',u)|\left|\int_{\xi}^{A''}f(x,u)\mathrm{d}x\right|\\
    &<M\left(\frac{\varepsilon}{2M}+\frac{\varepsilon}{2M}\right)=\varepsilon,\quad\forall u\in[\alpha,\beta].
  \end{aligned}
\]
再由参变量无穷积分的 Cauchy 收敛准则知, $\displaystyle\int_{a}^{+\infty}f(x,u)g(x,u)\mathrm{d}x$ 在 $u\in[\alpha,\beta]$ 上一致收敛.
\end{proof}

\begin{example}\label{ex:15.2.5}
证明: 无穷积分 $\displaystyle\int_{0}^{+\infty}\mathrm{e}^{-xu}\frac{\sin x}{x}\mathrm{d}x$ 在 $u\in[0,+\infty)$ 上一致收敛.
\end{example}

\begin{proof}
由\extref{例 6.5.12} 知, $\displaystyle\int_{0}^{+\infty}\frac{\sin x}{x}\mathrm{d}x$ 收敛 (当然关于 $u\in[0,+\infty)$ 是一致收敛的), 而函数 $\mathrm{e}^{-xu}$ 对 $x$ 单调减少, 且 $0\leqslant\mathrm{e}^{-xu}\leqslant1$, $\forall(x,u)\in[0,+\infty)\times[0,+\infty)$. 故根据参变量无穷积分的 Abel 判别法 (定理~\ref{thm:15.2.5}) 知, $\displaystyle\int_{0}^{+\infty}\mathrm{e}^{-xu}\frac{\sin x}{x}\mathrm{d}x$ 在 $u\in[0,+\infty)$ 中一致收敛.
\end{proof}

\begin{example}\label{ex:15.2.6}
证明: 无穷积分 $\displaystyle\int_{0}^{+\infty}\frac{x\sin ux}{a^{2}+x^{2}}\mathrm{d}x$ 在 $[\delta,+\infty)$ 中一致收敛, 其中 $a$, $\delta>0$ 为常数.
\end{example}

\begin{proof}
$\forall A>0$, 有
\[
  \left|\int_{0}^{A}\sin ux\,\mathrm{d}x\right|=\left|-\frac{\cos ux}{u}\right|_{0}^{A}\left|=\frac{1-\cos uA}{u}\leqslant\frac{2}{\delta},\quad\forall u\in[\delta,+\infty).
\]
另一方面, 由于
\[
  \left(\frac{x}{a^{2}+x^{2}}\right)'=\frac{a^{2}-x^{2}}{(a^{2}+x^{2})^{2}}\begin{cases}
    >0, & 0\leqslant x<a,\ y(x)\ \text{单调增加},\\
    =0, & x=a,\quad y(a)\ \text{达最大值},\\
    <0, & x>a,\ y(x)\ \text{单调减少},
\end{cases}
\]
故当 $x\geqslant a$ 时, $\dfrac{x}{a^{2}+x^{2}}$ 单调减少, 且
\[
  \lim_{x\to+\infty}\frac{x}{a^{2}+x^{2}}=0.
\]
根据参变量无穷积分的 Abel 判别法 (定理~\ref{thm:15.2.5}) 知, $\displaystyle\int_{a}^{+\infty}\frac{x\sin ux}{a^{2}+x^{2}}\mathrm{d}x$ 在 $[\delta,+\infty)$ 中一致收敛. 再由
\[
  \left|\int_{0}^{A}f(x,u)\mathrm{d}x-\int_{0}^{+\infty}f(x,u)\mathrm{d}x\right|=\left|\int_{a}^{A}f(x,u)\mathrm{d}x-\int_{a}^{+\infty}f(x,u)\mathrm{d}x\right|
\]
知, $\displaystyle\int_{0}^{+\infty}\frac{x\sin ux}{a^{2}+x^{2}}\mathrm{d}x$ 在 $[\delta,+\infty)$ 中一致收敛.
\end{proof}

\begin{remark}\label{rem:15.2.1}
例~\ref{ex:15.2.6} 可看到, 从
\[
  \left|\int_{a}^{A}f(x,u)\mathrm{d}x-\int_{a}^{+\infty}f(x,u)\mathrm{d}x\right|=\left|\int_{b}^{A}f(x,u)\mathrm{d}x-\int_{b}^{+\infty}f(x,u)\mathrm{d}x\right|,
\]
可推得
$\displaystyle\int_{a}^{+\infty}f(x,u)\mathrm{d}x$ 在 $u\in[\alpha,\beta]$ 上一致收敛 $\Leftrightarrow\displaystyle\int_{b}^{+\infty}f(x,u)\mathrm{d}x$ 在 $u\in[\alpha,\beta]$ 上一致收敛.
\end{remark}

\begin{example}\label{ex:15.2.7}
证明: 瑕积分
\[
  \int_{0}^{1}\frac{\sin\frac{1}{x}}{x^{p}}\mathrm{d}x,
\]
(1) 当 $p<1$ 时绝对收敛; 当 $1\leqslant p<2$ 时条件收敛; 当 $2\leqslant p<+\infty$ 时发散.

(2) 关于 $p\in(0,2-\delta]$ 一致收敛, 其中 $\delta>0$.

(3) 关于 $p\in(0,2)$ 非一致收敛.
\end{example}

\begin{proof}
作变换 $t=\dfrac{1}{x}$, 则 $x=\dfrac{1}{t}$, 且
\[
  \int_{0}^{1}\frac{\sin\frac{1}{x}}{x^{p}}\mathrm{d}x=\int_{+\infty}^{1}\frac{\sin t}{\left(\frac{1}{t}\right)^{p}}\cdot\frac{-\mathrm{d}t}{t^{2}}=\int_{1}^{+\infty}\frac{\sin t}{t^{2-p}}\mathrm{d}t.
\]
显然, 它的敛散性、一致收敛性与 $\displaystyle\int_{0}^{1}\frac{\sin\frac{1}{x}}{x^{p}}\mathrm{d}x$ 完全一致.

(1) 根据\extref{例 6.5.12} 知, 上述两广义积分当 $2-p>1$, 即 $p<1$ 时绝对收敛; 当 $0<2-p\leqslant1$, 即 $1\leqslant p<2$ 时条件收敛; $2-p\leqslant0$, 即 $2\leqslant p<+\infty$ 时发散.

(2) 因为 $0<p\leqslant2-\delta$, 所以 $2-p\geqslant\delta$, 且
\[
  0<\frac{1}{t^{2-p}}\leqslant\frac{1}{t^{\delta}}.
\]
由此知, 当 $t\to+\infty$ 时, $\dfrac{1}{t^{2-p}}$ 单调减少一致趋于 $0$. 另外
\[
  \left|\int_{1}^{A}\sin t\,\mathrm{d}t\right|=|\cos1-\cos A|\leqslant2
\]
(一致有界). 根据 Dirichlet 判别法, $\displaystyle\int_{1}^{+\infty}\frac{\sin t}{t^{2-p}}\mathrm{d}t$ 在 $p\in(0,2-\delta]$ 上一致收敛.

(3) (反证) 假设 $\displaystyle\int_{1}^{+\infty}\frac{\sin t}{t^{2-p}}\mathrm{d}t$ 在 $p\in(0,2)$ 上一致收敛, 则对 $\varepsilon_{0}=1$, $\exists A_{0}=A_{0}(\varepsilon_{0})>a=1$, 只要 $A''>A'$, $A_{0}$, 便有
\[
  \left|\int_{A'}^{A''}\frac{\sin t}{t^{2-p}}\mathrm{d}t\right|<\varepsilon_{0}=1,\quad\forall p\in(0,2).
\]
今取 $A'=2k\pi$, $A''=(2k+1)\pi$, 则当 $k$ 充分大时, $A''>A'>A_{0}$. 于是, 由上面不等式得
\[
  1=\varepsilon_{0}>\left|\int_{2k\pi}^{(2k+1)\pi}\frac{\sin t}{t^{2-p}}\mathrm{d}t\right|\geqslant\frac{1}{[(2k+1)\pi]^{2-p}}\int_{2k\pi}^{(2k+1)\pi}\sin t\,\mathrm{d}t=\frac{2}{[(2k+1)\pi]^{2-p}},
\]
$\forall p\in(0,2)$. 令 $p\to2^{-}$ 即得 $1\geqslant2>1$, 矛盾.
\end{proof}

% ===== PART-B =====

下面给出几个参变量广义积分一致收敛的应用.

\begin{theorem}[无穷区间上极限号与广义积分号可交换的充分条件]\label{thm:15.2.6}
设函数列 $\{f_{n}(x)\}$ 在 $[a,+\infty)$ 上收敛于 $f(x)$. 如果:
\begin{enumerate}
  \item $\forall A>a$, $\{f_{n}(x)\}$ 在 $[a,A]$ 上一致收敛于 $f(x)$, 即 $\forall\varepsilon>0$, $\exists N_{0}=N_{0}(\varepsilon)\in\mathbb{N}$, 当 $n>N_{0}$ 时, 有
  \[
    |f_{n}(x)-f(x)|<\varepsilon,\quad\forall x\in[a,A];
\]
  \item 无穷积分 $\displaystyle\int_{a}^{+\infty}f_{n}(x)\,\mathrm{d}x$ 对 $n\in\mathbb{N}$ 一致收敛.
\end{enumerate}
则无穷积分 $\displaystyle\int_{a}^{+\infty}f(x)\,\mathrm{d}x$ 收敛, 且
\[
  \int_{a}^{+\infty}f(x)\,\mathrm{d}x=\int_{a}^{+\infty}\lim_{n\to+\infty}f_{n}(x)\,\mathrm{d}x=\lim_{n\to+\infty}\int_{a}^{+\infty}f_{n}(x)\,\mathrm{d}x.
\]
\end{theorem}

\begin{proof}
由条件 (2), $\forall\varepsilon>0$, $\exists A_{0}=A_{0}(\varepsilon)>a$, 当 $A''>A'\geqslant A_{0}$ 时, 便有
\[
  \left|\int_{A'}^{A''}f_{n}(x)\,\mathrm{d}x\right|<\frac{\varepsilon}{2},\quad\forall n\in\mathbb{N}\quad(n\text{ 视为参数}).
\]
因为 $\{f_{n}(x)\}$ 在 $[A',A'']$ 上 Riemann 可积, 且由条件 (1) 推得 $\{f_{n}(x)\}$ 在 $[A',A'']$ 上一致收敛于 $f(x)$, 故由定理~\ref{thm:13.2.3}, 有
\[
  \begin{aligned}
    \left|\int_{A'}^{A''}f(x)\,\mathrm{d}x\right|&=\left|\int_{A'}^{A''}\lim_{n\to+\infty}f_{n}(x)\,\mathrm{d}x\right|=\left|\lim_{n\to+\infty}\int_{A'}^{A''}f_{n}(x)\,\mathrm{d}x\right|\\
    &=\lim_{n\to+\infty}\left|\int_{A'}^{A''}f_{n}(x)\,\mathrm{d}x\right|\leqslant\frac{\varepsilon}{2}<\varepsilon.
  \end{aligned}
\]
根据无穷积分收敛的 Cauchy 收敛准则, $\displaystyle\int_{a}^{+\infty}f(x)\,\mathrm{d}x$ 收敛.

由条件 (2) 无穷积分 $\displaystyle\int_{a}^{+\infty}f_{n}(x)\,\mathrm{d}x$ 对 $n\in\mathbb{N}$ 一致收敛及上面已证得 $\displaystyle\int_{a}^{+\infty}f(x)\,\mathrm{d}x$ 收敛, 故 $\forall\varepsilon>0$, $\exists A_{1}=A_{1}(\varepsilon)>a$, 使得
\[
  \left|\int_{A_{1}}^{+\infty}f_{n}(x)\,\mathrm{d}x\right|<\frac{\varepsilon}{3},\quad n\in\mathbb{N},
\]
\[
  \left|\int_{A_{1}}^{+\infty}f(x)\,\mathrm{d}x\right|<\frac{\varepsilon}{3}.
\]
对固定的 $A_{1}=A_{1}(\varepsilon)$, 由条件 (1) $\exists N\in\mathbb{N}$, 当 $n>N$ 时, 有
\[
  |f_{n}(x)-f(x)|<\frac{\varepsilon}{3(A_{1}-a)},\quad\forall x\in[a,A_{1}].
\]
因而, 当 $n>N$ 时, 有
\[
  \begin{aligned}
    &\left|\int_{a}^{+\infty}f_{n}(x)\,\mathrm{d}x-\int_{a}^{+\infty}f(x)\,\mathrm{d}x\right|\\
    &\quad\leqslant\int_{a}^{A_{1}}|f_{n}(x)-f(x)|\,\mathrm{d}x+\left|\int_{A_{1}}^{+\infty}f_{n}(x)\,\mathrm{d}x\right|+\left|\int_{A_{1}}^{+\infty}f(x)\,\mathrm{d}x\right|\\
    &\quad<\frac{\varepsilon}{3(A_{1}-a)}\cdot(A_{1}-a)+\frac{\varepsilon}{3}+\frac{\varepsilon}{3}=\varepsilon.
  \end{aligned}
\]
这就证明了
\[
  \lim_{n\to+\infty}\int_{a}^{+\infty}f_{n}(x)\,\mathrm{d}x=\int_{a}^{+\infty}f(x)\,\mathrm{d}x.
\]
\end{proof}

与定理~\ref{thm:15.2.6} 对偶的函数项级数的逐项广义积分定理为下面的形式.

\setcounter{theorem}{6}
\begin{theoremp}[函数项级数的逐项无穷积分——求和号与广义积分可交换]\label{thm:15.2.6p}
设 $u_{n}(x)$ 在 $[a,+\infty)$ 上为连续函数, $n\in\mathbb{N}$. $\displaystyle\sum_{n=1}^{\infty}u_{n}(x)$ 在 $[a,+\infty)$ 上收敛于函数 $S(x)$. 如果:
\begin{enumerate}
  \item $\forall A>a$, 函数项级数 $\displaystyle\sum_{n=1}^{\infty}u_{n}(x)$ 在 $[a,A]$ 中一致收敛;
  \item 无穷积分 $\displaystyle\int_{a}^{+\infty}\sum_{k=1}^{n}u_{k}(x)\,\mathrm{d}x$ 对 $n\in\mathbb{N}$ 一致收敛.
\end{enumerate}
则无穷积分 $\displaystyle\int_{a}^{+\infty}\sum_{n=1}^{\infty}u_{n}(x)\,\mathrm{d}x$ 收敛, 且
\[
  \int_{a}^{+\infty}\sum_{n=1}^{\infty}u_{n}(x)\,\mathrm{d}x=\sum_{n=1}^{\infty}\int_{a}^{+\infty}u_{n}(x)\,\mathrm{d}x.
\]
\end{theoremp}

\begin{proof}
令 $S_{n}(x)=\displaystyle\sum_{k=1}^{n}u_{k}(x)$, $S(x)=\displaystyle\sum_{k=1}^{\infty}u_{k}(x)$, 则 $S_{n}(x)$ 与 $S(x)$ 满足定理~\ref{thm:15.2.6} 中 $f_{n}(x)$ 与 $f(x)$ 所满足的条件, 故
\[
  \begin{aligned}
    \int_{a}^{+\infty}\sum_{k=1}^{\infty}u_{k}(x)\,\mathrm{d}x&=\int_{a}^{+\infty}S(x)\,\mathrm{d}x=\int_{a}^{+\infty}\lim_{n\to+\infty}S_{n}(x)\,\mathrm{d}x\\
    &=\lim_{n\to+\infty}\int_{a}^{+\infty}S_{n}(x)\,\mathrm{d}x=\lim_{n\to+\infty}\int_{a}^{+\infty}\sum_{k=1}^{n}u_{k}(x)\,\mathrm{d}x\\
    &=\lim_{n\to+\infty}\sum_{k=1}^{n}\int_{a}^{+\infty}u_{k}(x)\,\mathrm{d}x=\sum_{k=1}^{\infty}\int_{a}^{+\infty}u_{k}(x)\,\mathrm{d}x.
  \end{aligned}
\]
为使读者熟悉函数项级数的描述, 建议将定理~\ref{thm:15.2.6} 证明过程对偶地用函数项级数加以证明.
\end{proof}

\begin{theorem}\label{thm:15.2.7}
设 $\displaystyle\sum_{n=1}^{\infty}u_{n}(x)$ 在 $[a,+\infty)$ 上收敛于 $S(x)$. 如果:
\begin{enumerate}
  \item $u_{n}(x)$ 在 $[a,+\infty)$ 上非负连续;
  \item $S(x)$ 在 $[a,+\infty)$ 上连续;
  \item 无穷积分 $\displaystyle\int_{a}^{+\infty}S(x)\,\mathrm{d}x$ 收敛.
\end{enumerate}
则
\[
  \int_{a}^{+\infty}S(x)\,\mathrm{d}x=\int_{a}^{+\infty}\sum_{n=1}^{\infty}u_{n}(x)\,\mathrm{d}x=\sum_{n=1}^{\infty}\int_{a}^{+\infty}u_{n}(x)\,\mathrm{d}x
\]
(可逐项无穷积分).
\end{theorem}

\begin{proof}
由条件 (1)、(2) 及 Dini 定理~\ref{thm:13.2.5p} 知, 函数项级数 $\displaystyle\sum_{n=1}^{\infty}u_{n}(x)$ 在 $[a,A]$ 上一致收敛 (其中 $A>a$). 因 $u_{n}(x)\geqslant0$, 故
\[
  0\leqslant S_{n}(x)\leqslant S(x),\quad n\in\mathbb{N}.
\]
从条件 (3), 无穷积分 $\displaystyle\int_{a}^{+\infty}S(x)\,\mathrm{d}x$ 收敛. 根据 Weierstrass 判别法, 无穷积分 $\displaystyle\int_{a}^{+\infty}S_{n}(x)\,\mathrm{d}x$ 对 $n\in\mathbb{N}$ 一致收敛. 应用定理~\ref{thm:15.2.6} 得到
\[
  \int_{a}^{+\infty}\sum_{n=1}^{\infty}u_{n}(x)\,\mathrm{d}x=\sum_{n=1}^{\infty}\int_{a}^{+\infty}u_{n}(x)\,\mathrm{d}x.
\]
\end{proof}

\begin{remark}\label{rem:15.2.2}
定理~\ref{thm:15.2.7} 关于函数列的对偶定理为: 设函数列 $\{S_{n}(x)\}$ 在 $[a,+\infty)$ 上收敛于 $S(x)$. 如果:
\begin{enumerate}
  \item $S_{n}(x)$ 在 $[a,+\infty)$ 上连续, 且 $S_{n}(x)\geqslant S_{n-1}(x)$, $\forall n\in\mathbb{N}$;
  \item $S(x)$ 在 $[a,+\infty)$ 上连续;
  \item 无穷积分 $\displaystyle\int_{a}^{+\infty}S(x)\,\mathrm{d}x$ 收敛.
\end{enumerate}
则
\[
  \int_{a}^{+\infty}S(x)\,\mathrm{d}x=\int_{a}^{+\infty}\lim_{n\to+\infty}S_{n}(x)\,\mathrm{d}x=\lim_{n\to+\infty}\int_{a}^{+\infty}S_{n}(x)\,\mathrm{d}x.
\]
\end{remark}

\begin{remark}\label{rem:15.2.3}
在注~\ref{rem:15.2.2} 中, 如果 $S_{n}(x)\leqslant S_{n-1}(x)$, 则结论仍正确. 只需用 $-S_{n}(x)$ 代 $S_{n}(x)$.
\end{remark}

\begin{example}[Euler-Poisson 积分]\label{ex:15.2.8}
证明: Euler-Poisson 积分
\[
  \int_{0}^{+\infty}\mathrm{e}^{-x^{2}}\,\mathrm{d}x=\frac{\sqrt{\pi}}{2}.
\]
\end{example}

\begin{proof}
证法 1, 2 (参阅\extref{例 6.5.7}, \extref{例 10.6.2}).

证法 3 因为 $\left(1+\dfrac{1}{x}\right)^{x}$ ($x>0$) 为严格单调增加的函数, 所以 $S_{n}(x)=\left(1+\dfrac{x^{2}}{n}\right)^{-n}$ 对 $n$ 单调减少, 且每个 $n\in\mathbb{N}$, $S_{n}(x)=\left(1+\dfrac{x^{2}}{n}\right)^{-n}$ 关于 $x\in[0,+\infty)$ 为连续函数, 于是
\[
  S(x)=\lim_{n\to+\infty}S_{n}(x)=\lim_{n\to+\infty}\left(1+\frac{x^{2}}{n}\right)^{-n}=\begin{cases}
    \lim\limits_{n\to+\infty}1=1, & x=0,\\
    \lim\limits_{n\to+\infty}\left[\left(1+\dfrac{x^{2}}{n}\right)^{\frac{n}{x^{2}}}\right]^{-x^{2}}=\mathrm{e}^{-x^{2}}, & x\neq0
\end{cases}
\]
$=\mathrm{e}^{-x^{2}}$ 在 $[0,+\infty)$ 上连续. 此外, 由于
\[
  \lim_{x\to+\infty}\frac{\mathrm{e}^{-x^{2}}}{\frac{1}{1+x^{2}}}=\lim_{x\to+\infty}\frac{1+x^{2}}{\mathrm{e}^{x^{2}}}=\lim_{u\to+\infty}\frac{1+u}{\mathrm{e}^{u}}=\lim_{u\to+\infty}\frac{1}{\mathrm{e}^{u}}=0
\]
及 $\displaystyle\int_{0}^{+\infty}\frac{\mathrm{d}x}{1+x^{2}}$ 收敛, 所以根据比较判别法的极限形式知, $\displaystyle\int_{0}^{+\infty}\mathrm{e}^{-x^{2}}\,\mathrm{d}x$ 也收敛. 于是
\[
  \begin{aligned}
    \int_{0}^{+\infty}\mathrm{e}^{-x^{2}}\,\mathrm{d}x&=\int_{0}^{+\infty}\lim_{n\to+\infty}\left(1+\frac{x^{2}}{n}\right)^{-n}\mathrm{d}x\xlongequal{\text{注~\ref{rem:15.2.3}}}\lim_{n\to+\infty}\int_{0}^{+\infty}\left(1+\frac{x^{2}}{n}\right)^{-n}\mathrm{d}x\\
    &\xlongequal{x=\sqrt{n}t}\lim_{n\to+\infty}\int_{0}^{+\infty}(1+t^{2})^{-n}\sqrt{n}\,\mathrm{d}t=\lim_{n\to+\infty}\sqrt{n}\int_{0}^{+\infty}\frac{\mathrm{d}t}{(1+t^{2})^{n}}\\
    &\xlongequal{t=\tan\theta}\lim_{n\to+\infty}\sqrt{n}\int_{0}^{\frac{\pi}{2}}\frac{\sec^{2}\theta}{(1+\tan^{2}\theta)^{n}}\,\mathrm{d}\theta=\lim_{n\to+\infty}\sqrt{n}\int_{0}^{\frac{\pi}{2}}\cos^{2n-2}\theta\,\mathrm{d}\theta\\
    &=\lim_{n\to+\infty}\sqrt{n}\frac{(2n-3)!!}{(2n-2)!!}\frac{\pi}{2}=\lim_{n\to+\infty}\sqrt{\frac{n}{2n-1}}\sqrt{\frac{1}{\left[\frac{(2n-2)!!}{(2n-3)!!}\right]^{2}}\frac{1}{2n-1}\frac{\pi}{2}}\\
    &\xlongequal{\text{\extref{例 6.4.5}}}\frac{1}{\sqrt{2}}\sqrt{\frac{1}{\frac{\pi}{2}}}\cdot\frac{\pi}{2}=\frac{\sqrt{\pi}}{2}.
  \end{aligned}
\]
\end{proof}

\begin{remark}\label{rem:15.2.4}
上述计算中最关键的一步是
\[
  \int_{0}^{+\infty}\mathrm{e}^{-x^{2}}\,\mathrm{d}x=\int_{0}^{+\infty}\lim_{n\to+\infty}\left(1+\frac{x^{2}}{n}\right)^{-n}\mathrm{d}x=\lim_{n\to+\infty}\int_{0}^{+\infty}\left(1+\frac{x^{2}}{n}\right)^{-n}\mathrm{d}x.
\]
大胆地将极限号搬到无穷积分号中去, 然后计算出结果 $\dfrac{\sqrt{\pi}}{2}$. 最后, 再来检查是否满足注~\ref{rem:15.2.3} 中条件. 大胆设想, 严格检查是每一个研究数学的学生必须具备的素质.
\end{remark}

\begin{exblock}[练习题 15.2]
\begin{enumerate}
  \item 研究下列参变量广义积分在指定区间上的一致收敛性:
  \begin{enumerate}
    \item $\displaystyle\int_{0}^{+\infty}\mathrm{e}^{-ux}\sin x\,\mathrm{d}x$, $0<u_{0}\leqslant u<+\infty$;
    \item $\displaystyle\int_{-\infty}^{+\infty}\frac{x^{2}\cos ux}{1+x^{4}}\,\mathrm{d}x$, $-\infty<u<+\infty$;
    \item $\displaystyle\int_{0}^{+\infty}\frac{\mathrm{d}x}{1+(x+u)^{2}}$, $0\leqslant u<+\infty$;
    \item $\displaystyle\int_{1}^{+\infty}\mathrm{e}^{-\alpha x}\frac{\cos x}{\sqrt{x}}\,\mathrm{d}x$, $0\leqslant\alpha<+\infty$;
    \item $\displaystyle\int_{0}^{+\infty}\sqrt{u}\mathrm{e}^{-ux^{2}}\,\mathrm{d}x$, $0\leqslant u<+\infty$.
  \end{enumerate}
  \item 证明: $\displaystyle\int_{-\infty}^{+\infty}\mathrm{e}^{-(x-\alpha)^{2}}\,\mathrm{d}x$ 在 $a\leqslant\alpha\leqslant b$ 上一致收敛; 在 $-\infty<\alpha<+\infty$ 上非一致收敛.
  \item 证明: $\displaystyle\int_{0}^{+\infty}\frac{x\sin\alpha x}{\alpha(1+x^{2})}\,\mathrm{d}x$ 在 $0<\alpha<+\infty$ 上非一致收敛.
  \item 证明: $\displaystyle\int_{0}^{1}\frac{1}{x^{\alpha}}\sin\frac{1}{x}\,\mathrm{d}x$ 在 $0<\alpha<2$ 上非一致收敛.
  \item 证明: $\displaystyle\int_{0}^{1}(1+x+x^{2}+\cdots+x^{n})\left(\ln\frac{1}{x}\right)^{\frac{1}{2}}\mathrm{d}x$ 关于 $n=1,2,\cdots$ 一致收敛.
  \item 证明: $I(\alpha)=\displaystyle\int_{0}^{+\infty}\frac{\mathrm{e}^{-x}}{|\sin x|^{\alpha}}\,\mathrm{d}x$ 对 $\alpha\in[0,b]$ (其中 $0<b<1$) 一致收敛.
  \item 证明: 参变量无穷积分 $\displaystyle\int_{1}^{+\infty}\mathrm{e}^{-\frac{1}{\alpha^{2}}\left(x-\frac{1}{\alpha}\right)^{2}}\,\mathrm{d}x$ 在 $0<\alpha<1$ 上虽然绝对一致收敛, 但并不能用 Weierstrass 判别法进行判断.
  \item 证明: $\displaystyle\int_{0}^{+\infty}\frac{\cos x^{2}}{x^{p}}\,\mathrm{d}x$ 在 $|p|\leqslant p_{0}<1$ 上一致收敛.
  \item 设函数在 $x>0$ 时连续, 积分 $\displaystyle\int_{0}^{+\infty}x^{\alpha}f(x)\,\mathrm{d}x$ 在 $\alpha=a,\alpha=b$ ($a<b$) 时收敛. 证明: 该积分在 $\alpha\in[a,b]$ 上一致收敛.
\end{enumerate}
\end{exblock}

\begin{exblock}[思考题 15.2]
\begin{enumerate}
  \item 证明: 参变量广义积分 $\displaystyle\int_{0}^{+\infty}\frac{x\sin\alpha x}{\alpha(1+x^{2})}\,\mathrm{d}x$ 在 $\alpha\in[\alpha_{0},+\infty)$ 中一致收敛, 在 $\alpha\in(0,\delta)$ 中不一致收敛, 其中 $\alpha_{0}$ 与 $\delta$ 是任意正数.
  \item 证明: 参变量广义积分 $\displaystyle\int_{0}^{+\infty}\frac{\sin ux}{x}\,\mathrm{d}x$ 在任何不包含 $u=0$ 的闭区间 $[a,b]$ 上一致收敛, 在包含 $u=0$ 的闭区间上非一致收敛.
  \item 证明: 参变量广义积分 $\displaystyle\int_{0}^{+\infty}\frac{\sin3x}{x+u}\mathrm{e}^{-ux}\,\mathrm{d}x$ 在 $u\in[0,+\infty)$ 中一致收敛.
  \item 设 $f(x,u)$ 在 $[a,+\infty)\times[\alpha,\beta]$ 中连续. 如果对每个 $u\in[\alpha,\beta)$, 参变量广义积分 $\displaystyle\int_{a}^{+\infty}f(x,u)\,\mathrm{d}x$ 都收敛, 但广义积分 $\displaystyle\int_{a}^{+\infty}f(x,\beta)\,\mathrm{d}x$ 发散, 证明: $\displaystyle\int_{a}^{+\infty}f(x,u)\,\mathrm{d}x$ 在 $[\alpha,\beta)$ 上非一致收敛.
  \item 证明: 参变量广义积分 $\displaystyle\int_{a}^{+\infty}f(x,u)\,\mathrm{d}x$ 在 $[\alpha,\beta]$ 上一致收敛等价于对任一单调增加趋于 $+\infty$ 的数列 $\{A_{n}\}$ ($A_{1}=a$), 函数项级数
  \[
    \sum_{n=1}^{\infty}\int_{A_{n}}^{A_{n+1}}f(x,u)\,\mathrm{d}x
\]
  在 $[\alpha,\beta]$ 上一致收敛.
  \item 证明: $g(\alpha)=\displaystyle\int_{0}^{+\infty}\frac{\ln(1+x^{3})}{x^{\alpha}}\,\mathrm{d}x$ 在 $\alpha\in(1,4)$ 上收敛, 在任何内闭区间 $[a,b]$ 上一致收敛, $g(\alpha)$ 在 $(1,4)$ 上连续.
  \item 设 $\displaystyle\int_{-\infty}^{+\infty}|f(x)|\,\mathrm{d}x$ 收敛, 证明:
  \[
    g(\alpha)=\int_{-\infty}^{+\infty}f(x)\cos\alpha x\,\mathrm{d}x
\]
  在 $\alpha\in(-\infty,+\infty)$ 上一致收敛, 且 $g(\alpha)$ 在 $(-\infty,+\infty)$ 上为一致连续函数.
\end{enumerate}
\end{exblock}

% ====================================================================
%  徐森林、金亚东、薛春华《数学分析》（第三册）
%  \S15.3  含参变量广义积分的性质
%  源：原书印刷页 198 起（PDF 物理页 204 起）
%  底本：PaddleOCR-VL-1.6，见 chapters/_ch15_15_3_raw.md
%  约定：全西文标点；行间公式用一对方括号；多行公式用 align 或 aligned；
%        跨册引用一律用 extref（写作「定理 x.y.z」原文）；数学模式内的中文必须包text
% ====================================================================
\section{含参变量广义积分的性质}\label{sec:15.3}

设含参变量无穷积分 $\displaystyle\int_{a}^{+\infty}f(x,u)\,\mathrm{d}x$ 对 $[\alpha,\beta]$ 中每个 $u$ 都收敛, 我们要研究由它所确定的函数
\[
  \varphi(u)=\int_{a}^{+\infty}f(x,u)\,\mathrm{d}x
\]
的性质.

与连续函数项级数的一致收敛性保证了级数和函数的连续性一样, 无穷积分 $\displaystyle\int_{a}^{+\infty}f(x,u)\,\mathrm{d}x$ 关于 $u\in[\alpha,\beta]$ 的一致收敛性保证了 $\varphi(u)$ 的连续性.

\begin{theorem}[含参变量无穷积分保持连续性]\label{thm:15.3.1}
如果 $f(x,u)$ 在 $[a,+\infty)\times[\alpha,\beta]$ 中连续, 且无穷积分 $\displaystyle\int_{a}^{+\infty}f(x,u)\,\mathrm{d}x$ 在 $u\in[\alpha,\beta]$ 上一致收敛, 则函数
\[
  \varphi(u)=\int_{a}^{+\infty}f(x,u)\,\mathrm{d}x
\]
在 $[\alpha,\beta]$ 上连续, 即
\[
  \begin{aligned}
    \lim_{u\to u_{0}}\int_{a}^{+\infty}f(x,u)\,\mathrm{d}x&=\lim_{u\to u_{0}}\varphi(u)=\varphi(u_{0})\\
    &=\int_{a}^{+\infty}f(x,u_{0})\,\mathrm{d}x=\int_{a}^{+\infty}\left[\lim_{u\to u_{0}}f(x,u)\right]\mathrm{d}x
  \end{aligned}
\]
(即极限号与无穷积分号可以交换).
\end{theorem}

\begin{proof}
$\forall u_{0}\in[\alpha,\beta]$, $\forall\varepsilon>0$, 因为 $\displaystyle\int_{a}^{+\infty}f(x,u)\,\mathrm{d}x$ 在 $u\in[\alpha,\beta]$ 上一致收敛, 故 $\exists A_{0}=A_{0}(\varepsilon)>a$, 当 $A>A_{0}$ 时, $\forall u\in[\alpha,\beta]$, 有
\[
  \left|\int_{A}^{+\infty}f(x,u)\,\mathrm{d}x\right|<\frac{\varepsilon}{3}.
\]
对固定的 $A>A_{0}$, 由定理~\ref{thm:15.1.1} 知, $\displaystyle\int_{a}^{A}f(x,u)\,\mathrm{d}x$ 在 $u\in[\alpha,\beta]$ 上连续, 故 $\exists\delta>0$, 当 $|u-u_{0}|<\delta$ 时, 有
\[
  \left|\int_{a}^{A}f(x,u)\,\mathrm{d}x-\int_{a}^{A}f(x,u_{0})\,\mathrm{d}x\right|<\frac{\varepsilon}{3}.
\]
于是
\[
  \begin{aligned}
    |\varphi(u)-\varphi(u_{0})|&=\left|\int_{a}^{+\infty}f(x,u)\,\mathrm{d}x-\int_{a}^{+\infty}f(x,u_{0})\,\mathrm{d}x\right|\\
    &=\left|\int_{a}^{A}f(x,u)\,\mathrm{d}x-\int_{a}^{A}f(x,u_{0})\,\mathrm{d}x+\int_{A}^{+\infty}f(x,u)\,\mathrm{d}x-\int_{A}^{+\infty}f(x,u_{0})\,\mathrm{d}x\right|\\
    &\leqslant\left|\int_{a}^{A}f(x,u)\,\mathrm{d}x-\int_{a}^{A}f(x,u_{0})\,\mathrm{d}x\right|+\left|\int_{A}^{+\infty}f(x,u)\,\mathrm{d}x\right|+\left|\int_{A}^{+\infty}f(x,u_{0})\,\mathrm{d}x\right|\\
    &<\frac{\varepsilon}{3}+\frac{\varepsilon}{3}+\frac{\varepsilon}{3}=\varepsilon,
  \end{aligned}
\]
所以, $\varphi(u)$ 在 $u_{0}$ 处连续. 由 $u_{0}\in[\alpha,\beta]$ 的任意性, 故 $\varphi(u)$ 在 $[\alpha,\beta]$ 上连续.
\end{proof}

与函数项级数类似, 无穷积分的一致收敛只是保证 $\varphi(u)$ 连续的一个充分条件, 并不必要. 但在 $f(x,u)$ 非负的条件下, 无穷积分的一致收敛便是 $\varphi(u)$ 连续的充要条件. 与 Dini 定理~\ref{thm:13.2.5} 相对偶的是下面的定理.

\begin{theorem}[Dini]\label{thm:15.3.2}
设 $f(x,u)$ 在 $[a,+\infty)\times[\alpha,\beta]$ 中连续、非负. 如果
\[
  \varphi(u)=\int_{a}^{+\infty}f(x,u)\,\mathrm{d}x
\]
在 $u\in[\alpha,\beta]$ 上连续, 则无穷积分 $\displaystyle\int_{a}^{+\infty}f(x,u)\,\mathrm{d}x$ 在 $u\in[\alpha,\beta]$ 上一致收敛.
\end{theorem}

\begin{proof}
\textbf{证法 1} (利用函数项级数的 Dini 定理) 因为 $f(x,u)$ 连续及非负, 故积分
\[
  a_{n}(u)=\int_{a+n-1}^{a+n}f(x,u)\,\mathrm{d}x
\]
也连续、非负, 而 (和函数)
\[
  \varphi(u)=\int_{a}^{+\infty}f(x,u)\,\mathrm{d}x=\sum_{n=1}^{\infty}\int_{a+n-1}^{a+n}f(x,u)\,\mathrm{d}x=\sum_{n=1}^{\infty}a_{n}(u)
\]
在 $[\alpha,\beta]$ 上连续, 根据函数项级数的 Dini 定理~\ref{thm:13.2.5}, $\displaystyle\sum_{n=1}^{\infty}a_{n}(u)$ 在 $[\alpha,\beta]$ 上一致收敛. 所以, $\forall\varepsilon>0$, $\exists N=N(\varepsilon)\in\mathbb{N}$, 使得 $\forall u\in[\alpha,\beta]$, 有
\[
  \left|\sum_{n=1}^{N}a_{n}(u)-\varphi(u)\right|=\sum_{n=N+1}^{\infty}a_{n}(u)<\varepsilon.
\]
取 $A_{0}=a+N=a+N(\varepsilon)$, 由 $f(x,u)$ 非负, 故当 $A>A_{0}$ 时, 有
\[
  \begin{aligned}
    0&\leqslant\int_{A}^{+\infty}f(x,u)\,\mathrm{d}x\leqslant\int_{A_{0}}^{+\infty}f(x,u)\,\mathrm{d}x=\sum_{n=N+1}^{\infty}\int_{a+n-1}^{a+n}f(x,u)\,\mathrm{d}x\\
    &=\sum_{n=N+1}^{\infty}a_{n}(u)<\varepsilon,\quad\forall u\in[\alpha,\beta].
  \end{aligned}
\]
这就证明了 $\displaystyle\int_{a}^{+\infty}f(x,u)\,\mathrm{d}x$ 在 $u\in[\alpha,\beta]$ 上一致收敛.

\textbf{证法 2} (反证) 假设 $\varphi(u)=\displaystyle\int_{a}^{+\infty}f(x,u)\,\mathrm{d}x$ 在 $[\alpha,\beta]$ 上不一致收敛, 则 $\exists\varepsilon_{0}>0$, $\forall n\in\mathbb{N}$, $\exists u_{n}\in[\alpha,\beta]$, 有
\[
  \int_{a+n}^{+\infty}f(x,u_{n})\,\mathrm{d}x\geqslant\varepsilon_{0}\iff\int_{a}^{a+n}f(x,u_{n})\,\mathrm{d}x\leqslant\int_{a}^{+\infty}f(x,u_{n})\,\mathrm{d}x-\varepsilon_{0}=\varphi(u_{n})-\varepsilon_{0}.
\]
由于 $u_{n}\in[\alpha,\beta]$, $\forall n\in\mathbb{N}$, 根据 Bolzano-Weierstrass 序列紧性知, $\{u_{n}\}$ 有收敛子列 $\{u_{n_{k}}\}$, 设 $\lim\limits_{k\to+\infty}u_{n_{k}}=u_{0}\in[\alpha,\beta]$. 对固定的 $u_{0}$, 因为 $\displaystyle\int_{a}^{+\infty}f(x,u_{0})\,\mathrm{d}x$ 收敛, 故 $\exists m\in\mathbb{N}$, 使得
\[
  \int_{a+m}^{+\infty}f(x,u_{0})\,\mathrm{d}x<\varepsilon_{0}\iff\int_{a}^{a+m}f(x,u_{0})\,\mathrm{d}x>\int_{a}^{+\infty}f(x,u_{0})\,\mathrm{d}x-\varepsilon_{0}=\varphi(u_{0})-\varepsilon_{0}.
\]
但我们知道, 当 $k$ 充分大时, $n_{k}>m$, 即有
\[
  \varphi(u_{n_{k}})-\varepsilon_{0}\geqslant\int_{a}^{a+n_{k}}f(x,u_{n_{k}})\,\mathrm{d}x\xlongequal[f\text{ 非负}]{\quad\geqslant\quad}\int_{a}^{a+m}f(x,u_{n_{k}})\,\mathrm{d}x,
\]
两边令 $k\to+\infty$, 并由 $\varphi(u)$ 连续及定理~\ref{thm:15.1.1} 得到
\[
  \varphi(u_{0})-\varepsilon_{0}\geqslant\int_{a}^{a+m}f(x,u_{0})\,\mathrm{d}x>\varphi(u_{0})-\varepsilon_{0},
\]
矛盾.
\end{proof}

\begin{remark}\label{rem:15.3.1}
思考题 15.3 第 4 题表明, Dini 定理~\ref{thm:15.3.2} 中有界闭区间 $[\alpha,\beta]$ 改为开区间或无穷区间, 结论未必成立.
\end{remark}

关于 $\varphi(u)$ 的积分有下面的定理.

\begin{theorem}[有穷积分号与无穷积分号可交换]\label{thm:15.3.3}
如果 $f(x,u)$ 在 $[a,+\infty)\times[\alpha,\beta]$ 中连续, $\varphi(u)=\displaystyle\int_{a}^{+\infty}f(x,u)\,\mathrm{d}x$ 关于 $u\in[\alpha,\beta]$ 一致收敛, 则 $\varphi(u)$ 在 $[\alpha,\beta]$ 上 Riemann 可积, 且
\[
  \int_{\alpha}^{\beta}\varphi(u)\,\mathrm{d}u=\int_{\alpha}^{\beta}\left(\int_{a}^{+\infty}f(x,u)\,\mathrm{d}x\right)\mathrm{d}u=\int_{a}^{+\infty}\left(\int_{\alpha}^{\beta}f(x,u)\,\mathrm{d}u\right)\mathrm{d}x
\]
(即对 $x$ 与 $u$ 的积分次序可以交换).
\end{theorem}

\begin{proof}
\textbf{证法 1} 根据定理~\ref{thm:15.3.1}, $\varphi(u)$ 在 $[\alpha,\beta]$ 上连续, 当然在 $[\alpha,\beta]$ 上 Riemann 可积. 因为 $\displaystyle\int_{a}^{+\infty}f(x,u)\,\mathrm{d}x$ 在 $u\in[\alpha,\beta]$ 上一致收敛, 故 $\forall\varepsilon>0$, $\exists A_{0}=A_{0}(\varepsilon)>a$, 使得当 $A>A_{0}$ 时, 有
\[
  \left|\int_{A}^{+\infty}f(x,u)\,\mathrm{d}x\right|<\frac{\varepsilon}{\beta-\alpha},\quad\forall u\in[\alpha,\beta].
\]
取定 $A>A_{0}$, 由定理~\ref{thm:15.1.2} 知
\[
  \int_{\alpha}^{\beta}\left(\int_{a}^{A}f(x,u)\,\mathrm{d}x\right)\mathrm{d}u=\int_{a}^{A}\left(\int_{\alpha}^{\beta}f(x,u)\,\mathrm{d}u\right)\mathrm{d}x.
\]
于是
\[
  \begin{aligned}
    &\left|\int_{a}^{A}\left(\int_{\alpha}^{\beta}f(x,u)\,\mathrm{d}u\right)\mathrm{d}x-\int_{\alpha}^{\beta}\left(\int_{a}^{+\infty}f(x,u)\,\mathrm{d}x\right)\mathrm{d}u\right|\\
    &\quad=\left|\int_{a}^{A}\left(\int_{\alpha}^{\beta}f(x,u)\,\mathrm{d}u\right)\mathrm{d}x-\int_{\alpha}^{\beta}\left(\int_{a}^{A}f(x,u)\,\mathrm{d}x\right)\mathrm{d}u-\int_{\alpha}^{\beta}\left(\int_{A}^{+\infty}f(x,u)\,\mathrm{d}x\right)\mathrm{d}u\right|\\
    &\quad=\left|\int_{\alpha}^{\beta}\left(\int_{A}^{+\infty}f(x,u)\,\mathrm{d}x\right)\mathrm{d}u\right|\leqslant\int_{\alpha}^{\beta}\left|\int_{A}^{+\infty}f(x,u)\,\mathrm{d}x\right|\mathrm{d}u<\frac{\varepsilon}{\beta-\alpha}(\beta-\alpha)=\varepsilon.
  \end{aligned}
\]
这就证明了
\[
  \int_{a}^{+\infty}\left(\int_{\alpha}^{\beta}f(x,u)\,\mathrm{d}u\right)\mathrm{d}x=\int_{\alpha}^{\beta}\left(\int_{a}^{+\infty}f(x,u)\,\mathrm{d}x\right)\mathrm{d}u.
\]

\textbf{证法 2} 记
\[
  \int_{a}^{+\infty}f(x,u)\,\mathrm{d}x=\sum_{n=1}^{\infty}\int_{a+n-1}^{a+n}f(x,u)\,\mathrm{d}x.
\]
因为无穷积分 $\displaystyle\int_{a}^{+\infty}f(x,u)\,\mathrm{d}x$ 在 $u\in[\alpha,\beta]$ 中一致收敛, 故 $\forall\varepsilon>0$, $\exists A_{0}=A_{0}(\varepsilon)>a$, 当 $A>A_{0}$ 时, 有
\[
  \left|\int_{A}^{+\infty}f(x,u)\,\mathrm{d}x\right|<\varepsilon,\quad\forall u\in[\alpha,\beta].
\]
取 $N\in\mathbb{N}$, 使得 $N>A_{0}-a$, 即 $a+N>A_{0}$. 于是, $\forall u\in[\alpha,\beta]$, 有
\[
  \left|\sum_{n=N+1}^{\infty}\int_{a+n-1}^{a+n}f(x,u)\,\mathrm{d}x\right|=\left|\int_{a+N}^{+\infty}f(x,u)\,\mathrm{d}x\right|<\varepsilon.
\]
因此, $\displaystyle\sum_{n=1}^{\infty}\int_{a+n-1}^{a+n}f(x,u)\,\mathrm{d}x$ 在 $u\in[\alpha,\beta]$ 上一致收敛. 这就得到
\[
  \begin{aligned}
    \int_{\alpha}^{\beta}\left(\int_{a}^{+\infty}f(x,u)\,\mathrm{d}x\right)\mathrm{d}u&=\int_{\alpha}^{\beta}\left(\sum_{n=1}^{\infty}\int_{a+n-1}^{a+n}f(x,u)\,\mathrm{d}x\right)\mathrm{d}u\\
    &\xlongequal{\text{定理~\ref{thm:13.2.3p}}}\sum_{n=1}^{\infty}\int_{\alpha}^{\beta}\left(\int_{a+n-1}^{a+n}f(x,u)\,\mathrm{d}x\right)\mathrm{d}u\\
    &\xlongequal{\text{定理~\ref{thm:15.1.2}}}\sum_{n=1}^{\infty}\int_{a+n-1}^{a+n}\left(\int_{\alpha}^{\beta}f(x,u)\,\mathrm{d}u\right)\mathrm{d}x\\
    &=\lim_{N\to+\infty}\sum_{n=1}^{N}\int_{a+n-1}^{a+n}\left(\int_{\alpha}^{\beta}f(x,u)\,\mathrm{d}u\right)\mathrm{d}x\\
    &=\lim_{N\to+\infty}\int_{a}^{a+N}\left(\int_{\alpha}^{\beta}f(x,u)\,\mathrm{d}u\right)\mathrm{d}x\\
    &=\lim_{A\to+\infty}\int_{a}^{A}\left(\int_{\alpha}^{\beta}f(x,u)\,\mathrm{d}u\right)\mathrm{d}x\\
    &=\int_{a}^{+\infty}\left(\int_{\alpha}^{\beta}f(x,u)\,\mathrm{d}u\right)\mathrm{d}x,
  \end{aligned}
\]
其中最后一个等式成立是因为 $\displaystyle\int_{a}^{+\infty}f(x,u)\,\mathrm{d}x$ 关于 $u\in[\alpha,\beta]$ 一致收敛, 故 $\forall\varepsilon>0$, $\exists A_{1}=A_{1}(\varepsilon)>a$, 当 $A',A''>A_{1}$ 时, 有
\[
  \left|\int_{A'}^{A''}f(x,u)\,\mathrm{d}x\right|<\frac{\varepsilon}{\beta-\alpha},\quad\forall u\in[\alpha,\beta].
\]
因此
\[
  \begin{aligned}
    \left|\int_{A'}^{A''}\left(\int_{\alpha}^{\beta}f(x,u)\,\mathrm{d}u\right)\mathrm{d}x\right|&\xlongequal{\text{定理~\ref{thm:15.1.2}}}\left|\int_{\alpha}^{\beta}\left(\int_{A'}^{A''}f(x,u)\,\mathrm{d}x\right)\mathrm{d}u\right|\\
    &\leqslant\int_{\alpha}^{\beta}\left|\int_{A'}^{A''}f(x,u)\,\mathrm{d}x\right|\mathrm{d}u<\frac{\varepsilon}{\beta-\alpha}(\beta-\alpha)=\varepsilon.
  \end{aligned}
\]
这就证明了 $\displaystyle\int_{a}^{+\infty}\left(\int_{\alpha}^{\beta}f(x,u)\,\mathrm{d}u\right)\mathrm{d}x$ 收敛. 由此得到
\[
  \lim_{N\to+\infty}\int_{a}^{a+N}\left(\int_{\alpha}^{\beta}f(x,u)\,\mathrm{d}u\right)\mathrm{d}x=\lim_{A\to+\infty}\int_{a}^{A}\left(\int_{\alpha}^{\beta}f(x,u)\,\mathrm{d}u\right)\mathrm{d}x.
\]
\end{proof}

定理~\ref{thm:15.3.3} 中, 关于 $u$ 的积分区间 $[\alpha,\beta]$ 是有限的, 在很多的情况下, 往往需要知道两个无穷区间的积分是否可以交换. 对此, 有下面的定理.

\begin{theorem}[无穷积分号与无穷积分号可交换]\label{thm:15.3.4}
如果 $f(x,u)$ 满足:
\begin{enumerate}
  \item $f(x,u)$ 在 $[a,+\infty)\times[\alpha,+\infty)$ 上连续;
  \item 对任何固定的 $\beta>\alpha$, 无穷积分 $\displaystyle\int_{a}^{+\infty}f(x,u)\,\mathrm{d}x$ 关于 $u\in[\alpha,\beta]$ 一致收敛; 对任何固定的 $b>a$, 无穷积分 $\displaystyle\int_{\alpha}^{+\infty}f(x,u)\,\mathrm{d}u$ 关于 $x\in[a,b]$ 一致收敛;
  \item 无穷积分
  \[
    \int_{a}^{+\infty}\left(\int_{\alpha}^{+\infty}|f(x,u)|\,\mathrm{d}u\right)\mathrm{d}x,\quad\int_{\alpha}^{+\infty}\left(\int_{a}^{+\infty}|f(x,u)|\,\mathrm{d}x\right)\mathrm{d}u
\]
  中至少有一个收敛.
\end{enumerate}
则无穷积分
\[
  \int_{\alpha}^{+\infty}\left(\int_{a}^{+\infty}f(x,u)\,\mathrm{d}x\right)\mathrm{d}u,\quad\int_{a}^{+\infty}\left(\int_{\alpha}^{+\infty}f(x,u)\,\mathrm{d}u\right)\mathrm{d}x
\]
均收敛, 且
\[
  \int_{\alpha}^{+\infty}\left(\int_{a}^{+\infty}f(x,u)\,\mathrm{d}x\right)\mathrm{d}u=\int_{a}^{+\infty}\left(\int_{\alpha}^{+\infty}f(x,u)\,\mathrm{d}u\right)\mathrm{d}x.
\]
\end{theorem}

\begin{proof}
不失一般性, 设 $\displaystyle\int_{a}^{+\infty}\left(\int_{\alpha}^{+\infty}|f(x,u)|\,\mathrm{d}u\right)\mathrm{d}x$ 收敛. 根据无穷积分的 Cauchy 收敛准则, 有
\[
  \int_{a}^{+\infty}\left(\int_{\alpha}^{+\infty}f(x,u)\,\mathrm{d}u\right)\mathrm{d}x
\]
收敛. 于是
\[
  \begin{aligned}
    \int_{\alpha}^{+\infty}\left(\int_{a}^{+\infty}f(x,u)\,\mathrm{d}x\right)\mathrm{d}u&=\lim_{\beta\to+\infty}\int_{\alpha}^{\beta}\left(\int_{a}^{+\infty}f(x,u)\,\mathrm{d}x\right)\mathrm{d}u\\
    &\xlongequal[\text{定理~\ref{thm:15.3.3}}]{\text{条件 (2)}}\lim_{\beta\to+\infty}\int_{a}^{+\infty}\left(\int_{\alpha}^{\beta}f(x,u)\,\mathrm{d}u\right)\mathrm{d}x\\
    &=\lim_{\beta\to+\infty}\int_{a}^{+\infty}\left(\int_{\alpha}^{\beta}f(x,u)\,\mathrm{d}u\right)\mathrm{d}x\\
    &\quad+\lim_{\beta\to+\infty}\int_{a}^{+\infty}\left(\int_{\beta}^{+\infty}f(x,u)\,\mathrm{d}u\right)\mathrm{d}x\\
    &=\lim_{\beta\to+\infty}\int_{a}^{+\infty}\left[\int_{\alpha}^{\beta}f(x,u)\,\mathrm{d}u+\int_{\beta}^{+\infty}f(x,u)\,\mathrm{d}u\right]\mathrm{d}x\\
    &=\lim_{\beta\to+\infty}\int_{a}^{+\infty}\left(\int_{\alpha}^{+\infty}f(x,u)\,\mathrm{d}u\right)\mathrm{d}x\\
    &=\int_{a}^{+\infty}\left(\int_{\alpha}^{+\infty}f(x,u)\,\mathrm{d}u\right)\mathrm{d}x.
  \end{aligned}
\]
它等价于要证
\[
  \lim_{\beta\to+\infty}\int_{a}^{+\infty}\left(\int_{\beta}^{+\infty}f(x,u)\,\mathrm{d}u\right)\mathrm{d}x=0.
\]
事实上, $\forall\varepsilon>0$, 因为 $\displaystyle\int_{a}^{+\infty}\left(\int_{\alpha}^{+\infty}|f(x,u)|\,\mathrm{d}u\right)\mathrm{d}x$ 收敛, 故 $\exists b_{0}>a$, 使得
\[
  0\leqslant\int_{b_{0}}^{+\infty}\left(\int_{\alpha}^{+\infty}|f(x,u)|\,\mathrm{d}u\right)\mathrm{d}x<\frac{\varepsilon}{2}.
\]
对固定的 $b_{0}$, 由于 $\displaystyle\int_{\alpha}^{+\infty}f(x,u)\,\mathrm{d}u$ 关于 $x\in[a,b_{0}]$ 一致收敛, 故必 $\exists\beta_{0}>\alpha$, 使得当 $\beta>\beta_{0}$ 时, 有
\[
  \left|\int_{\beta}^{+\infty}f(x,u)\,\mathrm{d}u\right|<\frac{\varepsilon}{2(b_{0}-a)},\quad\forall x\in[a,b_{0}].
\]
于是, 当 $\beta>\beta_{0}$ 时, 有
\[
  \begin{aligned}
    \left|\int_{a}^{+\infty}\left(\int_{\beta}^{+\infty}f(x,u)\,\mathrm{d}u\right)\mathrm{d}x\right|&=\left|\int_{a}^{b_{0}}\left(\int_{\beta}^{+\infty}f(x,u)\,\mathrm{d}u\right)\mathrm{d}x+\int_{b_{0}}^{+\infty}\left(\int_{\beta}^{+\infty}f(x,u)\,\mathrm{d}u\right)\mathrm{d}x\right|\\
    &\leqslant\int_{a}^{b_{0}}\left|\int_{\beta}^{+\infty}f(x,u)\,\mathrm{d}u\right|\mathrm{d}x+\int_{b_{0}}^{+\infty}\left(\int_{\beta}^{+\infty}|f(x,u)|\,\mathrm{d}u\right)\mathrm{d}x\\
    &<\frac{\varepsilon}{2(b_{0}-a)}\cdot(b_{0}-a)+\int_{b_{0}}^{+\infty}\left(\int_{\alpha}^{+\infty}|f(x,u)|\,\mathrm{d}u\right)\mathrm{d}x\\
    &<\frac{\varepsilon}{2}+\frac{\varepsilon}{2}=\varepsilon.
  \end{aligned}
\]
这就证明了
\[
  \lim_{\beta\to+\infty}\int_{a}^{+\infty}\left(\int_{\beta}^{+\infty}f(x,u)\,\mathrm{d}u\right)\mathrm{d}x=0.
\]
\end{proof}

如果 $f(x,u)\geqslant0$, 利用无穷积分的 Dini 定理~\ref{thm:15.3.2}, 从定理~\ref{thm:15.3.4} 可得下面的结论.

\begin{theorem}[非负函数的无穷积分与无穷积分可交换]\label{thm:15.3.5}
如果 $f(x,u)$ 满足:
\begin{enumerate}
  \item $f(x,u)$ 在 $[a,+\infty)\times[\alpha,+\infty)$ 上连续且非负;
  \item 函数 $\displaystyle\int_{a}^{+\infty}f(x,u)\,\mathrm{d}x$ 在 $u\in[\alpha,+\infty)$ 上连续, $\displaystyle\int_{\alpha}^{+\infty}f(x,u)\,\mathrm{d}u$ 在 $x\in[a,+\infty)$ 上连续;
  \item 无穷积分
  \[
    \int_{\alpha}^{+\infty}\left(\int_{a}^{+\infty}f(x,u)\,\mathrm{d}x\right)\mathrm{d}u,\quad\int_{a}^{+\infty}\left(\int_{\alpha}^{+\infty}f(x,u)\,\mathrm{d}u\right)\mathrm{d}x
\]
  中至少有一个收敛.
\end{enumerate}
则条件 (3) 中另一个无穷积分也收敛, 且
\[
  \int_{\alpha}^{+\infty}\left(\int_{a}^{+\infty}f(x,u)\,\mathrm{d}x\right)\mathrm{d}u=\int_{a}^{+\infty}\left(\int_{\alpha}^{+\infty}f(x,u)\,\mathrm{d}u\right)\mathrm{d}x.
\]
\end{theorem}

\begin{proof}
由条件 (1)、(2) 及无穷积分的 Dini 定理~\ref{thm:15.3.2} 知:
$\forall\beta>\alpha$, $\displaystyle\int_{a}^{+\infty}f(x,u)\,\mathrm{d}x$ 关于 $u\in[\alpha,\beta]$ 一致收敛;
$\forall b>a$, $\displaystyle\int_{\alpha}^{+\infty}f(x,u)\,\mathrm{d}u$ 关于 $x\in[a,b]$ 一致收敛.
再由条件 (3) 及定理~\ref{thm:15.3.4} 可推得条件 (3) 中另一个无穷积分也收敛, 且
\[
  \int_{\alpha}^{+\infty}\left(\int_{a}^{+\infty}f(x,u)\,\mathrm{d}x\right)\mathrm{d}u=\int_{a}^{+\infty}\left(\int_{\alpha}^{+\infty}f(x,u)\,\mathrm{d}u\right)\mathrm{d}x.
\]
\end{proof}

由定理~\ref{thm:15.3.5} 立即得到.

\begin{corollary}[非负函数的无穷积分与无穷积分可交换]\label{cor:15.3.1}
如果 $f(x,u)$ 满足:
\begin{enumerate}
  \item $f(x,u)$ 在 $[a,+\infty)\times[\alpha,+\infty)$ 上连续且非负;
  \item 函数 $\displaystyle\int_{a}^{+\infty}f(x,u)\,\mathrm{d}x$ 在 $u\in[\alpha,+\infty)$ 上连续, $\displaystyle\int_{\alpha}^{+\infty}f(x,u)\,\mathrm{d}u$ 在 $x\in[a,+\infty)$ 上连续.
\end{enumerate}
则
\[
  \int_{\alpha}^{+\infty}\left(\int_{a}^{+\infty}f(x,u)\,\mathrm{d}x\right)\mathrm{d}u=\int_{a}^{+\infty}\left(\int_{\alpha}^{+\infty}f(x,u)\,\mathrm{d}u\right)\mathrm{d}x.
\]
此外, 如果无穷积分
\[
  \int_{\alpha}^{+\infty}\left(\int_{a}^{+\infty}f(x,u)\,\mathrm{d}x\right)\mathrm{d}u,\quad\int_{a}^{+\infty}\left(\int_{\alpha}^{+\infty}f(x,u)\,\mathrm{d}u\right)\mathrm{d}x
\]
中至少有一个收敛, 则上述另一个无穷积分也收敛.
如果上述无穷积分中有一个为 $+\infty$, 则另一个也为 $+\infty$, 即
\[
  \int_{\alpha}^{+\infty}\left(\int_{a}^{+\infty}f(x,u)\,\mathrm{d}x\right)\mathrm{d}u=\int_{a}^{+\infty}\left(\int_{\alpha}^{+\infty}f(x,u)\,\mathrm{d}u\right)\mathrm{d}x=+\infty.
\]
\end{corollary}

\begin{proof}
因为 $f(x,u)\geqslant0$, 所以有两种情形:
(1) 无穷积分中至少有一个为 $+\infty$, 则另一个必为 $+\infty$, 即
\[
  \int_{\alpha}^{+\infty}\left(\int_{a}^{+\infty}f(x,u)\,\mathrm{d}x\right)\mathrm{d}u=+\infty=\int_{a}^{+\infty}\left(\int_{\alpha}^{+\infty}f(x,u)\,\mathrm{d}u\right)\mathrm{d}x.
\]
(2) 无穷积分中至少有一个收敛, 则根据定理~\ref{thm:15.3.5} 知, 另一个无穷积分也收敛, 且
\[
  \int_{\alpha}^{+\infty}\left(\int_{a}^{+\infty}f(x,u)\,\mathrm{d}x\right)\mathrm{d}u=\int_{a}^{+\infty}\left(\int_{\alpha}^{+\infty}f(x,u)\,\mathrm{d}u\right)\mathrm{d}x.
\]
\end{proof}

现在来研究 $\displaystyle\int_{a}^{+\infty}f(x,u)\,\mathrm{d}x$ 的求导问题.

\begin{theorem}[无穷积分对参变量求导]\label{thm:15.3.6}
如果 $f(x,u)$ 满足:
\begin{enumerate}
  \item $f(x,u)$ 与 $\dfrac{\partial f(x,u)}{\partial u}$ 在 $[a,+\infty)\times[\alpha,\beta]$ 上连续;
  \item $\displaystyle\int_{a}^{+\infty}\frac{\partial f(x,u)}{\partial u}\,\mathrm{d}x$ 在 $u\in[\alpha,\beta]$ 上一致收敛;
  \item $\exists u_{0}\in[\alpha,\beta]$, 使 $\displaystyle\int_{a}^{+\infty}f(x,u_{0})\,\mathrm{d}x$ 收敛.
\end{enumerate}
则
\[
  \varphi(u)=\int_{a}^{+\infty}f(x,u)\,\mathrm{d}x
\]
在 $u\in[\alpha,\beta]$ 上可导, 且
\[
  \varphi'(u)=\frac{\mathrm{d}\varphi(u)}{\mathrm{d}u}=\frac{\mathrm{d}}{\mathrm{d}u}\left(\int_{a}^{+\infty}f(x,u)\,\mathrm{d}x\right)=\int_{a}^{+\infty}\frac{\partial f(x,u)}{\partial u}\,\mathrm{d}x.
\]
\end{theorem}

\begin{proof}
\textbf{证法 1} 设 $a_{n}(u)=\displaystyle\int_{a+n-1}^{a+n}f(x,u)\,\mathrm{d}x$, 显然由 $f(x,u)$ 连续知 $a_{n}(u)$ 为 $[\alpha,\beta]$ 上的连续函数, 再由 $\dfrac{\partial f(x,u)}{\partial u}$ 连续及定理~\ref{thm:15.1.3} 推得
\[
  a_{n}'(u)=\frac{\mathrm{d}}{\mathrm{d}u}\int_{a+n-1}^{a+n}f(x,u)\,\mathrm{d}x=\int_{a+n-1}^{a+n}\frac{\partial f(x,u)}{\partial u}\,\mathrm{d}x
\]
在 $[\alpha,\beta]$ 上连续. 从 $\displaystyle\int_{a}^{+\infty}\frac{\partial f(x,u)}{\partial u}\,\mathrm{d}x$ 在 $u\in[\alpha,\beta]$ 上一致收敛可推得
\[
  \sum_{n=1}^{\infty}a_{n}'(u)=\sum_{n=1}^{\infty}\int_{a+n-1}^{a+n}\frac{\partial f(x,u)}{\partial u}\,\mathrm{d}x
\]
在 $u\in[\alpha,\beta]$ 上一致收敛. 由条件 (3) 知 $\displaystyle\sum_{n=1}^{\infty}a_{n}(u_{0})=\sum_{n=1}^{\infty}\int_{a+n-1}^{a+n}f(x,u_{0})\,\mathrm{d}x=\int_{a}^{+\infty}f(x,u_{0})\,\mathrm{d}x$ 收敛. 根据函数项级数的逐项求导定理~\ref{thm:13.2.4p} 得到
\[
  \begin{aligned}
    \varphi'(u)&=\frac{\mathrm{d}}{\mathrm{d}u}\int_{a}^{+\infty}f(x,u)\,\mathrm{d}x=\frac{\mathrm{d}}{\mathrm{d}u}\sum_{n=1}^{\infty}\int_{a+n-1}^{a+n}f(x,u)\,\mathrm{d}x\\
    &\xlongequal{\text{定理~\ref{thm:13.2.4p}}}\sum_{n=1}^{\infty}\frac{\mathrm{d}}{\mathrm{d}u}\int_{a+n-1}^{a+n}f(x,u)\,\mathrm{d}x=\sum_{n=1}^{\infty}\int_{a+n-1}^{a+n}\frac{\partial f(x,u)}{\partial u}\,\mathrm{d}x\\
    &=\int_{a}^{+\infty}\frac{\partial f(x,u)}{\partial u}\,\mathrm{d}x.
  \end{aligned}
\]

\textbf{证法 2} 设 $\varphi_{n}(u)=\displaystyle\int_{a}^{a+n}f(x,u)\,\mathrm{d}x$, 则由条件 (1) 及定理~\ref{thm:15.1.3} 推得 $\varphi_{n}'(u)=\dfrac{\mathrm{d}}{\mathrm{d}u}\displaystyle\int_{a}^{a+n}f(x,u)\,\mathrm{d}x=\int_{a}^{a+n}\frac{\partial f(x,u)}{\partial u}\,\mathrm{d}x$ 在 $[\alpha,\beta]$ 上连续. 由条件 (2) 知 $\{\varphi_{n}'(u)\}$ 在 $[\alpha,\beta]$ 上一致收敛于 $\psi(u)=\displaystyle\int_{a}^{+\infty}\frac{\partial f(x,u)}{\partial u}\,\mathrm{d}x$. 再由条件 (3) 得到 $\{\varphi_{n}(u_{0})\}$ 收敛于 $\displaystyle\int_{a}^{+\infty}f(x,u_{0})\,\mathrm{d}x$. 根据定理~\ref{thm:13.2.4}, $\varphi(u)=\lim\limits_{n\to+\infty}\varphi_{n}(u)=\displaystyle\int_{a}^{+\infty}f(x,u)\,\mathrm{d}x$ 在 $u\in[\alpha,\beta]$ 上可导, 且
\[
  \begin{aligned}
    \varphi'(u)&=\frac{\mathrm{d}}{\mathrm{d}u}\left(\int_{a}^{+\infty}f(x,u)\,\mathrm{d}x\right)=\lim_{n\to+\infty}\varphi_{n}'(u)=\lim_{n\to+\infty}\int_{a}^{a+n}\frac{\partial f(x,u)}{\partial u}\,\mathrm{d}x\\
    &=\int_{a}^{+\infty}\frac{\partial f(x,u)}{\partial u}\,\mathrm{d}x.
  \end{aligned}
\]
\end{proof}

\begin{example}\label{ex:15.3.1}
讨论函数
\[
  \varphi(\alpha)=\int_{0}^{+\infty}\frac{\arctan x}{x^{\alpha}(2+x^{3})}\,\mathrm{d}x
\]
的连续性.
\end{example}

\begin{proof}[解]
因为
\[
  \frac{\arctan x}{x^{\alpha}(2+x^{3})}\sim\frac{1}{2}\frac{1}{x^{\alpha-1}}\quad(x\to0^{+}),
\]
所以, 当 $\alpha-1<1$, 即 $\alpha<2$ 时, 积分 $\displaystyle\int_{0}^{1}\frac{\arctan x}{x^{\alpha}(2+x^{3})}\,\mathrm{d}x$ 收敛; 又因为
\[
  \frac{\arctan x}{x^{\alpha}(2+x^{3})}\sim\frac{\pi}{2}\frac{1}{x^{\alpha+3}}\quad(x\to+\infty),
\]
故当 $\alpha+3>1$, 即 $\alpha>-2$ 时, 积分 $\displaystyle\int_{1}^{+\infty}\frac{\arctan x}{x^{\alpha}(2+x^{3})}\,\mathrm{d}x$ 收敛. 由此得到, $\varphi(\alpha)$ 的定义域应为 $(-2,2)$.

下面证明 $\varphi(\alpha)$ 在 $(-2,2)$ 上连续. 为此, 对任意 $[a,b]\subset(-2,2)$, 先证 $\displaystyle\int_{0}^{+\infty}\frac{\arctan x}{x^{\alpha}(2+x^{3})}\,\mathrm{d}x$ 在 $\alpha\in[a,b]$ 上一致收敛. 根据定理~\ref{thm:15.3.1}, $\varphi(\alpha)$ 在 $[a,b]$ 上连续. 再由于 $[a,b]$ 的任取性, $\varphi(\alpha)$ 在 $\alpha\in(-2,2)$ 上连续.

事实上, 当 $x\in(0,1)$ 时, 设 $a\leqslant\alpha\leqslant b<2$, 则存在常数 $c$, 使得
\[
  0<\frac{\arctan x}{x^{\alpha}(2+x^{3})}\leqslant\frac{c}{x^{\alpha-1}}\leqslant\frac{c}{x^{b-1}},
\]
而 $b-1<1$, 故由 Weierstrass 判别法知, 积分 $\displaystyle\int_{0}^{1}\frac{\arctan x}{x^{\alpha}(2+x^{3})}\,\mathrm{d}x$ 在 $\alpha\in(-\infty,b]$ 上一致收敛.

当 $x\in[1,+\infty)$ 时, 设 $-2<a\leqslant\alpha$, 则
\[
  0<\frac{\arctan x}{x^{\alpha}(2+x^{3})}\leqslant\frac{\pi}{2}\frac{1}{x^{\alpha+3}}\leqslant\frac{\pi}{2}\frac{1}{x^{a+3}},
\]
而 $a+3>1$, 故由 Weierstrass 判别法知, 积分 $\displaystyle\int_{1}^{+\infty}\frac{\arctan x}{x^{\alpha}(2+x^{3})}\,\mathrm{d}x$ 在 $\alpha\in[a,+\infty)$ 上一致收敛.

将两个积分合在一起, 即知 $\displaystyle\int_{0}^{+\infty}\frac{\arctan x}{x^{\alpha}(2+x^{3})}\,\mathrm{d}x$ 在 $\alpha\in[a,b]\subset(-2,2)$ 上一致收敛.
\end{proof}

\begin{example}\label{ex:15.3.2}
证明: $\displaystyle\int_{0}^{+\infty}\frac{\mathrm{e}^{-ax}-\mathrm{e}^{-bx}}{x}\,\mathrm{d}x=\ln\frac{b}{a}$ ($b>a>0$).
\end{example}

\begin{proof}
\textbf{证法 1} 在\extref{例 6.5.19} 中, 取 $f(x)=\mathrm{e}^{-x}$, 则 $k=\lim\limits_{x\to+\infty}f(x)=\lim\limits_{x\to+\infty}\mathrm{e}^{-x}=0$, $f(0)=1$. 于是
\[
  \int_{0}^{+\infty}\frac{\mathrm{e}^{-ax}-\mathrm{e}^{-bx}}{x}\,\mathrm{d}x=[f(0)-k]\ln\frac{b}{a}=(1-0)\ln\frac{b}{a}=\ln\frac{b}{a}.
\]

\textbf{证法 2}
\[
  \begin{aligned}
    \int_{0}^{+\infty}\frac{\mathrm{e}^{-ax}-\mathrm{e}^{-bx}}{x}\,\mathrm{d}x&=\int_{0}^{+\infty}\left(\int_{a}^{b}\mathrm{e}^{-ux}\,\mathrm{d}u\right)\mathrm{d}x\xlongequal[\text{积分可交换}]{\text{定理~\ref{thm:15.3.3}}}\int_{a}^{b}\left(\int_{0}^{+\infty}\mathrm{e}^{-ux}\,\mathrm{d}x\right)\mathrm{d}u\\
    &=\int_{a}^{b}-\frac{\mathrm{e}^{-ux}}{u}\bigg|_{x=0}^{+\infty}\mathrm{d}u=\int_{a}^{b}\frac{\mathrm{d}u}{u}=\ln u\bigg|_{a}^{b}=\ln\frac{b}{a}.
  \end{aligned}
\]
在上述计算中, 关键的是将
\[
  \frac{\mathrm{e}^{-ax}-\mathrm{e}^{-bx}}{x}
\]
表示成积分 $\displaystyle\int_{a}^{b}\mathrm{e}^{-ux}\,\mathrm{d}u$. 然后大胆地 (千万不要缩手缩脚!) 交换积分号. 最后计算出其结果为 $\ln\dfrac{b}{a}$. 但是, 还必须验证积分确实可交换.

事实上, 因为 $f(x,u)=\mathrm{e}^{-ux}$ 在 $[0,+\infty)\times[a,b]$ 上连续, 又
\[
  |\mathrm{e}^{-ux}|\leqslant\mathrm{e}^{-ax},\quad\forall(x,u)\in[0,+\infty)\times[a,b],
\]
且 $\displaystyle\int_{0}^{+\infty}\mathrm{e}^{-ax}\,\mathrm{d}x$ 收敛. 根据 Weierstrass 判别法知, $\displaystyle\int_{0}^{+\infty}\mathrm{e}^{-ux}\,\mathrm{d}x$ 在 $u\in[a,b]$ 中一致收敛, 从而有
\[
  \int_{0}^{+\infty}\left(\int_{a}^{b}\mathrm{e}^{-ux}\,\mathrm{d}u\right)\mathrm{d}x=\int_{a}^{b}\left(\int_{0}^{+\infty}\mathrm{e}^{-ux}\,\mathrm{d}x\right)\mathrm{d}u.
\]

\textbf{证法 3} 视 $a$ 为参数, 令
\[
  I(a)=\int_{0}^{+\infty}\frac{\mathrm{e}^{-ax}-\mathrm{e}^{-bx}}{x}\,\mathrm{d}x,
\]
则
\[
  \begin{aligned}
    I'(a)&=\frac{\mathrm{d}}{\mathrm{d}a}\int_{0}^{+\infty}\frac{\mathrm{e}^{-ax}-\mathrm{e}^{-bx}}{x}\,\mathrm{d}x\xlongequal[\text{积分号下求导}]{\text{定理~\ref{thm:15.3.6}}}\int_{0}^{+\infty}\frac{\partial}{\partial a}\left(\frac{\mathrm{e}^{-ax}-\mathrm{e}^{-bx}}{x}\right)\mathrm{d}x\\
    &=-\int_{0}^{+\infty}\mathrm{e}^{-ax}\,\mathrm{d}x=\frac{1}{a}\mathrm{e}^{-ax}\bigg|_{0}^{+\infty}=-\frac{1}{a},
  \end{aligned}
\]
\[
  I(a)=\int_{b}^{a}I'(t)\,\mathrm{d}t+I(b)=-\int_{b}^{a}\frac{\mathrm{d}t}{t}+0=-\ln t\bigg|_{b}^{a}=\ln\frac{b}{a}.
\]
由于
\[
  f(x,a)=\begin{cases}
    \dfrac{\mathrm{e}^{-ax}-\mathrm{e}^{-bx}}{x}, & x\neq0,\\
    b-a, & x=0,
\end{cases}
\]
与
\[
  \frac{\partial f(x,a)}{\partial a}=\begin{cases}
    -\mathrm{e}^{-ax}, & x\neq0,\\
    -1, & x=0
\end{cases}=-\mathrm{e}^{-ax}
\]
在 $[0,+\infty)\times[0,+\infty)$ 上都连续. 另一方面, 任取 $\delta>0$, 有
\[
  \left|\frac{\partial f(x,a)}{\partial a}\right|=|-\mathrm{e}^{-ax}|\leqslant\mathrm{e}^{-\delta x},\quad\forall x\in[0,+\infty),\forall a\in[\delta,+\infty).
\]
由于 $\displaystyle\int_{0}^{+\infty}\mathrm{e}^{-\delta x}\,\mathrm{d}x$ 收敛, 故根据 Weierstrass 判别法知, 无穷积分 $\displaystyle\int_{0}^{+\infty}\frac{\partial f(x,a)}{\partial a}\,\mathrm{d}x=\int_{0}^{+\infty}(-\mathrm{e}^{-ax})\,\mathrm{d}x$ 在 $a\in[\delta,+\infty)$ 上一致收敛. 此外, 显然 $f(x,a)$ 连续, 且由
\[
  \lim_{x\to+\infty}\frac{f(x,a)}{\frac{1}{x^{2}}}=\lim_{x\to+\infty}\frac{\mathrm{e}^{-ax}-\mathrm{e}^{-bx}}{\frac{1}{x}}=\lim_{x\to+\infty}\left(\frac{x}{\mathrm{e}^{ax}}-\frac{x}{\mathrm{e}^{bx}}\right)=0-0=0
\]
以及 $\displaystyle\int_{1}^{+\infty}\frac{\mathrm{d}x}{x^{2}}$ 收敛立知, $\displaystyle\int_{0}^{+\infty}f(x,a)\,\mathrm{d}x=\int_{0}^{+\infty}\frac{\mathrm{e}^{-ax}-\mathrm{e}^{-bx}}{x}\,\mathrm{d}x$ 收敛. 综合以上讨论及应用定理~\ref{thm:15.3.6} 得到
\[
  \frac{\mathrm{d}}{\mathrm{d}a}\int_{0}^{+\infty}\frac{\mathrm{e}^{-ax}-\mathrm{e}^{-bx}}{x}\,\mathrm{d}x=\int_{0}^{+\infty}\frac{\partial}{\partial a}\left(\frac{\mathrm{e}^{-ax}-\mathrm{e}^{-bx}}{x}\right)\mathrm{d}x,\quad\forall a\in[\delta,+\infty).
\]
由于 $\delta>0$ 任取, 故上式在 $a\in(0,+\infty)$ 均成立.
\end{proof}

更一般地, 有下面的例题.

\begin{example}\label{ex:15.3.3}
证明: $\displaystyle\int_{0}^{+\infty}\frac{\mathrm{e}^{-ax}-\mathrm{e}^{-bx}}{x}\cos mx\,\mathrm{d}x=\frac{1}{2}\ln\frac{b^{2}+m^{2}}{a^{2}+m^{2}}$ ($b>a>0$). 当 $m=0$ 时, 就是
\[
  \int_{0}^{+\infty}\frac{\mathrm{e}^{-ax}-\mathrm{e}^{-bx}}{x}\,\mathrm{d}x=\ln\frac{b}{a}.
\]
\end{example}

\begin{proof}
\textbf{证法 1}
\[
  \begin{aligned}
    \int_{0}^{+\infty}\frac{\mathrm{e}^{-ax}-\mathrm{e}^{-bx}}{x}\cos mx\,\mathrm{d}x&=\int_{0}^{+\infty}\left(\int_{a}^{b}\mathrm{e}^{-ux}\cos mx\,\mathrm{d}u\right)\mathrm{d}x\\
    &\xlongequal[\text{积分可交换}]{\text{定理~\ref{thm:15.3.3}}}\int_{a}^{b}\left(\int_{0}^{+\infty}\mathrm{e}^{-ux}\cos mx\,\mathrm{d}x\right)\mathrm{d}u\\
    &=\int_{a}^{b}\frac{u}{u^{2}+m^{2}}\,\mathrm{d}u\\
    &=\frac{1}{2}\ln(u^{2}+m^{2})\bigg|_{a}^{b}=\frac{1}{2}\ln\frac{b^{2}+m^{2}}{a^{2}+m^{2}},
  \end{aligned}
\]
其中积分可交换是因为
\[
  |\mathrm{e}^{-ux}\cos mx|\leqslant\mathrm{e}^{-ux}\leqslant\mathrm{e}^{-ax},
\]
而 $\displaystyle\int_{0}^{+\infty}\mathrm{e}^{-ax}\,\mathrm{d}x$ 收敛, 根据 Weierstrass 判别法知, $\displaystyle\int_{0}^{+\infty}\mathrm{e}^{-ux}\cos mx\,\mathrm{d}x$ 在 $u\in[a,b]$ 上一致收敛. 从而有
\[
  \int_{0}^{+\infty}\left(\int_{a}^{b}\mathrm{e}^{-ux}\cos mx\,\mathrm{d}u\right)\mathrm{d}x=\int_{a}^{b}\left(\int_{0}^{+\infty}\mathrm{e}^{-ux}\cos mx\,\mathrm{d}x\right)\mathrm{d}u.
\]

\textbf{证法 2} 视 $a$ 为参数, 令
\[
  I(a)=\int_{0}^{+\infty}\frac{\mathrm{e}^{-ax}-\mathrm{e}^{-bx}}{x}\cos mx\,\mathrm{d}x,
\]
则
\[
  \begin{aligned}
    I'(a)&=\frac{\mathrm{d}}{\mathrm{d}a}\int_{0}^{+\infty}\frac{\mathrm{e}^{-ax}-\mathrm{e}^{-bx}}{x}\cos mx\,\mathrm{d}x\\
    &\xlongequal[\text{积分号下求导}]{\text{定理~\ref{thm:15.3.6}}}\int_{0}^{+\infty}\frac{\partial}{\partial a}\left(\frac{\mathrm{e}^{-ax}-\mathrm{e}^{-bx}}{x}\cos mx\right)\mathrm{d}x\\
    &=-\int_{0}^{+\infty}\mathrm{e}^{-ax}\cos mx\,\mathrm{d}x=-\frac{a}{a^{2}+m^{2}},
  \end{aligned}
\]
\[
  \begin{aligned}
    I(a)&=\int_{b}^{a}I'(t)\,\mathrm{d}t+I(b)=-\int_{b}^{a}\frac{t}{t^{2}+m^{2}}\,\mathrm{d}t\\
    &=-\frac{1}{2}\ln(t^{2}+m^{2})\bigg|_{b}^{a}=\frac{1}{2}\ln\frac{b^{2}+m^{2}}{a^{2}+m^{2}}.
  \end{aligned}
\]
由于
\[
  f(x,a)=\begin{cases}
    \dfrac{\mathrm{e}^{-ax}-\mathrm{e}^{-bx}}{x}\cos mx, & x\neq0,\\
    b-a, & x=0,
\end{cases}
\]
\[
  \frac{\partial f(x,a)}{\partial a}=\begin{cases}
    -\mathrm{e}^{-ax}\cos mx, & x\neq0,\\
    -1, & x=0
\end{cases}=-\mathrm{e}^{-ax}\cos mx
\]
在 $[0,+\infty)\times[0,+\infty)$ 上连续. 另一方面, 任取 $\delta>0$, 有
\[
  \left|\frac{\partial f(x,a)}{\partial a}\right|=|-\mathrm{e}^{-ax}\cos mx|\leqslant\mathrm{e}^{-\delta x},\quad\forall x\in[0,+\infty),\forall a\in[\delta,+\infty).
\]
由于 $\displaystyle\int_{0}^{+\infty}\mathrm{e}^{-\delta x}\,\mathrm{d}x$ 收敛, 故根据 Weierstrass 判别法知, 无穷积分 $\displaystyle\int_{0}^{+\infty}\frac{\partial f(x,a)}{\partial a}\,\mathrm{d}x=\int_{0}^{+\infty}(-\mathrm{e}^{-ax}\cos mx)\,\mathrm{d}x$ 在 $a\in[\delta,+\infty)$ 上一致收敛. 此外, 显然 $f(x,a)$ 连续, 且由
\[
  \lim_{x\to+\infty}\frac{f(x,a)}{\frac{1}{x^{2}}}=\lim_{x\to+\infty}\frac{\mathrm{e}^{-ax}-\mathrm{e}^{-bx}}{\frac{1}{x}}\cos mx=\lim_{x\to+\infty}\left(\frac{x}{\mathrm{e}^{ax}}-\frac{x}{\mathrm{e}^{bx}}\right)\cos mx=0
\]
以及 $\displaystyle\int_{1}^{+\infty}\frac{\mathrm{d}x}{x^{2}}$ 收敛可知, $\displaystyle\int_{0}^{+\infty}f(x,a)\,\mathrm{d}x=\int_{0}^{+\infty}\frac{\mathrm{e}^{-ax}-\mathrm{e}^{-bx}}{x}\cos mx\,\mathrm{d}x$ 收敛. 综合以上讨论及应用定理~\ref{thm:15.3.6} 得到
\[
  \frac{\mathrm{d}}{\mathrm{d}a}\int_{0}^{+\infty}\frac{\mathrm{e}^{-ax}-\mathrm{e}^{-bx}}{x}\cos mx\,\mathrm{d}x=\int_{0}^{+\infty}\frac{\partial}{\partial a}\left(\frac{\mathrm{e}^{-ax}-\mathrm{e}^{-bx}}{x}\cos mx\right)\mathrm{d}x.
\]
\end{proof}

\begin{example}\label{ex:15.3.4}
证明: $\displaystyle\int_{0}^{+\infty}\mathrm{e}^{-ax^{2}}\cos bx\,\mathrm{d}x=\frac{1}{2}\sqrt{\frac{\pi}{a}}\mathrm{e}^{-\frac{b^{2}}{4a}}$ ($a>0$).

当 $b=0$ 时, $\displaystyle\int_{0}^{+\infty}\mathrm{e}^{-ax^{2}}\,\mathrm{d}x=\frac{1}{2}\sqrt{\frac{\pi}{a}}$, $\displaystyle\int_{0}^{+\infty}\mathrm{e}^{-x^{2}}\,\mathrm{d}x=\frac{\sqrt{\pi}}{2}$.
\end{example}

\begin{proof}
视 $b$ 为参数, 令 $I(b)=\displaystyle\int_{0}^{+\infty}\mathrm{e}^{-ax^{2}}\cos bx\,\mathrm{d}x$, $f(x,b)=\mathrm{e}^{-ax^{2}}\cos bx$ 在 $[0,+\infty)\times[0,+\infty)$ 中连续, 且
\[
  |f(x,b)|=|\mathrm{e}^{-ax^{2}}\cos bx|\leqslant\mathrm{e}^{-ax^{2}},\quad\forall b\in[0,+\infty),
\]
而 $\displaystyle\int_{0}^{+\infty}\mathrm{e}^{-ax^{2}}\,\mathrm{d}x$ 收敛, 根据 Weierstrass 判别法, $\displaystyle\int_{0}^{+\infty}\mathrm{e}^{-ax^{2}}\cos bx\,\mathrm{d}x$ 在 $b\in[0,+\infty)$ 上一致收敛. 根据定理~\ref{thm:15.3.1}, $I(b)$ 在 $[0,+\infty)$ 上连续.

又因为
\[
  \left|\frac{\partial f(x,b)}{\partial b}\right|=|-x\mathrm{e}^{-ax^{2}}\sin bx|\leqslant x\mathrm{e}^{-ax^{2}},
\]
而 $\displaystyle\int_{0}^{+\infty}x\mathrm{e}^{-ax^{2}}\,\mathrm{d}x$ 收敛, 根据 Weierstrass 判别法知, $\displaystyle\int_{0}^{+\infty}\frac{\partial f(x,b)}{\partial b}\,\mathrm{d}x=-\int_{0}^{+\infty}x\mathrm{e}^{-ax^{2}}\sin bx\,\mathrm{d}x$ 在 $b\in[0,+\infty)$ 中一致收敛. 应用无穷积分号下求导定理~\ref{thm:15.3.6} 得到
\[
  \begin{aligned}
    I'(b)&=-\int_{0}^{+\infty}x\mathrm{e}^{-ax^{2}}\sin bx\,\mathrm{d}x=\frac{1}{2a}\int_{0}^{+\infty}\sin bx\,\mathrm{d}\mathrm{e}^{-ax^{2}}\\
    &=\frac{1}{2a}\left(\mathrm{e}^{-ax^{2}}\sin bx\bigg|_{0}^{+\infty}-\int_{0}^{+\infty}b\mathrm{e}^{-ax^{2}}\cos bx\,\mathrm{d}x\right)\\
    &=-\frac{b}{2a}I(b),
  \end{aligned}
\]
\[
  (\ln|I(b)|)'=\frac{I'(b)}{I(b)}=-\frac{b}{2a},\quad\ln|I(b)|=\ln|I(0)|-\frac{b^{2}}{4a},
\]
\[
  I(b)=I(0)\mathrm{e}^{-\frac{b^{2}}{4a}}=\frac{1}{2}\sqrt{\frac{\pi}{a}}\mathrm{e}^{-\frac{b^{2}}{4a}},
\]
其中 $I(0)$
\[
  \int_{0}^{+\infty}\mathrm{e}^{-ax^{2}}\,\mathrm{d}x\xlongequal{t=\sqrt{a}x}\int_{0}^{+\infty}\mathrm{e}^{-t^{2}}\frac{\mathrm{d}t}{\sqrt{a}}=\frac{1}{\sqrt{a}}\cdot\frac{\sqrt{\pi}}{2}=\frac{1}{2}\sqrt{\frac{\pi}{a}}.
\]
\end{proof}

% ===== PART-B =====

\begin{example}\label{ex:15.3.5}
证明: $\displaystyle\int_{0}^{+\infty}\mathrm{e}^{-x^{2}}\,\mathrm{d}x=\frac{\sqrt{\pi}}{2}$.
\end{example}

\begin{proof}
证法 1$\sim$3 参阅\extref{例 6.5.7}, \extref{例 10.6.2}, 例~\ref{ex:15.2.8}.

\textbf{证法 4} 由 $I=\displaystyle\int_{0}^{+\infty}\mathrm{e}^{-x^{2}}\,\mathrm{d}x\xlongequal[u>0]{x=ut}\int_{0}^{+\infty}u\mathrm{e}^{-u^{2}t^{2}}\,\mathrm{d}t$ 得到
\[
  \begin{aligned}
    I^{2}&=I\int_{0}^{+\infty}\mathrm{e}^{-u^{2}}\,\mathrm{d}u=\int_{0}^{+\infty}I\mathrm{e}^{-u^{2}}\,\mathrm{d}u=\int_{0}^{+\infty}\left(\int_{0}^{+\infty}u\mathrm{e}^{-u^{2}t^{2}}\,\mathrm{d}t\right)\mathrm{e}^{-u^{2}}\,\mathrm{d}u\\
    &=\int_{0}^{+\infty}\left(\int_{0}^{+\infty}u\mathrm{e}^{-u^{2}(1+t^{2})}\,\mathrm{d}t\right)\mathrm{d}u\\
    &\xlongequal[\text{无穷积分号交换}]{\text{定理~\ref{thm:15.3.5}}}\int_{0}^{+\infty}\left(\int_{0}^{+\infty}u\mathrm{e}^{-u^{2}(1+t^{2})}\,\mathrm{d}u\right)\mathrm{d}t\\
    &=\int_{0}^{+\infty}-\frac{\mathrm{e}^{-u^{2}(1+t^{2})}}{2(1+t^{2})}\bigg|_{u=0}^{+\infty}\,\mathrm{d}t=\frac{1}{2}\int_{0}^{+\infty}\frac{\mathrm{d}t}{1+t^{2}}=\frac{1}{2}\arctan t\bigg|_{0}^{+\infty}=\frac{\pi}{4},
  \end{aligned}
\]
可得到
\[
  I=\frac{\sqrt{\pi}}{2}.
\]

接下来证无穷积分交换次序的合理性. 因为 $f(t,u)=u\mathrm{e}^{-u^{2}(1+t^{2})}$ 在 $[0,+\infty)\times[0,+\infty)$ 上连续且非负, 故
\[
  \varphi(u)=\int_{0}^{+\infty}u\mathrm{e}^{-u^{2}(1+t^{2})}\,\mathrm{d}t=\mathrm{e}^{-u^{2}}\int_{0}^{+\infty}\mathrm{e}^{-(ut)^{2}}\,\mathrm{d}(ut)=I\mathrm{e}^{-u^{2}}
\]
为 $[0,+\infty)$ 上的连续函数,
\[
  \psi(t)=\int_{0}^{+\infty}u\mathrm{e}^{-u^{2}(1+t^{2})}\,\mathrm{d}u=\frac{1}{2}\frac{1}{1+t^{2}}
\]
也为 $[0,+\infty)$ 上的连续函数. 且
\[
  \int_{0}^{+\infty}\varphi(u)\,\mathrm{d}u=\int_{0}^{+\infty}I\mathrm{e}^{-u^{2}}\,\mathrm{d}u=I\int_{0}^{+\infty}\mathrm{e}^{-u^{2}}\,\mathrm{d}u=I^{2}\quad(\text{正实数})
\]
收敛. 根据定理~\ref{thm:15.3.5}, 有
\[
  \int_{0}^{+\infty}\left(\int_{0}^{+\infty}u\mathrm{e}^{-u^{2}(1+t^{2})}\,\mathrm{d}t\right)\mathrm{d}u=\int_{0}^{+\infty}\left(\int_{0}^{+\infty}u\mathrm{e}^{-u^{2}(1+t^{2})}\,\mathrm{d}u\right)\mathrm{d}t.
\]
\end{proof}

\begin{example}\label{ex:15.3.6}
证明:
\[
  \int_{0}^{+\infty}\frac{\mathrm{e}^{-ax^{2}}-\mathrm{e}^{-bx^{2}}}{x^{2}}\,\mathrm{d}x=\sqrt{\pi}(\sqrt{b}-\sqrt{a}),\quad a>0,b>0.
\]
\end{example}

\begin{proof}
\textbf{证法 1} 视 $a$ 为参数, 令
\[
  I(a)=\int_{0}^{+\infty}\frac{\mathrm{e}^{-ax^{2}}-\mathrm{e}^{-bx^{2}}}{x^{2}}\,\mathrm{d}x,
\]
则
\[
  \begin{aligned}
    I'(a)&=\frac{\mathrm{d}}{\mathrm{d}a}\int_{0}^{+\infty}\frac{\mathrm{e}^{-ax^{2}}-\mathrm{e}^{-bx^{2}}}{x^{2}}\,\mathrm{d}x\xlongequal[\text{无穷积分号下求导}]{\text{定理~\ref{thm:15.3.6}}}\int_{0}^{+\infty}\frac{\partial}{\partial a}\left(\frac{\mathrm{e}^{-ax^{2}}-\mathrm{e}^{-bx^{2}}}{x^{2}}\right)\mathrm{d}x\\
    &=-\int_{0}^{+\infty}\mathrm{e}^{-ax^{2}}\,\mathrm{d}x\xlongequal{t=\sqrt{a}x}-\int_{0}^{+\infty}\mathrm{e}^{-t^{2}}\frac{\mathrm{d}t}{\sqrt{a}}=-\frac{\sqrt{\pi}}{2\sqrt{a}},
  \end{aligned}
\]
\[
  I(a)=\int_{b}^{a}I'(t)\,\mathrm{d}t+I(b)=-\int_{b}^{a}\frac{\sqrt{\pi}}{2\sqrt{t}}\,\mathrm{d}t+I(b)=-\sqrt{\pi}\sqrt{t}\bigg|_{b}^{a}+0=\sqrt{\pi}(\sqrt{b}-\sqrt{a}).
\]
现在再来验证无穷积分号下求导的合理性. 事实上, 由于
\[
  f(x,a)=\begin{cases}
    \dfrac{\mathrm{e}^{-ax^{2}}-\mathrm{e}^{-bx^{2}}}{x^{2}}, & x\neq0,\\
    b-a, & x=0
\end{cases}
\]
与
\[
  \frac{\partial f(x,a)}{\partial a}=\begin{cases}
    -\mathrm{e}^{-ax^{2}}, & x\neq0,\\
    -1, & x=0
\end{cases}=-\mathrm{e}^{-ax^{2}}
\]
在 $[0,+\infty)\times[0,+\infty)$ 上都连续. 另一方面, 任取 $\delta>0$, 有
\[
  \left|\frac{\partial f(x,a)}{\partial a}\right|=|-\mathrm{e}^{-ax^{2}}|\leqslant\mathrm{e}^{-\delta x^{2}},\quad\forall x\in[0,+\infty),\forall a\in[\delta,+\infty).
\]
由于 $\displaystyle\int_{0}^{+\infty}\mathrm{e}^{-\delta x^{2}}\,\mathrm{d}x$ 收敛, 故根据 Weierstrass 判别法知, 无穷积分 $\displaystyle\int_{0}^{+\infty}\frac{\partial f(x,a)}{\partial a}\,\mathrm{d}x=\int_{0}^{+\infty}(-\mathrm{e}^{-ax^{2}})\,\mathrm{d}x$ 在 $a\in[\delta,+\infty)$ 上一致收敛. 此外, 显然 $f(x,a)$ 连续, 且由
\[
  \lim_{x\to+\infty}\frac{f(x,a)}{\frac{1}{x^{2}}}=\lim_{x\to+\infty}(\mathrm{e}^{-ax^{2}}-\mathrm{e}^{-bx^{2}})=0
\]
以及 $\displaystyle\int_{1}^{+\infty}\frac{\mathrm{d}x}{x^{2}}$ 收敛立知, $\displaystyle\int_{0}^{+\infty}f(x,a)\,\mathrm{d}x=\int_{0}^{+\infty}\frac{\mathrm{e}^{-ax^{2}}-\mathrm{e}^{-bx^{2}}}{x^{2}}\,\mathrm{d}x$ 收敛. 综上及应用定理~\ref{thm:15.3.6} 得到
\[
  \frac{\mathrm{d}}{\mathrm{d}a}\int_{0}^{+\infty}\frac{\mathrm{e}^{-ax^{2}}-\mathrm{e}^{-bx^{2}}}{x^{2}}\,\mathrm{d}x=\int_{0}^{+\infty}\frac{\partial}{\partial a}\left(\frac{\mathrm{e}^{-ax^{2}}-\mathrm{e}^{-bx^{2}}}{x^{2}}\right)\mathrm{d}x,\quad\forall a\in[\delta,+\infty).
\]
由于 $\delta>0$ 任取, 故上式在 $a\in(0,+\infty)$ 均成立.

\textbf{证法 2}
\[
  \begin{aligned}
    \int_{0}^{+\infty}\frac{\mathrm{e}^{-ax^{2}}-\mathrm{e}^{-bx^{2}}}{x^{2}}\,\mathrm{d}x&=\int_{0}^{+\infty}\left(\int_{a}^{b}\mathrm{e}^{-ux^{2}}\,\mathrm{d}u\right)\mathrm{d}x\xlongequal[\text{积分可交换}]{\text{定理~\ref{thm:15.3.3}}}\int_{a}^{b}\left(\int_{0}^{+\infty}\mathrm{e}^{-ux^{2}}\,\mathrm{d}x\right)\mathrm{d}u\\
    &\xlongequal{t=\sqrt{u}x}\int_{a}^{b}\left(\int_{0}^{+\infty}\mathrm{e}^{-t^{2}}\frac{\mathrm{d}t}{\sqrt{u}}\right)\mathrm{d}u=\int_{a}^{b}\frac{\sqrt{\pi}}{2\sqrt{u}}\,\mathrm{d}u=\sqrt{\pi u}\bigg|_{a}^{b}\\
    &=\sqrt{\pi}(\sqrt{b}-\sqrt{a}).
  \end{aligned}
\]
因为 $f(x,u)=\mathrm{e}^{-ux^{2}}$ 在 $[0,+\infty)\times[a,b]$ 上连续, 又因为
\[
  |\mathrm{e}^{-ux^{2}}|\leqslant\mathrm{e}^{-ax^{2}},\quad\forall(x,u)\in[0,+\infty)\times[a,b],
\]
且 $\displaystyle\int_{0}^{+\infty}\mathrm{e}^{-ax^{2}}\,\mathrm{d}x$ 收敛, 根据 Weierstrass 判别法知, $\displaystyle\int_{0}^{+\infty}\mathrm{e}^{-ux^{2}}\,\mathrm{d}x$ 在 $u\in[a,b]$ 中一致收敛 (不失一般性, 令 $a<b$), 从而有
\[
  \int_{0}^{+\infty}\left(\int_{a}^{b}\mathrm{e}^{-ux^{2}}\,\mathrm{d}u\right)\mathrm{d}x=\int_{a}^{b}\left(\int_{0}^{+\infty}\mathrm{e}^{-ux^{2}}\,\mathrm{d}x\right)\mathrm{d}u.
\]
\end{proof}

\begin{example}\label{ex:15.3.7}
证明: $\displaystyle\int_{0}^{+\infty}\mathrm{e}^{-\left(x^{2}+\frac{a^{2}}{x^{2}}\right)}\,\mathrm{d}x=\frac{\sqrt{\pi}}{2}\mathrm{e}^{-2a}$ ($a\geqslant0$).
\end{example}

\begin{proof}
\textbf{证法 1} 视 $a$ 为参数, 令 $\varphi(a)=\displaystyle\int_{0}^{+\infty}\mathrm{e}^{-\left(x^{2}+\frac{a^{2}}{x^{2}}\right)}\,\mathrm{d}x$ ($a\geqslant0$), 则
\[
  f(x,a)=\begin{cases}
    \mathrm{e}^{-\left(x^{2}+\frac{a^{2}}{x^{2}}\right)}, & x\neq0,\\
    0, & x=0,
\end{cases}\quad
  \frac{\partial f(x,a)}{\partial a}=\begin{cases}
    -\dfrac{2a}{x^{2}}\mathrm{e}^{-\left(x^{2}+\frac{a^{2}}{x^{2}}\right)}, & x\neq0,\\
    0, & x=0
\end{cases}
\]
在 $[0,+\infty)\times[0,+\infty)$ 上都连续. 因为
\[
  |f(x,a)|\leqslant\mathrm{e}^{-x^{2}},\quad\forall x\in[0,+\infty),\forall a\in[0,+\infty)
\]
及 $\displaystyle\int_{0}^{+\infty}\mathrm{e}^{-x^{2}}\,\mathrm{d}x$ 收敛, 故根据 Weierstrass 判别法知, $\displaystyle\int_{0}^{+\infty}f(x,a)\,\mathrm{d}x=\int_{0}^{+\infty}\mathrm{e}^{-\left(x^{2}+\frac{a^{2}}{x^{2}}\right)}\,\mathrm{d}x$ 收敛 ($a\geqslant0$). 又因为
\[
  \left|\frac{\partial f(x,a)}{\partial a}\right|=\left|-\frac{2a}{x^{2}}\mathrm{e}^{-\left(x^{2}+\frac{a^{2}}{x^{2}}\right)}\right|\leqslant\frac{2\beta}{x^{2}}\min\left\{\mathrm{e}^{-x^{2}},\mathrm{e}^{-\frac{\alpha^{2}}{x^{2}}}\right\},
\]
所以, $\displaystyle\int_{0}^{+\infty}\frac{\partial f(x,a)}{\partial a}\,\mathrm{d}x$ 在 $a\in[\alpha,\beta]$ 上一致收敛 ($\beta>\alpha>0$). 于是, 应用定理~\ref{thm:15.3.6} 得到
\[
  \begin{aligned}
    \varphi'(a)&=\frac{\mathrm{d}}{\mathrm{d}a}\int_{0}^{+\infty}\mathrm{e}^{-\left(x^{2}+\frac{a^{2}}{x^{2}}\right)}\,\mathrm{d}x=\int_{0}^{+\infty}\frac{\partial}{\partial a}\mathrm{e}^{-\left(x^{2}+\frac{a^{2}}{x^{2}}\right)}\,\mathrm{d}x\\
    &=-2a\int_{0}^{+\infty}\frac{1}{x^{2}}\mathrm{e}^{-\left(x^{2}+\frac{a^{2}}{x^{2}}\right)}\,\mathrm{d}x\xlongequal{x=\frac{a}{t}}-2a\int_{+\infty}^{0}\frac{t^{2}}{a^{2}}\mathrm{e}^{-\left(\frac{a^{2}}{t^{2}}+t^{2}\right)}\frac{-a}{t^{2}}\,\mathrm{d}t\\
    &=-2\int_{0}^{+\infty}\mathrm{e}^{-\left(t^{2}+\frac{a^{2}}{t^{2}}\right)}\,\mathrm{d}t=-2\varphi(a),
  \end{aligned}
\]
\[
  \frac{\varphi'(a)}{\varphi(a)}=-2,\quad\ln\varphi(a)=-2a+\ln\varphi(0),
\]
\[
  \varphi(a)=\varphi(0)\mathrm{e}^{-2a}=\frac{\sqrt{\pi}}{2}\mathrm{e}^{-2a}.
\]

\textbf{证法 2} 由
\[
  \begin{aligned}
    \varphi(a)&=\int_{0}^{+\infty}\mathrm{e}^{-\left(x^{2}+\frac{a^{2}}{x^{2}}\right)}\,\mathrm{d}x\xlongequal[a>0]{x=\frac{a}{t}}\int_{+\infty}^{0}\mathrm{e}^{-\left(\frac{a^{2}}{t^{2}}+t^{2}\right)}\frac{-a}{t^{2}}\,\mathrm{d}t\\
    &=\int_{0}^{+\infty}\frac{a}{x^{2}}\mathrm{e}^{-\left(x^{2}+\frac{a^{2}}{x^{2}}\right)}\,\mathrm{d}x,
  \end{aligned}
\]
得到
\[
  \begin{aligned}
    2\varphi(a)&=\int_{0}^{+\infty}\left(1+\frac{a}{x^{2}}\right)\mathrm{e}^{-\left(x^{2}+\frac{a^{2}}{x^{2}}\right)}\,\mathrm{d}x=\mathrm{e}^{-2a}\int_{0}^{+\infty}\mathrm{e}^{-\left(x-\frac{a}{x}\right)^{2}}\,\mathrm{d}\left(x-\frac{a}{x}\right)\\
    &\xlongequal{t=x-\frac{a}{x}}\mathrm{e}^{-2a}\int_{-\infty}^{+\infty}\mathrm{e}^{-t^{2}}\,\mathrm{d}t=2\mathrm{e}^{-2a}\int_{0}^{+\infty}\mathrm{e}^{-t^{2}}\,\mathrm{d}t=2\mathrm{e}^{-2a}\cdot\frac{\sqrt{\pi}}{2},
  \end{aligned}
\]
即 $\varphi(a)=\dfrac{\sqrt{\pi}}{2}\mathrm{e}^{-2a}$.
\end{proof}

\begin{example}\label{ex:15.3.8}
证明:
\[
  \int_{0}^{+\infty}\frac{\sin\beta x}{x}\,\mathrm{d}x=\begin{cases}
    \dfrac{\pi}{2}, & \beta>0,\\
    0, & \beta=0,\\
    -\dfrac{\pi}{2}, & \beta<0.
\end{cases}
\]
\end{example}

\begin{proof}
用 Dirichlet 判别法, 由\extref{例 6.5.12} 知, $\displaystyle\int_{0}^{+\infty}\frac{\sin x}{x}\,\mathrm{d}x$ 是收敛的 (但非绝对收敛, 故条件收敛). 为计算它的值, 应引进收敛因子 $\mathrm{e}^{-\alpha x}$, 考虑含参变量 $\alpha$ 的积分
\[
  I(\alpha)=\int_{0}^{+\infty}\mathrm{e}^{-\alpha x}\frac{\sin x}{x}\,\mathrm{d}x.
\]
例~\ref{ex:15.2.5} 已经证明了这个积分在 $[0,+\infty)$ 上一致收敛. 而被积函数
\[
  f(x,\alpha)=\begin{cases}
    \mathrm{e}^{-\alpha x}\dfrac{\sin x}{x}, & x\neq0,\\
    1, & x=0
\end{cases}
\]
在 $[0,+\infty)\times[0,+\infty)$ 上连续, 根据定理~\ref{thm:15.3.1}, $I(\alpha)$ 在 $[0,+\infty)$ 上连续.

由于当 $\alpha\geqslant\delta>0$ 时, 有
\[
  \left|\frac{\partial f(x,\alpha)}{\partial\alpha}\right|=|-\mathrm{e}^{-\alpha x}\sin x|\leqslant\mathrm{e}^{-\delta x},
\]
且 $\displaystyle\int_{0}^{+\infty}\mathrm{e}^{-\delta x}\,\mathrm{d}x$ 收敛, 由 Weierstrass 判别法知道, 积分
\[
  \int_{0}^{+\infty}\frac{\partial f(x,\alpha)}{\partial\alpha}\,\mathrm{d}x=-\int_{0}^{+\infty}\mathrm{e}^{-\alpha x}\sin x\,\mathrm{d}x
\]
在 $\alpha\in[\delta,+\infty)$ 上一致收敛. 于是, 由定理~\ref{thm:15.3.6} 得到
\[
  I'(\alpha)=-\int_{0}^{+\infty}\mathrm{e}^{-\alpha x}\sin x\,\mathrm{d}x\xlongequal{\text{两次分部积分}}-\frac{1}{1+\alpha^{2}},\quad0<\alpha<+\infty,
\]
\[
  I(\alpha)=-\arctan\alpha+c.
\]
由于
\[
  |I(\alpha)|=\left|\int_{0}^{+\infty}\mathrm{e}^{-\alpha x}\frac{\sin x}{x}\,\mathrm{d}x\right|\leqslant\int_{0}^{+\infty}\mathrm{e}^{-\alpha x}\,\mathrm{d}x=\frac{1}{\alpha},
\]
故 $\lim\limits_{\alpha\to+\infty}I(\alpha)=0$, 且
\[
  c=\lim_{\alpha\to+\infty}c=\lim_{\alpha\to+\infty}(I(\alpha)+\arctan\alpha)=\lim_{\alpha\to+\infty}I(\alpha)+\frac{\pi}{2}=0+\frac{\pi}{2}=\frac{\pi}{2},
\]
\[
  I(\alpha)=-\arctan\alpha+\frac{\pi}{2},
\]
\[
  \int_{0}^{+\infty}\frac{\sin x}{x}\,\mathrm{d}x=I(0)=0+\frac{\pi}{2}=\frac{\pi}{2}.
\]
由此得到
\[
  \int_{0}^{+\infty}\frac{\sin\beta x}{x}\,\mathrm{d}x=\begin{cases}
    \displaystyle\int_{0}^{+\infty}\frac{\sin t}{t}\frac{\mathrm{d}t}{\beta}=\int_{0}^{+\infty}\frac{\sin t}{t}\,\mathrm{d}t=\frac{\pi}{2}, & \beta>0,\\
    0, & \beta=0,\\
    -\displaystyle\int_{0}^{+\infty}\frac{\sin|\beta|x}{x}\,\mathrm{d}x=-\frac{\pi}{2}, & \beta<0.
\end{cases}
\]
\end{proof}

\begin{example}[Fresnel 积分]\label{ex:15.3.9}
证明: Fresnel 积分
\[
  \int_{0}^{+\infty}\sin x^{2}\,\mathrm{d}x=\int_{0}^{+\infty}\cos x^{2}\,\mathrm{d}x=\frac{1}{2}\sqrt{\frac{\pi}{2}}.
\]
\end{example}

\begin{proof}
先计算 $\displaystyle\int_{0}^{+\infty}\sin x^{2}\,\mathrm{d}x$:
\[
  \begin{aligned}
    \int_{0}^{+\infty}\sin x^{2}\,\mathrm{d}x&\xlongequal[\mathrm{d}t=2x\,\mathrm{d}x]{t=x^{2}}\frac{1}{2}\int_{0}^{+\infty}\frac{\sin t}{\sqrt{t}}\,\mathrm{d}t=\frac{1}{2}\int_{0}^{+\infty}\sin t\left(\frac{2}{\sqrt{\pi}}\int_{0}^{+\infty}\mathrm{e}^{-tu^{2}}\,\mathrm{d}u\right)\mathrm{d}t\\
    &=\frac{1}{\sqrt{\pi}}\int_{0}^{+\infty}\left(\int_{0}^{+\infty}\mathrm{e}^{-tu^{2}}\sin t\,\mathrm{d}u\right)\mathrm{d}t\\
    &\xlongequal{\text{积分号交换}}\frac{1}{\sqrt{\pi}}\int_{0}^{+\infty}\left(\int_{0}^{+\infty}\mathrm{e}^{-tu^{2}}\sin t\,\mathrm{d}t\right)\mathrm{d}u\\
    &\xlongequal{\text{分部积分}}\frac{1}{\sqrt{\pi}}\int_{0}^{+\infty}\frac{\mathrm{d}u}{1+u^{4}}\xlongequal{\text{\extref{例 6.5.4}}}\frac{1}{\sqrt{\pi}}\cdot\frac{\pi}{2\sqrt{2}}=\frac{1}{2}\sqrt{\frac{\pi}{2}}.
  \end{aligned}
\]
积分值虽然算出来了, 但要证明积分次序交换的合法性却并不容易! 因此, 上面证明严格来说还未完成.

为了摆脱这一困境, 我们引进收敛因子 $\mathrm{e}^{-\alpha t}$ ($\alpha>0$). 考虑积分
\[
  \int_{0}^{+\infty}\mathrm{e}^{-\alpha t}\frac{\sin t}{\sqrt{t}}\,\mathrm{d}t,
\]
有
\[
  \begin{aligned}
    \int_{0}^{+\infty}\mathrm{e}^{-\alpha t}\frac{\sin t}{\sqrt{t}}\,\mathrm{d}t&=\frac{2}{\sqrt{\pi}}\int_{0}^{+\infty}\left(\int_{0}^{+\infty}\mathrm{e}^{-t(u^{2}+\alpha)}\sin t\,\mathrm{d}u\right)\mathrm{d}t\\
    &\xlongequal{\text{积分号交换}}\frac{2}{\sqrt{\pi}}\int_{0}^{+\infty}\left(\int_{0}^{+\infty}\mathrm{e}^{-t(u^{2}+\alpha)}\sin t\,\mathrm{d}t\right)\mathrm{d}u\\
    &\xlongequal{\text{两次分部积分}}\frac{2}{\sqrt{\pi}}\int_{0}^{+\infty}\frac{\mathrm{d}u}{1+(u^{2}+\alpha)^{2}},\quad\alpha>0.
  \end{aligned}
\]
由例~\ref{ex:15.2.3} 与例~\ref{ex:15.2.4}, 无穷积分 $\displaystyle\int_{0}^{+\infty}\mathrm{e}^{-t(u^{2}+\alpha)}\sin t\,\mathrm{d}t$ 在 $u\in[0,+\infty)$ 上一致收敛; 无穷积分 $\displaystyle\int_{0}^{+\infty}\mathrm{e}^{-t(u^{2}+\alpha)}\sin t\,\mathrm{d}u$ 在 $t\in[0,+\infty)$ 上一致收敛. 因为无穷积分
\[
  \int_{0}^{+\infty}\left(\int_{0}^{+\infty}\mathrm{e}^{-t(u^{2}+\alpha)}\,\mathrm{d}t\right)\mathrm{d}u=\int_{0}^{+\infty}\frac{\mathrm{e}^{-t(u^{2}+\alpha)}}{-(u^{2}+\alpha)}\Bigg|_{0}^{+\infty}\,\mathrm{d}u=\int_{0}^{+\infty}\frac{\mathrm{d}u}{u^{2}+\alpha}
\]
收敛, 故
\[
  \int_{0}^{+\infty}\left(\int_{0}^{+\infty}\mathrm{e}^{-t(u^{2}+\alpha)}|\sin t|\,\mathrm{d}t\right)\mathrm{d}u
\]
也收敛. 根据定理~\ref{thm:15.3.4} 知, 上述交换积分次序是合理的.

从
\[
  0<\frac{1}{1+(u^{2}+\alpha)^{2}}<\frac{1}{1+u^{4}}
\]
及 $\displaystyle\int_{0}^{+\infty}\frac{\mathrm{d}u}{1+u^{4}}$ 收敛, 并应用 Weierstrass 判别法知, $\displaystyle\int_{0}^{+\infty}\frac{\mathrm{d}u}{1+(u^{2}+\alpha)^{2}}$ 在 $\alpha\in[0,+\infty)$ 上一致收敛. 根据\extref{例 6.5.12}, $\displaystyle\int_{0}^{+\infty}\frac{\sin t}{\sqrt{t}}\,\mathrm{d}t$ 收敛, 当然在 $\alpha\in[0,+\infty)$ 上也是一致收敛的. 又因 $\mathrm{e}^{-\alpha t}$ 对 $t$ 单调, 且 $|\mathrm{e}^{-\alpha t}|\leqslant1$ 关于 $\alpha\in[0,+\infty)$ 一致有界, 根据 Abel 判别法知 $\displaystyle\int_{0}^{+\infty}\mathrm{e}^{-\alpha t}\frac{\sin t}{\sqrt{t}}\,\mathrm{d}t$ 在 $\alpha\in[0,+\infty)$ 上一致收敛. 于是
\[
  \begin{aligned}
    \int_{0}^{+\infty}\sin x^{2}\,\mathrm{d}x&\xlongequal{t=x^{2}}\frac{1}{2}\int_{0}^{+\infty}\frac{\sin t}{\sqrt{t}}\,\mathrm{d}t=\frac{1}{2}\int_{0}^{+\infty}\lim_{\alpha\to0^{+}}\mathrm{e}^{-\alpha t}\frac{\sin t}{\sqrt{t}}\,\mathrm{d}t\\
    &\xlongequal{\text{定理~\ref{thm:15.3.1}}}\frac{1}{2}\lim_{\alpha\to0^{+}}\int_{0}^{+\infty}\mathrm{e}^{-\alpha t}\frac{\sin t}{\sqrt{t}}\,\mathrm{d}t=\frac{1}{\sqrt{\pi}}\lim_{\alpha\to0^{+}}\int_{0}^{+\infty}\frac{\mathrm{d}u}{1+(u^{2}+\alpha)^{2}}\\
    &\xlongequal{\text{定理~\ref{thm:15.3.1}}}\frac{1}{\sqrt{\pi}}\int_{0}^{+\infty}\lim_{\alpha\to0^{+}}\frac{1}{1+(u^{2}+\alpha)^{2}}\,\mathrm{d}u=\frac{1}{\sqrt{\pi}}\int_{0}^{+\infty}\frac{\mathrm{d}u}{1+u^{4}}=\frac{1}{2}\sqrt{\frac{\pi}{2}}.
  \end{aligned}
\]

对于 $\displaystyle\int_{0}^{+\infty}\cos x^{2}\,\mathrm{d}x=\frac{1}{2}\sqrt{\frac{\pi}{2}}$, 能否类似 $\displaystyle\int_{0}^{+\infty}\sin x^{2}\,\mathrm{d}x=\frac{1}{2}\sqrt{\frac{\pi}{2}}$ 来证明呢? 应用 Dirichlet 判别法知
\[
  \int_{0}^{+\infty}\cos x^{2}\,\mathrm{d}x\xlongequal{t=x^{2}}\int_{0}^{+\infty}\frac{\cos t}{2\sqrt{t}}\,\mathrm{d}t
\]
是收敛的. 再乘收敛因子 $\mathrm{e}^{-\alpha t}$ 并应用 $\dfrac{1}{\sqrt{t}}=\dfrac{2}{\sqrt{\pi}}\displaystyle\int_{0}^{+\infty}\mathrm{e}^{-tu^{2}}\,\mathrm{d}u$ 得到
\[
  \begin{aligned}
    \int_{0}^{+\infty}\mathrm{e}^{-\alpha t}\frac{\cos t}{\sqrt{t}}\,\mathrm{d}t&=\frac{2}{\sqrt{\pi}}\int_{0}^{+\infty}\left(\int_{0}^{+\infty}\mathrm{e}^{-t(u^{2}+\alpha)}\cos t\,\mathrm{d}u\right)\mathrm{d}t\\
    &\xlongequal{\text{积分号交换}}\frac{2}{\sqrt{\pi}}\int_{0}^{+\infty}\left(\int_{0}^{+\infty}\mathrm{e}^{-t(u^{2}+\alpha)}\cos t\,\mathrm{d}t\right)\mathrm{d}u\\
    &\xlongequal{\text{分部积分}}\frac{2}{\sqrt{\pi}}\int_{0}^{+\infty}\frac{u^{2}+\alpha}{1+(u^{2}+\alpha)^{2}}\,\mathrm{d}u.
  \end{aligned}
\]
因为 $\displaystyle\int_{0}^{+\infty}\frac{\cos t}{\sqrt{t}}\,\mathrm{d}t$ 收敛, $\mathrm{e}^{-\alpha t}$ 对固定的 $\alpha$ 关于 $t$ 单调且关于 $\alpha\in[0,+\infty)$ 一致有界 ($|\mathrm{e}^{-\alpha t}|\leqslant1$), 故由 Abel 判别法, $\displaystyle\int_{0}^{+\infty}\mathrm{e}^{-\alpha t}\frac{\cos t}{\sqrt{t}}\,\mathrm{d}t$ 关于 $\alpha\in[0,+\infty)$ 一致收敛. 再由
\[
  \left|\frac{u^{2}+\alpha}{1+(u^{2}+\alpha)^{2}}\right|\leqslant\frac{u^{2}+1}{1+u^{4}},\quad\alpha\in[0,1]
\]
知 $\displaystyle\int_{0}^{+\infty}\frac{u^{2}+\alpha}{1+(u^{2}+\alpha)^{2}}\,\mathrm{d}u$ 关于 $\alpha\in[0,1]$ 一致收敛. 于是
\[
  \begin{aligned}
    \int_{0}^{+\infty}\cos x^{2}\,\mathrm{d}x&=\frac{1}{2}\int_{0}^{+\infty}\frac{\cos t}{\sqrt{t}}\,\mathrm{d}t=\frac{1}{2}\int_{0}^{+\infty}\lim_{\alpha\to0^{+}}\mathrm{e}^{-\alpha t}\frac{\cos t}{\sqrt{t}}\,\mathrm{d}t\\
    &\xlongequal{\text{定理~\ref{thm:15.3.1}}}\frac{1}{2}\lim_{\alpha\to0^{+}}\int_{0}^{+\infty}\mathrm{e}^{-\alpha t}\frac{\cos t}{\sqrt{t}}\,\mathrm{d}t=\frac{1}{\sqrt{\pi}}\lim_{\alpha\to0^{+}}\int_{0}^{+\infty}\frac{u^{2}+\alpha}{1+(u^{2}+\alpha)^{2}}\,\mathrm{d}u\\
    &\xlongequal{\text{定理~\ref{thm:15.3.1}}}\frac{1}{\sqrt{\pi}}\int_{0}^{+\infty}\lim_{\alpha\to0^{+}}\frac{u^{2}+\alpha}{1+(u^{2}+\alpha)^{2}}\,\mathrm{d}u=\frac{1}{\sqrt{\pi}}\int_{0}^{+\infty}\frac{u^{2}}{1+u^{4}}\,\mathrm{d}u\\
    &=\frac{1}{\sqrt{\pi}}\left(\int_{0}^{+\infty}\frac{u^{2}+1}{1+u^{4}}\,\mathrm{d}u-\int_{0}^{+\infty}\frac{\mathrm{d}u}{1+u^{4}}\right)\\
    &=\frac{1}{\sqrt{\pi}}\left(\int_{0}^{+\infty}\frac{\mathrm{d}\left(u-\frac{1}{u}\right)}{\left(u-\frac{1}{u}\right)^{2}+2}-\frac{\pi}{2\sqrt{2}}\right)\\
    &=\frac{1}{\sqrt{\pi}}\left(\frac{1}{\sqrt{2}}\arctan\frac{u-\frac{1}{u}}{\sqrt{2}}\Bigg|_{0^{+}}^{+\infty}-\frac{\pi}{2\sqrt{2}}\right)=\frac{1}{2}\sqrt{\frac{\pi}{2}}.
  \end{aligned}
\]
需要验证积分号次序可交换的条件是否成立? 不幸, 回答是否定的!

易见, 由 $|\mathrm{e}^{-t(u^{2}+\alpha)}\cos t|\leqslant\mathrm{e}^{-t\alpha}$ ($\alpha>0$) 知, $\displaystyle\int_{0}^{+\infty}\mathrm{e}^{-t(u^{2}+\alpha)}\cos t\,\mathrm{d}t$ 关于 $u\in[0,+\infty)$ 是一致收敛的. 但是, $\displaystyle\int_{0}^{+\infty}\mathrm{e}^{-t(u^{2}+\alpha)}\cos t\,\mathrm{d}u$ 关于 $t\in[0,a]$ 是非一致收敛的. 事实上, (反证) 假设该无穷积分关于 $t\in[0,a]$ 一致收敛, 则对 $\varepsilon_{0}=1$, $\exists A_{0}=A_{0}(\varepsilon_{0})>0$, 当 $A>A_{0}$ 时, $\forall t\in[0,a]$, 有
\[
  \left|\int_{A}^{+\infty}\mathrm{e}^{-t(u^{2}+\alpha)}\cos t\,\mathrm{d}u\right|<\varepsilon_{0}=1.
\]
作变换 $\theta=\sqrt{t}u$, 得
\[
  \left|\int_{A}^{+\infty}\mathrm{e}^{-t(u^{2}+\alpha)}\cos t\,\mathrm{d}u\right|=\left|\int_{\sqrt{t}A}^{+\infty}\mathrm{e}^{-\theta^{2}-\alpha t}\frac{\cos t}{\sqrt{t}}\,\mathrm{d}\theta\right|=\left|\mathrm{e}^{-\alpha t}\frac{\cos t}{\sqrt{t}}\int_{\sqrt{t}A}^{+\infty}\mathrm{e}^{-\theta^{2}}\,\mathrm{d}\theta\right|<1.
\]
令 $t\to0^{+}$ 得到 $+\infty\leqslant1$, 矛盾.

为了正确地证明 $\displaystyle\int_{0}^{+\infty}\cos x^{2}\,\mathrm{d}x=\frac{1}{2}\sqrt{\frac{\pi}{2}}$, 我们应该进行更细致的估计.
\[
  \int_{0}^{+\infty}\left[\int_{\delta}^{+\infty}\mathrm{e}^{-t(u^{2}+\alpha)}|\cos t|\,\mathrm{d}t\right]\mathrm{d}u
\]
收敛. 于是
\[
  \begin{aligned}
    \int_{\delta}^{+\infty}\mathrm{e}^{-\alpha t}\frac{\cos t}{\sqrt{t}}\,\mathrm{d}t&=\frac{2}{\sqrt{\pi}}\int_{\delta}^{+\infty}\left[\int_{0}^{+\infty}\mathrm{e}^{-t(u^{2}+\alpha)}\cos t\,\mathrm{d}u\right]\mathrm{d}t\\
    &\xlongequal{\text{积分可交换}}\frac{2}{\sqrt{\pi}}\int_{0}^{+\infty}\left[\int_{\delta}^{+\infty}\mathrm{e}^{-t(u^{2}+\alpha)}\cos t\,\mathrm{d}t\right]\mathrm{d}u\\
    &=\frac{2}{\sqrt{\pi}}\int_{0}^{+\infty}\frac{(u^{2}+\alpha)\cos\delta-\sin\delta}{(u^{2}+\alpha)^{2}+1}\mathrm{e}^{-\delta(u^{2}+\alpha)}\,\mathrm{d}u.
  \end{aligned}
\]
从
\[
  \int_{0}^{+\infty}\frac{\mathrm{d}u}{(u^{2}+\alpha)^{2}+1}\quad\text{与}\quad\int_{0}^{+\infty}\frac{u^{2}+\alpha}{(u^{2}+\alpha)^{2}+1}\,\mathrm{d}u
\]
均收敛推得上式右边积分关于 $\delta\in[0,1]$ 是一致收敛的. 令 $\delta\to0^{+}$ 得到
\[
  \int_{0}^{+\infty}\mathrm{e}^{-\alpha t}\frac{\cos t}{\sqrt{t}}\,\mathrm{d}t=\frac{2}{\sqrt{\pi}}\int_{0}^{+\infty}\frac{u^{2}+\alpha}{(u^{2}+\alpha)^{2}+1}\,\mathrm{d}u.
\]
由此式及前面的论证结果, 立即有
\[
  \int_{0}^{+\infty}\cos x^{2}\,\mathrm{d}x=\frac{1}{2}\sqrt{\frac{\pi}{2}}.
\]
\end{proof}

\begin{exblock}[练习题 15.3]
\begin{enumerate}
  \item 研究下列函数在指定区间上的连续性:
  \begin{enumerate}
    \item $f(x)=\displaystyle\int_{0}^{+\infty}\frac{t}{2+t^{x}}\,\mathrm{d}t$, $x\in(2,+\infty)$;
    \item $\varphi(\alpha)=\displaystyle\int_{0}^{\pi}\frac{\sin x}{x^{\alpha}(\pi-x)^{\alpha}}\,\mathrm{d}x$, $\alpha\in(0,2)$;
    \item $f(\alpha)=\displaystyle\int_{1}^{+\infty}\frac{\sin x}{x^{\alpha}}\,\mathrm{d}x$, $\alpha\in(0,+\infty)$.
  \end{enumerate}
  \item 利用公式 $\displaystyle\int_{0}^{1}x^{\alpha-1}\,\mathrm{d}x=\frac{1}{\alpha}$ ($\alpha>0$), 计算积分 $\displaystyle\int_{0}^{1}x^{\alpha-1}(\ln x)^{m}\,\mathrm{d}x$, 其中 $m$ 为自然数.
  \item 计算参变量广义积分, 其中 $a>0$, $b>0$:
  \begin{enumerate}
    \item $\displaystyle\int_{0}^{+\infty}\frac{\mathrm{e}^{-ax}-\mathrm{e}^{-bx}}{x}\sin mx\,\mathrm{d}x$;
    \item $\displaystyle\int_{0}^{+\infty}\frac{\mathrm{e}^{-ax}-\mathrm{e}^{-bx}}{x}\cos mx\,\mathrm{d}x$.
  \end{enumerate}
  \item 计算参变量广义积分
  \[
    \int_{0}^{+\infty}\frac{\ln(\alpha^{2}+x^{2})}{\beta^{2}+x^{2}}\,\mathrm{d}x,\quad\beta\neq0.
\]
  \item 利用已知的积分值计算下列积分:
  \begin{enumerate}
    \item $\displaystyle\int_{-\infty}^{+\infty}\mathrm{e}^{-(2x^{2}+x+2)}\,\mathrm{d}x$;
    \item $\displaystyle\int_{0}^{+\infty}\frac{\sin x^{2}}{x}\,\mathrm{d}x$;
    \item $\displaystyle\int_{0}^{+\infty}\frac{\sin^{3}x}{x}\,\mathrm{d}x$;
    \item $\displaystyle\int_{0}^{+\infty}\frac{\sin ax\cos\beta x}{x}\,\mathrm{d}x$.
  \end{enumerate}
  \item 设 $\{f_{n}(x)\}$ 为 $[0,+\infty)$ 上的连续函数列, 满足:
  \begin{enumerate}
    \item 在 $[0,+\infty)$ 上, $|f_{n}(x)|\leqslant g(x)$, 且 $\displaystyle\int_{0}^{+\infty}g(x)\,\mathrm{d}x$ 收敛;
    \item 在任何有限区间 $[0,A]$ 上 ($A>0$), 函数列 $\{f_{n}(x)\}$ 一致收敛于 $f(x)$.
  \end{enumerate}
  证明:
  \[
    \lim_{n\to+\infty}\int_{0}^{+\infty}f_{n}(x)\,\mathrm{d}x=\int_{0}^{+\infty}f(x)\,\mathrm{d}x.
\]
  \item (1) 证明: $\displaystyle\int_{0}^{+\infty}\frac{\sin2x}{x+\alpha}\mathrm{e}^{-\alpha x}\,\mathrm{d}x$ 在 $\alpha\in[0,b]$ 上一致收敛 ($b>0$);\\
  (2) 求极限
  \[
    \lim_{\alpha\to0^{+}}\int_{0}^{+\infty}\frac{\sin2x}{x+\alpha}\mathrm{e}^{-\alpha x}\,\mathrm{d}x.
\]
\end{enumerate}
\end{exblock}

\begin{exblock}[思考题 15.3]
\begin{enumerate}
  \item 利用已知积分计算下列积分:
  \begin{enumerate}
    \item $\displaystyle\int_{0}^{+\infty}\frac{\sin^{4}x}{x^{2}}\,\mathrm{d}x$;
    \item $\displaystyle\int_{0}^{+\infty}\mathrm{e}^{-\left(x^{2}+\frac{a^{2}}{x^{2}}\right)}\,\mathrm{d}x$, $a>0$;
    \item $\displaystyle\int_{0}^{+\infty}\frac{\mathrm{e}^{-ax^{2}}-\cos bx}{x^{2}}\,\mathrm{d}x$, $a>0$;
    \item $\displaystyle\int_{0}^{+\infty}\frac{\sin^{4}\alpha x-\sin^{4}\beta x}{x}\,\mathrm{d}x$, $\alpha>0,\beta>0$.
  \end{enumerate}
  \item 证明: 函数
  \[
    f(\alpha)=\int_{0}^{+\infty}\frac{\mathrm{e}^{-t}}{|\sin t|^{\alpha}}\,\mathrm{d}t
\]
  在 $(0,1)$ 中连续.
  \item 设 $\displaystyle\int_{-\infty}^{+\infty}|f(x)|\,\mathrm{d}x$ 收敛, 证明: 函数
  \[
    \varphi(u)=\int_{-\infty}^{+\infty}f(x)\cos ux\,\mathrm{d}x
\]
  在 $(-\infty,+\infty)$ 上一致连续.
  \item 通过积分
  \[
    \varphi(u)=\int_{0}^{+\infty}u\mathrm{e}^{-ux}\,\mathrm{d}x,\quad0<u\leqslant1
\]
  与
  \[
    \psi(u)=\int_{0}^{+\infty}u\mathrm{e}^{u(u-x)}\,\mathrm{d}x,\quad1\leqslant u<+\infty,
\]
  说明含参变量广义积分的 Dini 定理~\ref{thm:15.3.2} 在开区间或无穷区间上不成立.
  \item 设
  \[
    g(\alpha)=\int_{1}^{+\infty}\frac{\arctan\alpha x}{x^{2}\sqrt{x^{2}-1}}\,\mathrm{d}x,
\]
  证明:
  \begin{enumerate}
    \item $g'(\alpha)=\dfrac{\pi}{2}\left(1-\dfrac{|\alpha|}{\sqrt{1+\alpha^{2}}}\right)$;
    \item $g(\alpha)=\dfrac{\pi}{2}(|\alpha|+1-\sqrt{1+\alpha^{2}})\operatorname{sgn}\alpha$ ($-\infty<\alpha<+\infty$).
  \end{enumerate}
  \item 设
  \[
    g(\alpha,\beta)=\int_{0}^{+\infty}\frac{\arctan\alpha x\arctan\beta x}{x^{2}}\,\mathrm{d}x,\quad\alpha>0,\beta>0,
\]
  证明:
  \begin{enumerate}
    \item $g_{\alpha\beta}''(\alpha,\beta)=\dfrac{\pi}{2(\alpha+\beta)}$, $\alpha>0,\beta>0$;
    \item $g_{\alpha}'(\alpha,\beta)=\dfrac{\pi}{2}\ln\dfrac{\alpha+\beta}{\alpha}$, $\alpha>0,\beta>0$;
    \item $g(\alpha,\beta)=\dfrac{\pi}{2}\ln\dfrac{(\alpha+\beta)^{\alpha+\beta}}{\alpha^{\alpha}\beta^{\beta}}$, $\alpha>0,\beta>0$;
    \item
    \[
      g(\alpha,\beta)=\begin{cases}
        \dfrac{\pi}{2}\operatorname{sgn}\alpha\cdot\operatorname{sgn}\beta\cdot\ln\dfrac{(|\alpha|+|\beta|)^{|\alpha|+|\beta|}}{|\alpha|^{|\alpha|}|\beta|^{|\beta|}}, & \alpha\beta\neq0,\\
        0, & \alpha\beta=0.
\end{cases}
\]
  \end{enumerate}
  \item 设 $F(x)=\mathrm{e}^{\frac{x^{2}}{2}}\displaystyle\int_{x}^{+\infty}\mathrm{e}^{-\frac{t^{2}}{2}}\,\mathrm{d}t$, $x\in[0,+\infty)$. 证明:
  \begin{enumerate}
    \item $\lim\limits_{x\to+\infty}F(x)=0$;
    \item $F(x)$ 在 $[0,+\infty)$ 内单调减少.
  \end{enumerate}
  \item 按下列步骤证明: Euler-Poisson 概率积分 $\displaystyle\int_{0}^{+\infty}\mathrm{e}^{-x^{2}}\,\mathrm{d}x=\frac{\sqrt{\pi}}{2}$.
  \begin{enumerate}
    \item $\displaystyle\int_{0}^{+\infty}\mathrm{e}^{-x^{2}}\,\mathrm{d}x=\int_{0}^{+\infty}\alpha\mathrm{e}^{-\alpha^{2}t^{2}}\,\mathrm{d}t$;
    \item 设 $I=\displaystyle\int_{0}^{+\infty}\mathrm{e}^{-\alpha^{2}}\,\mathrm{d}\alpha$, 则
    \[
      I^{2}=\int_{0}^{+\infty}\mathrm{d}\alpha\int_{0}^{+\infty}\alpha\mathrm{e}^{-\alpha^{2}(1+t^{2})}\,\mathrm{d}t=\frac{\pi}{4};
\]
    \item $\displaystyle\int_{0}^{+\infty}\mathrm{e}^{-x^{2}}\,\mathrm{d}x=I=\frac{\sqrt{\pi}}{2}$ (对照\extref{例 6.5.7}, \extref{例 10.6.2}).
  \end{enumerate}
  \item 应用等式 $\mathrm{e}^{-x^{2}}=\lim\limits_{n\to+\infty}\left(1+\dfrac{x^{2}}{n}\right)^{-n}$ 证明:
  \[
    \int_{0}^{+\infty}\mathrm{e}^{-x^{2}}\,\mathrm{d}x=\lim_{n\to+\infty}\sqrt{n}\int_{0}^{\frac{\pi}{2}}\sin^{2n-2}t\,\mathrm{d}t=\frac{\sqrt{\pi}}{2}.
\]
  \item 设
  \[
    \varphi(x)=\int_{0}^{1}\frac{\mathrm{e}^{-x^{2}(1+u^{2})}}{1+u^{2}}\,\mathrm{d}u,\quad f(x)=\left(\int_{0}^{x}\mathrm{e}^{-t^{2}}\,\mathrm{d}t\right)^{2},\quad I=\int_{0}^{+\infty}\mathrm{e}^{-x^{2}}\,\mathrm{d}x,
\]
  证明:
  \begin{enumerate}
    \item $f'(x)=-\varphi'(x)$, $f(x)=\dfrac{\pi}{4}-\varphi(x)$, $x\geqslant0$;
    \item $I^{2}=\lim\limits_{x\to+\infty}f(x)=\dfrac{\pi}{4}-\lim\limits_{x\to+\infty}\varphi(x)=\dfrac{\pi}{4}$, $I=\dfrac{\sqrt{\pi}}{2}$.
  \end{enumerate}
  \item 设
  \[
    g(\alpha)=\int_{0}^{+\infty}\mathrm{e}^{-x^{2}}\cos2\alpha x\,\mathrm{d}x,
\]
  证明:
  \begin{enumerate}
    \item $g'(\alpha)=-2\alpha g(\alpha)$, $g(\alpha)=\dfrac{\sqrt{\pi}}{2}\mathrm{e}^{-\alpha^{2}}$ (已知 $g(0)=\displaystyle\int_{0}^{+\infty}\mathrm{e}^{-x^{2}}\,\mathrm{d}x=\frac{\sqrt{\pi}}{2}$);
    \item $g(\alpha)$ 有任意阶导数, 因而 $g(\alpha)$ 是 $C^{\infty}$ 函数, 且
    \[
      \int_{0}^{+\infty}x^{2k}\mathrm{e}^{-x^{2}}\cos2\alpha x\,\mathrm{d}x=(-1)^{k}\frac{\sqrt{\pi}}{2^{2k+1}}\frac{\mathrm{d}^{2k}}{\mathrm{d}\alpha^{2k}}\mathrm{e}^{-\alpha^{2}}.
\]
  \end{enumerate}
\end{enumerate}
\end{exblock}

% ====================================================================
%  徐森林、金亚东、薛春华《数学分析》（第三册）
%  \S15.4  Gamma 函数与 B 函数
%  源：原书印刷页 222 起（PDF 物理页 228 起）
%  底本：PaddleOCR-VL-1.6，见 chapters/_ch15_15_4_raw.md
%  约定：全西文标点；行间公式用一对方括号；多行公式用 align 或 aligned；
%        跨册引用一律用 extref（写作「定理 x.y.z」原文）；数学模式内的中文必须包text
% ====================================================================
\section{\texorpdfstring{$\Gamma$ 函数与 $\mathrm{B}$ 函数}{Gamma 函数与 B 函数}}\label{sec:15.4}

在\extref{例 6.5.4}证明 $\displaystyle\int_{0}^{+\infty}\frac{\mathrm{d}x}{1+x^{4}}=\frac{\pi}{2\sqrt{2}}$ 的证法 2 中, 我们借用了 $\Gamma$ 函数 $\displaystyle\Gamma(s)=\int_{0}^{+\infty}t^{s-1}\mathrm{e}^{-t}\,\mathrm{d}t$ 与 $\mathrm{B}$ 函数 $\displaystyle\mathrm{B}(p,q)=\int_{0}^{1}t^{p-1}(1-t)^{q-1}\,\mathrm{d}t$. 前者是含一个参变量 $s$ 的广义积分, 它确定的函数 $\Gamma(s)$ 的定义域为 $\{s\mid 1-s<1\}=\{s\mid s>0\}=(0,+\infty)$. 后者是含两个参变量 $p,q$ 的广义积分, 它确定的二元函数 $\mathrm{B}(p,q)$ 的定义域为 $\{(p,q)\mid 1-p<1, 1-q<1\}=\{(p,q)\mid p>0, q>0\}$.

\begin{theorem}\label{thm:15.4.1}
$\Gamma(s)$ 在 $(0,+\infty)$ 上连续, 且有各阶连续导数.
\end{theorem}

\begin{proof}
对任意 $\beta>\alpha>0$, 由于当 $s\in[\alpha,\beta]$ 时
\[
  0<t^{s-1}\mathrm{e}^{-t}\leqslant t^{\alpha-1}\mathrm{e}^{-t},\quad\forall t\in(0,1)
\]
及 $\displaystyle\int_{0}^{1}t^{\alpha-1}\mathrm{e}^{-t}\,\mathrm{d}t$ 收敛, 所以 $\displaystyle\int_{0}^{1}t^{s-1}\mathrm{e}^{-t}\,\mathrm{d}t$ 在 $s\in[\alpha,\beta]$ 上一致收敛, 因而 $\displaystyle\int_{0}^{1}t^{s-1}\mathrm{e}^{-t}\,\mathrm{d}t$ 在 $s\in[\alpha,\beta]$ 上为连续函数. 因为
\[
  0<t^{s-1}\mathrm{e}^{-t}\leqslant t^{\beta-1}\mathrm{e}^{-t},\quad\forall t\in[1,+\infty),\ \forall s\in[\alpha,\beta]
\]
及 $\displaystyle\int_{1}^{+\infty}t^{\beta-1}\mathrm{e}^{-t}\,\mathrm{d}t$ 收敛, 根据 Weierstrass 判别法知, $\displaystyle\int_{1}^{+\infty}t^{s-1}\mathrm{e}^{-t}\,\mathrm{d}t$ 在 $s\in[\alpha,\beta]$ 上一致收敛, 因而 $\displaystyle\int_{1}^{+\infty}t^{s-1}\mathrm{e}^{-t}\,\mathrm{d}t$ 为 $[\alpha,\beta]$ 上的连续函数. 由 $\beta>\alpha>0$ 的任意性知, $\Gamma(s)$ 在 $(0,+\infty)$ 上连续.

对 $k\in\mathbb{N}$, 令 $\alpha=2\delta$, 则由
\[
  \lim_{t\to 0^{+}}\left|\frac{(\ln t)^{k}}{\left(\frac{1}{t}\right)^\delta}\mathrm{e}^{-t}\right|=\lim_{u\to+\infty}\left|\frac{(-\ln u)^{k}}{u^\delta}\mathrm{e}^{-\frac{1}{u}}\right|=0,
\]
\[
  \left|\frac{(\ln t)^{k}}{\frac{1}{t}}\mathrm{e}^{-t}\right|\leqslant M_{k},\quad\forall t\in(0,+\infty), % 原书此处分母印为 1/t
\]
\[
  \left|t^{s-1}(\ln t)^{k}\mathrm{e}^{-t}\right|\leqslant\left|t^{\alpha-1}(\ln t)^{k}\mathrm{e}^{-t}\right|=\left|t^{\delta-1}\frac{(\ln t)^{k}}{\left(\frac{1}{t}\right)^{\delta}}\mathrm{e}^{-t}\right|\leqslant\frac{M_{k}}{t^{1-\delta}},\quad\forall s\in[\alpha,\beta],
\]
再由 $\displaystyle\int_{0}^{1}\frac{\mathrm{d}t}{t^{1-\delta}}$ 收敛及 Weierstrass 判别法推得 $\displaystyle\int_{0}^{1}t^{s-1}(\ln t)^{k}\mathrm{e}^{-t}\,\mathrm{d}t$ 在 $s\in[\alpha,\beta]$ 上一致收敛, 因而
\[
  \frac{\mathrm{d}^{l}}{\mathrm{d}s^{l}}\left(\int_{0}^{1}t^{s-1}\mathrm{e}^{-t}\,\mathrm{d}t\right)=\int_{0}^{1}\frac{\mathrm{d}^{l}}{\mathrm{d}s^{l}}t^{s-1}\mathrm{e}^{-t}\,\mathrm{d}t=\int_{0}^{1}t^{s-1}(\ln t)^{l}\mathrm{e}^{-t}\,\mathrm{d}t,\quad l=1,2,\cdots.
\]
因为
\[
  \lim_{t\to+\infty}\frac{\left|t^{\beta-1}(\ln t)^{k}\mathrm{e}^{-t}\right|}{\frac{1}{t^{2}}}=\lim_{t\to+\infty}\frac{\left|t^{\beta+1}(\ln t)^{k}\right|}{\mathrm{e}^{t}}=0,\quad\int_{1}^{+\infty}\frac{\mathrm{d}t}{t^{2}}\ \text{收敛},
\]
所以, 根据\extref{定理 6.5.3$'$}(2)知, $\displaystyle\int_{1}^{+\infty}\left|t^{\beta-1}(\ln t)^{k}\mathrm{e}^{-t}\right|\mathrm{d}t$ 收敛. 又因为
\[
  \left|t^{s-1}(\ln t)^{k}\mathrm{e}^{-t}\right|\leqslant\left|t^{\beta-1}(\ln t)^{k}\mathrm{e}^{-t}\right|,\quad\forall t\in[1,+\infty),\ \forall s\in[\alpha,\beta]
\]
故 $\displaystyle\int_{1}^{+\infty}t^{s-1}(\ln t)^{k}\mathrm{e}^{-t}\,\mathrm{d}t$ 在 $s\in[\alpha,\beta]$ 上一致收敛. 因而由定理~\ref{thm:15.3.6}知
\[
  \frac{\mathrm{d}^{l}}{\mathrm{d}s^{l}}\left(\int_{1}^{+\infty}t^{s-1}\mathrm{e}^{-t}\,\mathrm{d}t\right)=\int_{1}^{+\infty}\frac{\mathrm{d}^{l}}{\mathrm{d}s^{l}}t^{s-1}\mathrm{e}^{-t}\,\mathrm{d}t=\int_{1}^{+\infty}t^{s-1}(\ln t)^{l}\mathrm{e}^{-t}\,\mathrm{d}t,\quad l=1,2,\cdots.
\]
综上可知, 当 $\beta>\alpha>0$ 时,
\[
  \begin{aligned}
    \Gamma^{(l)}(s)&=\left(\int_{0}^{1}t^{s-1}\mathrm{e}^{-t}\,\mathrm{d}t\right)^{(l)}+\left(\int_{1}^{+\infty}t^{s-1}\mathrm{e}^{-t}\,\mathrm{d}t\right)^{(l)}\\
    &=\int_{0}^{1}t^{s-1}(\ln t)^{l}\mathrm{e}^{-t}\,\mathrm{d}t+\int_{1}^{+\infty}t^{s-1}(\ln t)^{l}\mathrm{e}^{-t}\,\mathrm{d}t\\
    &=\int_{0}^{+\infty}t^{s-1}(\ln t)^{l}\mathrm{e}^{-t}\,\mathrm{d}t.
  \end{aligned}
\]
由 $\beta>\alpha>0$ 的任意性知, $\displaystyle\Gamma(s)=\int_{0}^{+\infty}t^{s-1}\mathrm{e}^{-t}\,\mathrm{d}t$ 在 $(0,+\infty)$ 上有各阶导数, 因而有各阶连续导数.
\end{proof}

\begin{theorem}\label{thm:15.4.2}
$\Gamma$ 函数具有下面的三条性质:
\begin{enumerate}
  \item $\forall s>0$, $\Gamma(s)>0$, 且 $\Gamma(1)=1$;
  \item $\Gamma(s+1)=s\Gamma(s)$, $\forall s>0$;
  \item $\ln\Gamma(s)$ 为 $(0,+\infty)$ 上的凸函数.
\end{enumerate}
\end{theorem}

\begin{proof}
(1) 因为 $t^{s-1}\mathrm{e}^{-t}$ 在 $t\in(0,+\infty)$ 上连续且恒大于 0, 故 $\displaystyle\Gamma(s)=\int_{0}^{+\infty}t^{s-1}\mathrm{e}^{-t}\,\mathrm{d}t>0$.

又
\[
  \Gamma(1)=\int_{0}^{+\infty}t^{1-1}\mathrm{e}^{-t}\,\mathrm{d}t=\int_{0}^{+\infty}\mathrm{e}^{-t}\,\mathrm{d}t=-\left.\mathrm{e}^{-t}\right|_{0}^{+\infty}=1.
\]

(2)
\[
  \Gamma(s+1)=\int_{0}^{+\infty}t^{s}\mathrm{e}^{-t}\,\mathrm{d}t\xlongequal{\text{分部积分}}-\left.t^{s}\mathrm{e}^{-t}\right|_{0}^{+\infty}+s\int_{0}^{+\infty}t^{s-1}\mathrm{e}^{-t}\,\mathrm{d}t=s\Gamma(s).
\]

(3) 对 $p\in(1,+\infty)$, $\displaystyle\frac{1}{p}+\frac{1}{q}=1$, $s_{1},s_{2}\in(0,+\infty)$, 有
\[
  \begin{aligned}
    \Gamma\left(\frac{s_{1}}{p}+\frac{s_{2}}{q}\right)&=\int_{0}^{+\infty}t^{\frac{s_{1}}{p}+\frac{s_{2}}{q}-1}\mathrm{e}^{-t}\,\mathrm{d}t=\int_{0}^{+\infty}\left(t^{\frac{s_{1}-1}{p}}\mathrm{e}^{-\frac{t}{p}}\right)\left(t^{\frac{s_{2}-1}{q}}\mathrm{e}^{-\frac{t}{q}}\right)\mathrm{d}t\\
    &\stackrel{\mathrm{H\ddot{o}lder}}{\leqslant}\left(\int_{0}^{+\infty}t^{s_{1}-1}\mathrm{e}^{-t}\,\mathrm{d}t\right)^{\frac{1}{p}}\left(\int_{0}^{+\infty}t^{s_{2}-1}\mathrm{e}^{-t}\,\mathrm{d}t\right)^{\frac{1}{q}}=\Gamma^{\frac{1}{p}}(s_{1})\Gamma^{\frac{1}{q}}(s_{2}),
  \end{aligned}
\]
\[
  \ln\Gamma\left(\frac{s_{1}}{p}+\frac{s_{2}}{q}\right)\leqslant\ln\left[\Gamma^{\frac{1}{p}}(s_{1})\Gamma^{\frac{1}{q}}(s_{2})\right]=\frac{1}{p}\ln\Gamma(s_{1})+\frac{1}{q}\ln\Gamma(s_{2}),
\]
这就证明了 $\ln\Gamma(s)$ 为 $(0,+\infty)$ 上的凸函数.
\end{proof}

\begin{corollary}\label{cor:15.4.1}
\begin{enumerate}
  \item 设 $n\in\mathbb{N}$, 则 $\Gamma(n+1)=n!$ (故 $\Gamma$ 函数可视作阶乘函数的推广);
  \item 设 $n\in\mathbb{N}$, $n<s\leqslant n+1$, 即 $0<s-n\leqslant 1$, 则 $\Gamma(s+1)=s(s-1)\cdots(s-n)\Gamma(s-n)$, 这表明 $\Gamma$ 的值完全由 $\Gamma$ 在 $(0,1]$ 中的值确定.
\end{enumerate}
\end{corollary}

\begin{proof}
(1)
\[
  \begin{aligned}
    \Gamma(n+1)&\xlongequal{\text{定理~\ref{thm:15.4.2}(2)}}n\Gamma(n)\xlongequal{\text{定理~\ref{thm:15.4.2}(2)}}n(n-1)\Gamma(n-1)\\
    &=\cdots=n(n-1)\cdots 1\cdot\Gamma(1)=n!.
  \end{aligned}
\]

(2)
\[
  \Gamma(s+1)\xlongequal{\text{定理~\ref{thm:15.4.2}(2)}}s\Gamma(s)=\cdots=s(s-1)\cdots(s-n)\Gamma(s-n),\quad\text{其中 }0<s-n\leqslant 1.
\]
\end{proof}

出乎意料的是, Bohr 与 Mollerup 首先发现了: 在 $(0,+\infty)$ 上满足定理~\ref{thm:15.4.2} 中 3 条性质的函数一定是 $\Gamma$ 函数.

\begin{theorem}[Bohr-Mollerup]\label{thm:15.4.3}
如果 $(0,+\infty)$ 上的函数 $f(x)$ 满足:
\begin{enumerate}
  \item $\forall x>0$, $f(x)>0$ 且 $f(1)=1$;
  \item $f(x+1)=xf(x)$, $\forall x>0$;
  \item $\ln f(x)$ 为 $(0,+\infty)$ 上的凸函数.
\end{enumerate}
则 $f(x)=\Gamma(x)$, $\forall x\in(0,+\infty)$.
\end{theorem}

\begin{proof}
由条件 (1) 与条件 (2) 知, $f(n)=(n-1)f(n-2)=\cdots=(n-1)!f(1)=(n-1)!$.

设 $x\in(0,1)$, 由条件 (3), $\ln f(x)$ 为 $(0,+\infty)$ 上的凸函数, 根据\extref{定理 3.6.1}得到
\[
  \frac{\ln f(n)-\ln f(n-1)}{n-(n-1)}\leqslant\frac{\ln f(n+x)-\ln f(n)}{(n+x)-n}\leqslant\frac{\ln f(n+1)-\ln f(n)}{(n+1)-n},\quad n=2,3,\cdots.
\]
于是
\[
  \ln(n-1)!-\ln(n-2)!\leqslant\frac{\ln f(n+x)-\ln(n-1)!}{x}\leqslant\ln n!-\ln(n-1)!,
\]
即
\[
  x\ln(n-1)\leqslant\ln f(n+x)-\ln(n-1)!\leqslant x\ln n,
\]
\[
  x\ln(n-1)+\ln(n-1)!\leqslant\ln f(n+x)\leqslant x\ln n+\ln(n-1)!,
\]
\[
  \ln\left[(n-1)^{x}(n-1)!\right]\leqslant\ln f(n+x)\leqslant\ln\left[n^{x}(n-1)!\right],
\]
\[
  (n-1)^{x}(n-1)!\leqslant f(n+x)\xlongequal{\text{由条件 (2)}}(n-1+x)\cdots(1+x)xf(x)\leqslant n^{x}(n-1)!,
\]
\[
  \frac{(n-1)^{x}(n-1)!}{x(x+1)\cdots(x+n-1)}\leqslant f(x)\leqslant\frac{n^{x}(n-1)!}{x(x+1)\cdots(x+n-1)}.
\]
将左边不等式中的 $n-1$ 换成 $n$, 不等式仍然成立. 则有
\[
  \frac{n^{x}n!}{x(x+1)\cdots(x+n)}\leqslant f(x)\leqslant\frac{n^{x}n!}{x(x+1)\cdots(x+n)}\frac{x+n}{n},
\]
\[
  \frac{n}{x+n}f(x)\leqslant\frac{n^{x}n!}{x(x+1)\cdots(x+n)}\leqslant f(x).
\]
由于 $\displaystyle\lim_{n\to+\infty}\frac{n}{x+n}f(x)=f(x)=\lim_{n\to+\infty}f(x)$ 及夹逼定理得到
\[
  \lim_{n\to+\infty}\frac{n^{x}n!}{x(x+1)\cdots(x+n)}=f(x).
\]
这就证明了 $f(x)$ ($x\in(0,1)$) 被左边的极限所惟一确定. 再由条件 (2) 推得 $f(x)$ ($x\in(0,+\infty)$) 由条件 (1)、(2)、(3) 所惟一确定, 即 $f(x)=\Gamma(x)$, $x\in(0,+\infty)$.
\end{proof}

\begin{theorem}[$\Gamma$ 函数的另一表达式]\label{thm:15.4.4}
$\forall x>0$, 有
\[
  \Gamma(x)=\lim_{n\to+\infty}\frac{n^{x}n!}{x(x+1)\cdots(x+n)}.
\]
\end{theorem}

\begin{proof}
定理~\ref{thm:15.4.3} 的证明中, 已经证明当 $0<x\leqslant 1$ 时, 等式是成立的. 如果记
\[
  g(x)=\lim_{n\to+\infty}\frac{n^{x}n!}{x(x+1)\cdots(x+n)},
\]
由于
\[
  \begin{aligned}
    g(x+1)&=\lim_{n\to+\infty}\frac{n^{x+1}n!}{(x+1)(x+2)\cdots(x+n+1)}\\
    &=x\lim_{n\to+\infty}\frac{n^{x}n!}{x(x+1)\cdots(x+n)}\cdot\frac{n}{x+n+1}=xg(x),
  \end{aligned}
\]
由此即知 $g(x)$ 在 $(0,1]$ 中的值完全确定了 $g(x)$ 在 $(0,+\infty)$ 上的值. 于是, 从 $g(x)=\Gamma(x)$, $x\in(0,1]$, $\Gamma(x+1)=x\Gamma(x)$, $g(x+1)=xg(x)$ 立刻知, $g(x)=\Gamma(x)$, $x>0$.
\end{proof}

表面上似乎不相关联的 $\Gamma$ 函数与 $\mathrm{B}$ 函数实际上有着密切的联系.

\begin{theorem}[$\Gamma$ 函数与 $\mathrm{B}$ 函数的关联公式]\label{thm:15.4.5}
$\forall p>0, q>0$, 有
\[
  \mathrm{B}(p,q)=\frac{\Gamma(p)\Gamma(q)}{\Gamma(p+q)}.
\]
\end{theorem}

\begin{proof}
\textbf{证法 1}\quad 固定 $q>0$, 令
\[
  f(p)=\frac{\Gamma(p+q)\mathrm{B}(p,q)}{\Gamma(q)}.
\]
先证 $\mathrm{B}$ 函数有递推公式:
\[
  \mathrm{B}(p+1,q)=\frac{p}{p+q}\mathrm{B}(p,q).
\]
事实上
\[
  \begin{aligned}
    \mathrm{B}(p+1,q)&=\int_{0}^{1}t^{p}(1-t)^{q-1}\,\mathrm{d}t=\int_{0}^{1}\left(\frac{t}{1-t}\right)^{p}(1-t)^{p+q-1}\,\mathrm{d}t\\
    &=\int_{0}^{1}\left(\frac{t}{1-t}\right)^{p}\mathrm{d}\frac{-(1-t)^{p+q}}{p+q}=\left.\left(\frac{t}{1-t}\right)^{p}\cdot\frac{-(1-t)^{p+q}}{p+q}\right|_{0}^{1}\\
    &\quad+\int_{0}^{1}p\left(\frac{t}{1-t}\right)^{p-1}\cdot\frac{1}{(1-t)^{2}}\frac{(1-t)^{p+q}}{p+q}\,\mathrm{d}t\\
    &=\frac{p}{p+q}\int_{0}^{1}t^{p-1}(1-t)^{q-1}\,\mathrm{d}t=\frac{p}{p+q}\mathrm{B}(p,q).
  \end{aligned}
\]
或者
\[
  \begin{aligned}
    \mathrm{B}(p+1,q)&=\int_{0}^{1}t^{p}(1-t)^{q-1}\,\mathrm{d}t=\int_{0}^{1}t^{p}\mathrm{d}\frac{-(1-t)^{q}}{q}\\
    &=\left.t^{p}\frac{-(1-t)^{q}}{q}\right|_{0}^{1}+\int_{0}^{1}\frac{p}{q}t^{p-1}(1-t)^{q}\,\mathrm{d}t\\
    &=\frac{p}{q}\int_{0}^{1}t^{p-1}(1-t)^{q}\,\mathrm{d}t=\frac{p}{q}\int_{0}^{1}t^{p-1}\left[(1-t)^{q-1}-(1-t)^{q-1}t\right]\mathrm{d}t\\
    &=\frac{p}{q}\mathrm{B}(p,q)-\frac{p}{q}\mathrm{B}(p+1,q),
  \end{aligned}
\]
\[
  \mathrm{B}(p+1,q)=\frac{p}{p+q}\mathrm{B}(p,q).
\]
应用上式与 $\Gamma$ 函数的递推关系即得
\[
  \begin{aligned}
    f(p+1)&=\frac{1}{\Gamma(q)}\Gamma(p+1+q)\mathrm{B}(p+1,q)\\
    &=\frac{1}{\Gamma(q)}(p+q)\Gamma(p+q)\frac{p}{p+q}\mathrm{B}(p,q)\\
    &=pf(p).
  \end{aligned}
\]
所以 $f(x)$ 具有定理~\ref{thm:15.4.3} 中的性质 (2). 注意到
\[
  \mathrm{B}(1,q)=\int_{0}^{1}(1-t)^{q-1}\,\mathrm{d}t=-\left.\frac{(1-t)^{q}}{q}\right|_{0}^{1}=\frac{1}{q},
\]
可得
\[
  f(1)=\frac{\Gamma(1+q)\mathrm{B}(1,q)}{\Gamma(q)}=\frac{\Gamma(q+1)}{q\Gamma(q)}=\frac{\Gamma(q+1)}{\Gamma(q+1)}=1.
\]
此外显然有 $f(p)>0$, 可见 $f(x)$ 具有定理~\ref{thm:15.4.3} 中的性质 (1). 再证 $f(x)$ 具有定理~\ref{thm:15.4.3} 中的性质 (3). 由于 $q$ 固定, 故 $-\ln\Gamma(q)$ 为常值函数, 它为 $p$ 在 $(0,+\infty)$ 上的凸函数. 由定理~\ref{thm:15.4.2} 中的 (3) 知, $\ln\Gamma(p+q)$ 为 $p\in(0,+\infty)$ 上的凸函数. 对 $u\in(1,+\infty)$, $\displaystyle\frac{1}{u}+\frac{1}{v}=1$, $p_{1},p_{2}\in(0,+\infty)$, 有
\[
  \begin{aligned}
    \mathrm{B}\left(\frac{p_{1}}{u}+\frac{p_{2}}{v},q\right)&=\int_{0}^{1}t^{\frac{p_{1}}{u}+\frac{p_{2}}{v}-1}(1-t)^{q-1}\,\mathrm{d}t\\
    &=\int_{0}^{1}t^{\frac{p_{1}-1}{u}}(1-t)^{\frac{q-1}{u}}t^{\frac{p_{2}-1}{v}}(1-t)^{\frac{q-1}{v}}\mathrm{d}t\\
    &\stackrel{\mathrm{H\ddot{o}lder}}{\leqslant}\left(\int_{0}^{1}t^{p_{1}-1}(1-t)^{q-1}\,\mathrm{d}t\right)^{\frac{1}{u}}\left(\int_{0}^{1}t^{p_{2}-1}(1-t)^{q-1}\,\mathrm{d}t\right)^{\frac{1}{v}}\\
    &=\mathrm{B}^{\frac{1}{u}}(p_{1},q)\mathrm{B}^{\frac{1}{v}}(p_{2},q),
  \end{aligned}
\]
\[
  \ln\mathrm{B}\left(\frac{p_{1}}{u}+\frac{p_{2}}{v},q\right)\leqslant\ln\left[\mathrm{B}^{\frac{1}{u}}(p_{1},q)\mathrm{B}^{\frac{1}{v}}(p_{2},q)\right]=\frac{1}{u}\ln\mathrm{B}(p_{1},q)+\frac{1}{v}\ln\mathrm{B}(p_{2},q),
\]
这就证明了 $\ln\mathrm{B}(p,q)$ 为 $p\in(0,+\infty)$ 上的凸函数. 综上得到
\[
  \ln f(p)=\ln\frac{\Gamma(p+q)\mathrm{B}(p,q)}{\Gamma(q)}=\ln\Gamma(p+q)+\ln\mathrm{B}(p,q)-\ln\Gamma(q)
\]
为三个关于 $p\in(0,+\infty)$ 的凸函数之和, 应该仍为 $p\in(0,+\infty)$ 的凸函数. 由此推得 $f(p)$ 满足定理~\ref{thm:15.4.3} 中的条件 (1)、(2)、(3), 故 $f(p)=\Gamma(p)$, $\forall p\in(0,+\infty)$, 即
\[
  \frac{\Gamma(p+q)\mathrm{B}(p,q)}{\Gamma(q)}=f(p)=\Gamma(p),
\]
也就是
\[
  \mathrm{B}(p,q)=\frac{\Gamma(p)\Gamma(q)}{\Gamma(p+q)},\quad p>0,q>0.
\]

\textbf{证法 2}\quad 参阅\extref{例 10.6.3}.
\end{proof}

从 $\Gamma$ 函数的性质, 可直接推出 $\mathrm{B}$ 函数的一些性质.

\begin{theorem}\label{thm:15.4.6}
\begin{enumerate}
  \item $\mathrm{B}(p,q)$ 在 $(0,+\infty)\times(0,+\infty)$ 上连续且有各阶连续偏导数;
  \item $\mathrm{B}(p,q)=\mathrm{B}(q,p)$, $p>0, q>0$ (对称性);
  \item $\displaystyle\mathrm{B}(p+1,q)=\frac{p}{p+q}\mathrm{B}(p,q)$;
  \item $\displaystyle\mathrm{B}(p+1,q+1)=\frac{pq}{(p+q+1)(p+q)}\mathrm{B}(p,q)$, $p>0, q>0$.
\end{enumerate}
\end{theorem}

\begin{proof}
(1) 根据定理~\ref{thm:15.4.5}, 有
\[
  \mathrm{B}(p,q)=\frac{\Gamma(p)\Gamma(q)}{\Gamma(p+q)},
\]
再由定理~\ref{thm:15.4.1} 推得 $\mathrm{B}(p,q)$ 在 $(p,q)\in(0,+\infty)\times(0,+\infty)$ 上连续且有各阶连续偏导数.

或者类似定理~\ref{thm:15.4.1} 直接证明 $\mathrm{B}(p,q)$ 在 $(0,+\infty)\times(0,+\infty)$ 上连续且有各阶连续偏导数.

(2)
\[
  \mathrm{B}(p,q)=\frac{\Gamma(p)\Gamma(q)}{\Gamma(p+q)}=\frac{\Gamma(q)\Gamma(p)}{\Gamma(q+p)}=\mathrm{B}(q,p).
\]
或者
\[
  \mathrm{B}(p,q)=\int_{0}^{1}t^{p-1}(1-t)^{q-1}\,\mathrm{d}t\xlongequal{t=1-s}\int_{1}^{0}(1-s)^{p-1}s^{q-1}(-\mathrm{d}s)=\int_{0}^{1}s^{q-1}(1-s)^{p-1}\,\mathrm{d}s=\mathrm{B}(q,p).
\]

(3) 在定理~\ref{thm:15.4.5} 证法 1 中已证明.

(4)
\[
  \begin{aligned}
    \mathrm{B}(p+1,q+1)&\xlongequal{(3)}\frac{p}{p+(q+1)}\mathrm{B}(p,q+1)\xlongequal{(2),(3)}\frac{p}{p+q+1}\cdot\frac{q}{q+p}\mathrm{B}(p,q)\\
    &=\frac{pq}{(p+q+1)(p+q)}\mathrm{B}(p,q).
  \end{aligned}
\]
\end{proof}

$\Gamma$ 函数还有加倍公式与余元公式这两个重要公式.

\begin{theorem}[加倍公式]\label{thm:15.4.7}
$\forall s>0$, 有
\[
  \Gamma(2s)=\frac{2^{2s-1}}{\sqrt{\pi}}\Gamma(s)\Gamma\left(s+\frac{1}{2}\right).
\]
\end{theorem}

\begin{proof}
显然
\[
  \Gamma(2s)=\frac{2^{2s-1}}{\sqrt{\pi}}\Gamma(s)\Gamma\left(s+\frac{1}{2}\right) \iff \Gamma(p)=\frac{2^{p-1}}{\sqrt{\pi}}\Gamma\left(\frac{p}{2}\right)\Gamma\left(\frac{p+1}{2}\right),\quad p=2s.
\]
为证后一式, 令
\[
  f(p)=\frac{2^{p-1}}{\sqrt{\pi}}\Gamma\left(\frac{p}{2}\right)\Gamma\left(\frac{p+1}{2}\right),
\]
则
\[
  f(1)=\frac{1}{\sqrt{\pi}}\Gamma\left(\frac{1}{2}\right)\Gamma(1)=\frac{1}{\sqrt{\pi}}\sqrt{\pi}\cdot 1=1,
\]
\[
  \begin{aligned}
    f(p+1)&=\frac{2^{p}}{\sqrt{\pi}}\Gamma\left(\frac{p+1}{2}\right)\Gamma\left(\frac{p+2}{2}\right)=\frac{2^{p}}{\sqrt{\pi}}\frac{p}{2}\Gamma\left(\frac{p}{2}\right)\Gamma\left(\frac{p+1}{2}\right)\\
    &=pf(p).
  \end{aligned}
\]
又因为
\[
  \ln f(p)=\ln\frac{2^{p-1}}{\sqrt{\pi}}+\ln\Gamma\left(\frac{p}{2}\right)+\ln\Gamma\left(\frac{p+1}{2}\right)
\]
为 $p\in(0,+\infty)$ 上的三个凸函数之和, 所以 $\ln f(p)$ 也为 $p\in(0,+\infty)$ 上的凸函数. 根据定理~\ref{thm:15.4.3} 可知
\[
  \Gamma(p)=f(p)=\frac{2^{p-1}}{\sqrt{\pi}}\Gamma\left(\frac{p}{2}\right)\Gamma\left(\frac{p+1}{2}\right).
\]
\end{proof}

\begin{theorem}[余元公式]\label{thm:15.4.8}
$\forall p\in(0,1)$, 有
\[
  \Gamma(p)\Gamma(1-p)=\frac{\pi}{\sin p\pi}.
\]
特别当 $p=\frac{1}{2}$ 时, $\Gamma^{2}\left(\frac{1}{2}\right)=\pi$, $\Gamma\left(\frac{1}{2}\right)=\sqrt{\pi}$.
\[
  \begin{aligned}
    \int_{0}^{+\infty}\mathrm{e}^{-x^{2}}\,\mathrm{d}x&\xlongequal{t=x^{2}}\int_{0}^{+\infty}\mathrm{e}^{-t}\frac{\mathrm{d}t}{2\sqrt{t}}\\
    &=\frac{1}{2}\int_{0}^{+\infty}t^{\frac{1}{2}-1}\mathrm{e}^{-t}\,\mathrm{d}t=\frac{1}{2}\Gamma\left(\frac{1}{2}\right)=\frac{\sqrt{\pi}}{2}.
  \end{aligned}
\]
根据余元公式, 从 $\Gamma$ 函数在 $\left(0,\frac{1}{2}\right)$ 中的值, 便能算出 $\Gamma$ 函数在 $(0,1)$ 中的值, 从而算出 $\Gamma$ 函数在 $(0,+\infty)$ 中的值.
\end{theorem}

\begin{proof}
\[
  \begin{aligned}
    \Gamma(p)\Gamma(1-p)&=\mathrm{B}(1-p,p)\Gamma(1-p+p)\\
    &=\mathrm{B}(1-p,p)\Gamma(1)=\int_{0}^{1}t^{-p}(1-t)^{p-1}\,\mathrm{d}t\\
    &\xlongequal[\displaystyle 1-t=\frac{x}{1+x}]{\displaystyle t=\frac{1}{1+x}}\int_{+\infty}^{0}\left(\frac{1}{1+x}\right)^{-p}\left(\frac{x}{1+x}\right)^{p-1}\frac{-\mathrm{d}x}{(1+x)^{2}}\\
    &=\int_{0}^{+\infty}\frac{x^{p-1}}{1+x}\,\mathrm{d}x\xlongequal{\text{\extref{例 16.1.5}}}\frac{\pi}{\sin p\pi}.
  \end{aligned}
\]
\end{proof}

为了处理与 $n!$ 有关的问题, 我们来给出 $n!$ 的一个渐近表达式, 进而证明精细 Stirling 公式.

\begin{lemma}\label{lem:15.4.1}
\[
  \int_{0}^{n}\frac{[t]-t+\frac{1}{2}}{t+x}\,\mathrm{d}t=\left(n+x+\frac{1}{2}\right)\ln(n+x)-\left(\frac{1}{2}+x\right)\ln x-\sum_{k=1}^{n}\ln(k+x)-n.
\]
\end{lemma}

\begin{proof}
\[
  \begin{aligned}
    \int_{0}^{n}\frac{[t]-t+\frac{1}{2}}{t+x}\,\mathrm{d}t&=\sum_{k=0}^{n-1}\int_{k}^{k+1}\frac{k-t+\frac{1}{2}}{t-x}\,\mathrm{d}t=\sum_{k=0}^{n-1}\int_{k}^{k+1}\left(\frac{k+\frac{1}{2}+x}{t+x}-1\right)\mathrm{d}t\\ % 原书第一项分母印为 t-x
    &=\sum_{k=0}^{n-1}\left[\left(k+\frac{1}{2}+x\right)(\ln(k+1+x)-\ln(k+x))-1\right]\\
    &=\sum_{k=1}^{n}\left(k+x-\frac{1}{2}\right)\ln(k+x)-\sum_{k=0}^{n-1}\left(k+\frac{1}{2}+x\right)\ln(k+x)-n\\
    &=\sum_{k=1}^{n}\left(k+x+\frac{1}{2}\right)\ln(k+x)-\sum_{k=1}^{n}\ln(k+x)-\sum_{k=0}^{n-1}\left(k+x+\frac{1}{2}\right)\ln(k+x)-n\\
    &=\left(n+x+\frac{1}{2}\right)\ln(n+x)-\left(\frac{1}{2}+x\right)\ln x-\sum_{k=1}^{n}\ln(k+x)-n.
  \end{aligned}
\]
\end{proof}

\begin{lemma}\label{lem:15.4.2}
设 $u>0$, 则
\[
  0<\left(u+\frac{1}{2}\right)\ln\left(1+\frac{1}{u}\right)-1<\frac{1}{12}\left(\frac{1}{u}-\frac{1}{u+1}\right).
\]
\end{lemma}

\begin{proof}
\[
  \begin{aligned}
    0&<\left(u+\frac{1}{2}\right)\ln\left(1+\frac{1}{u}\right)-1=\frac{2u+1}{2}\ln\frac{1+\frac{1}{2u+1}}{1-\frac{1}{2u+1}}-1\\
    &=\frac{2u+1}{2}\sum_{n=1}^{\infty}\frac{(-1)^{n-1}}{n}\left[\left(\frac{1}{2u+1}\right)^{n}-\left(-\frac{1}{2u+1}\right)^{n}\right]-1\\
    &=\sum_{n=1}^{\infty}\frac{(-1)^{n-1}+1}{2n}\left(\frac{1}{2u+1}\right)^{n-1}-1=\sum_{k=1}^{\infty}\frac{1}{2k+1}\frac{1}{(2u+1)^{2k}}\\
    &=\frac{1}{3}\sum_{k=1}^{\infty}\frac{3}{2k+1}\frac{1}{(2u+1)^{2k}}<\frac{1}{3}\sum_{k=1}^{\infty}\frac{1}{(2u+1)^{2k}}\\
    &=\frac{1}{3(2u+1)^{2}}\frac{1}{1-\frac{1}{(2u+1)^{2}}}=\frac{1}{12u(u+1)}=\frac{1}{12}\left(\frac{1}{u}-\frac{1}{u+1}\right).
  \end{aligned}
\]
\end{proof}

\begin{lemma}[Stirling 渐近表达式]\label{lem:15.4.3}
\[
  n!\sim\sqrt{2\pi n}\left(\frac{n}{\mathrm{e}}\right)^{n}\quad(n\to+\infty).
\]
\end{lemma}

\begin{proof}
应用\extref{例 6.4.5}(2)中 Wallis 公式
\[
  \lim_{n\to+\infty}\left[\frac{(2n)!!}{(2n-1)!!}\right]\frac{1}{2n+1}=\frac{\pi}{2} % 原书此处漏印平方，照录
\]
得到
\[
  \begin{aligned}
    \lim_{n\to+\infty}\frac{(n!)^{2}2^{2n}}{(2n)!\sqrt{n}}&=\lim_{n\to+\infty}\frac{[(2n)!!]^{2}}{(2n-1)!!(2n)!!\sqrt{n}}=\lim_{n\to+\infty}\frac{(2n)!!}{(2n-1)!!\sqrt{n}}\\
    &=\lim_{n\to+\infty}\sqrt{\left[\frac{(2n)!!}{(2n-1)!!}\right]^{2}\frac{1}{2n+1}}\sqrt{\frac{2n+1}{n}}\\
    &=\sqrt{\frac{\pi}{2}}\cdot\sqrt{2}=\sqrt{\pi}.
  \end{aligned}
\]
考虑数列 $\{a_{n}\}$:
\[
  a_{n}=\frac{n!\mathrm{e}^{n}}{n^{n+\frac{1}{2}}},\quad n=1,2,\cdots,
\]
则
\[
  \frac{a_{n}}{a_{n+1}}=\frac{\frac{n!\mathrm{e}^{n}}{n^{n+\frac{1}{2}}}}{\frac{(n+1)!\mathrm{e}^{n+1}}{(n+1)^{n+1+\frac{1}{2}}}}=\frac{1}{\mathrm{e}}\left(1+\frac{1}{n}\right)^{n+\frac{1}{2}}.
\]
取对数, 得
\[
  \ln\frac{a_{n}}{a_{n+1}}=\left(n+\frac{1}{2}\right)\ln\left(1+\frac{1}{n}\right)-1.
\]
根据引理~\ref{lem:15.4.2}, 有
\[
  \begin{gathered}
    0<\ln\frac{a_{n}}{a_{n+1}}=\left(n+\frac{1}{2}\right)\ln\left(1+\frac{1}{n}\right)-1<\frac{1}{12}\left(\frac{1}{n}-\frac{1}{n+1}\right),\\
    1<\frac{a_{n}}{a_{n+1}}<\mathrm{e}^{\frac{1}{12}\left(\frac{1}{n}-\frac{1}{n+1}\right)}.
  \end{gathered}
\]
现在用 $n+j$ ($j=0,1,\cdots,k-1$) 代上述不等式中的 $n$ 得出
\[
  \begin{aligned}
    1&<\frac{a_{n}}{a_{n+1}}<\mathrm{e}^{\frac{1}{12}\left(\frac{1}{n}-\frac{1}{n+1}\right)},\\
    1&<\frac{a_{n+1}}{a_{n+2}}<\mathrm{e}^{\frac{1}{12}\left(\frac{1}{n+1}-\frac{1}{n+2}\right)},\\
    &\ \vdots\\
    1&<\frac{a_{n+k-1}}{a_{n+k}}<\mathrm{e}^{\frac{1}{12}\left(\frac{1}{n+k-1}-\frac{1}{n+k}\right)},
  \end{aligned}
\]
将上述 $k$ 个不等式乘起来便有
\[
  1<\frac{a_{n}}{a_{n+k}}<\mathrm{e}^{\frac{1}{12}\left(\frac{1}{n}-\frac{1}{n+k}\right)}.
\]
$\displaystyle\frac{a_{n}}{a_{n+1}}>1$ 表明数列 $\{a_{n}\}$ 严格单调减少且有下界 0, 所以应当收敛于某个数 $\alpha$. 于是, 令 $k\to+\infty$, 得到
\[
  1<\frac{a_{n}}{\alpha}\leqslant\mathrm{e}^{\frac{1}{12n}}.
\]
由此可知 $\alpha>0$. 进而有
\[
  \begin{aligned}
    \lim_{n\to+\infty}\frac{n!\mathrm{e}^{n}}{n^{n+\frac{1}{2}}}&=\lim_{n\to+\infty}a_{n}=\alpha=\lim_{n\to+\infty}\frac{a_{n}^{2}}{a_{2n}}=\lim_{n\to+\infty}\frac{\left(\frac{n!\mathrm{e}^{n}}{n^{n+\frac{1}{2}}}\right)^{2}}{\frac{(2n)!\mathrm{e}^{2n}}{(2n)^{2n+\frac{1}{2}}}}\\
    &=\lim_{n\to+\infty}\frac{(n!)^{2}\cdot 2^{2n+\frac{1}{2}}}{(2n)!\sqrt{n}}=\sqrt{2\pi},
  \end{aligned}
\]
即
\[
  \lim_{n\to+\infty}\frac{n!}{\sqrt{2n\pi}\left(\frac{n}{\mathrm{e}}\right)^{n}}=1,
\]
\[
  n!\sim\sqrt{2n\pi}\left(\frac{n}{\mathrm{e}}\right)^{n},\quad n\to+\infty.
\]
\end{proof}

\begin{theorem}[Stirling 公式]\label{thm:15.4.9}
$\forall x>0$, $\exists\theta(x)\in(0,1)$, 使得
\[
  \Gamma(x+1)=\sqrt{2\pi x}\left(\frac{x}{\mathrm{e}}\right)^{x}\mathrm{e}^{\frac{\theta(x)}{12x}}.
\]
\end{theorem}

\begin{proof}
\[
  \begin{aligned}
    \ln\Gamma(x)&\xlongequal{\text{定理~\ref{thm:15.4.4}}}\ln\lim_{n\to+\infty}\frac{n^{x}n!}{x(x+1)\cdots(x+n)}=\lim_{n\to+\infty}\ln\frac{n^{x}n!}{x(x+1)\cdots(x+n)}\\
    &=\lim_{n\to+\infty}\left[\ln n!+x\ln n-\sum_{k=0}^{n}\ln(x+k)\right]\\
    &=\lim_{n\to+\infty}\left[\ln n!+n-\left(n+\frac{1}{2}\right)\ln n+\left(x-\frac{1}{2}\right)\ln x-\left(n+x+\frac{1}{2}\right)\ln\left(1+\frac{x}{n}\right)\right.\\
    &\quad\left.+\left(n+x+\frac{1}{2}\right)\ln(n+x)-\left(\frac{1}{2}+x\right)\ln x-\sum_{k=1}^{n}\ln(k+x)-n\right]\\
    &=\lim_{n\to+\infty}\left\{\left[\ln n!+n-\left(n+\frac{1}{2}\right)\ln n\right]+\left(x-\frac{1}{2}\right)\ln x-\left(n+x+\frac{1}{2}\right)\ln\left(1+\frac{x}{n}\right)\right.\\
    &\quad\left.+\int_{0}^{n}\frac{[t]-t+\frac{1}{2}}{t+x}\,\mathrm{d}t\right\}\\
    &=\lim_{n\to+\infty}\left\{\left[\ln n!+n-\left(n+\frac{1}{2}\right)\ln n\right]+\left(x-\frac{1}{2}\right)\ln x-\ln\left[\left(1+\frac{x}{n}\right)^{\frac{n}{x}}\right]^{\frac{n+x+\frac{1}{2}}{n}x}\right.\\
    &\quad\left.+\int_{0}^{n}\frac{[t]-t+\frac{1}{2}}{t+x}\,\mathrm{d}t\right\}\\
    &\xlongequal{\text{引理~\ref{lem:15.4.3}}}\ln\sqrt{2\pi}+\left(x-\frac{1}{2}\right)\ln x-x+\int_{0}^{+\infty}\frac{[t]-t+\frac{1}{2}}{t+x}\,\mathrm{d}t\\
    &=\ln\sqrt{2\pi}+\left(x-\frac{1}{2}\right)\ln x-x+\frac{\theta(x)}{12x},
  \end{aligned}
\]
即
\[
  \Gamma(x)=\sqrt{2\pi}x^{x-\frac{1}{2}}\mathrm{e}^{-x}\mathrm{e}^{\frac{\theta(x)}{12x}},
\]
于是
\[
  \Gamma(x+1)=x\Gamma(x)=\sqrt{2\pi x}\left(\frac{x}{\mathrm{e}}\right)^{x}\mathrm{e}^{\frac{\theta(x)}{12x}}.
\]
剩下的要证
\[
  \theta(x)=12x\int_{0}^{+\infty}\frac{[t]-t+\frac{1}{2}}{t+x}\,\mathrm{d}t\in(0,1).
\]
事实上, 由引理~\ref{lem:15.4.3}, 并取 $u=k+x$ 得到 % 原书此处笔误印为引理 15.4.3，实为引理 15.4.2，照录
\[
  \begin{aligned}
    0<\theta(x)&=12x\int_{0}^{+\infty}\frac{[t]-t+\frac{1}{2}}{t+x}\,\mathrm{d}t=12x\sum_{k=0}^{\infty}\int_{k}^{k+1}\frac{k-t+\frac{1}{2}}{t+x}\,\mathrm{d}t\\
    &=12x\sum_{k=0}^{\infty}\int_{k}^{k+1}\left(\frac{k+\frac{1}{2}+x}{t+x}-1\right)\mathrm{d}t\\
    &=12x\sum_{k=0}^{\infty}\left[\left(k+\frac{1}{2}+x\right)\ln\left(1+\frac{1}{k+x}\right)-1\right]\\
    &<12x\cdot\frac{1}{12}\sum_{k=0}^{\infty}\left(\frac{1}{k+x}-\frac{1}{k+x+1}\right)=x\cdot\frac{1}{x}=1.
  \end{aligned}
\]
\end{proof}

% ===== PART-B =====

\begin{example}[$\mathrm{B}(p,q)$ 的其他形式]\label{ex:15.4.1}
设 $p>0, q>0$, 则有
\begin{enumerate}
  \item $\displaystyle\mathrm{B}(p,q)=2\int_{0}^{\frac{\pi}{2}}\sin^{2q-1}\varphi\cos^{2p-1}\varphi\,\mathrm{d}\varphi$;
  \item $\displaystyle\mathrm{B}(p,q)=\int_{0}^{+\infty}\frac{y^{p-1}}{(1+y)^{p+q}}\,\mathrm{d}y$;
  \item $\displaystyle\mathrm{B}(p,q)=\int_{0}^{1}\frac{y^{p-1}+y^{q-1}}{(1+y)^{p+q}}\,\mathrm{d}y$.
\end{enumerate}
\end{example}
\begin{proof}
(1) $\displaystyle\mathrm{B}(p,q)=\int_{0}^{1}t^{p-1}(1-t)^{q-1}\,\mathrm{d}t$
\[
  \xlongequal{t=\cos^{2}\varphi}\int_{\frac{\pi}{2}}^{0}\cos^{2p-2}\varphi\sin^{2q-2}\varphi\cdot 2\cos\varphi(-\sin\varphi)\,\mathrm{d}\varphi
  =2\int_{0}^{\frac{\pi}{2}}\sin^{2q-1}\varphi\cos^{2p-1}\varphi\,\mathrm{d}\varphi.
\]

(2) 令 $t=\frac{y}{1+y}$, $1-t=\frac{1}{1+y}$, $\mathrm{d}t=\frac{\mathrm{d}y}{(1+y)^{2}}$, 则有
\[
  \begin{aligned}
    \mathrm{B}(p,q)&=\int_{0}^{1}t^{p-1}(1-t)^{q-1}\,\mathrm{d}t=\int_{0}^{+\infty}\left(\frac{y}{1+y}\right)^{p-1}\left(\frac{1}{1+y}\right)^{q-1}\frac{\mathrm{d}y}{(1+y)^{2}}\\
    &=\int_{0}^{+\infty}\frac{y^{p-1}}{(1+y)^{p+q}}\,\mathrm{d}y.
  \end{aligned}
\]

(3) $\displaystyle\mathrm{B}(p,q)\xlongequal{(2)}\int_{0}^{+\infty}\frac{y^{p-1}}{(1+y)^{p+q}}\,\mathrm{d}y=\int_{0}^{1}\frac{y^{p-1}}{(1+y)^{p+q}}\,\mathrm{d}y+\int_{1}^{+\infty}\frac{y^{p-1}}{(1+y)^{p+q}}\,\mathrm{d}y$
\[
  \begin{aligned}
    &\xlongequal{y=\frac{1}{t}}\int_{0}^{1}\frac{y^{p-1}}{(1+y)^{p+q}}\,\mathrm{d}y+\int_{1}^{0}\frac{\left(\frac{1}{t}\right)^{p-1}}{\left(1+\frac{1}{t}\right)^{p+q}}\cdot\left(-\frac{\mathrm{d}t}{t^{2}}\right)\\
    &=\int_{0}^{1}\frac{y^{p-1}}{(1+y)^{p+q}}\,\mathrm{d}y+\int_{0}^{1}\frac{t^{q-1}}{(1+t)^{p+q}}\,\mathrm{d}t=\int_{0}^{1}\frac{y^{p-1}+y^{q-1}}{(1+y)^{p+q}}\,\mathrm{d}y.
  \end{aligned}
\]
\end{proof}

\begin{example}\label{ex:15.4.2}
证明：$\displaystyle\mathrm{B}(m,n)=\frac{(m-1)!(n-1)!}{(m+n-1)!}$.
\end{example}
\begin{proof}[证法 1]
\[
  \mathrm{B}(m,n)=\frac{\Gamma(m)\Gamma(n)}{\Gamma(m+n)}=\frac{(m-1)!(n-1)!}{(m+n-1)!}.
\]
\end{proof}
\begin{proof}[证法 2]
\[
  \begin{aligned}
    \mathrm{B}(m,n)&\xlongequal{\text{定理~\ref{thm:15.4.6}(3)}}\frac{m-1}{m+n-1}\cdot\frac{m-2}{m+n-2}\cdots\frac{1}{n+1}\mathrm{B}(1,n)\\
    &=\frac{(m-1)!n!}{(m+n-1)!}\int_{0}^{1}(1-t)^{n-1}\,\mathrm{d}t=\left.\frac{(m-1)!n!}{(m+n-1)!}\frac{(1-t)^{n}}{-n}\right|_{0}^{1}\\
    &=\frac{(m-1)!(n-1)!}{(m+n-1)!}.
  \end{aligned}
\]
\end{proof}

\begin{example}\label{ex:15.4.3}
证明：
\begin{enumerate}
  \item $\displaystyle\int_{0}^{\frac{\pi}{2}}\cos^{\alpha}x\sin^{\beta}x\,\mathrm{d}x=\frac{1}{2}\mathrm{B}\left(\frac{\beta+1}{2},\frac{\alpha+1}{2}\right)=\frac{1}{2}\frac{\Gamma\left(\frac{\beta+1}{2}\right)\Gamma\left(\frac{\alpha+1}{2}\right)}{\Gamma\left(\frac{\alpha+\beta}{2}+1\right)}$;
  \item $\displaystyle\int_{0}^{\frac{\pi}{2}}\cos^{6}x\sin^{4}x\,\mathrm{d}x=\frac{3\pi}{512}$.
\end{enumerate}
\end{example}
\begin{proof}
(1)
\[
  \begin{aligned}
    \int_{0}^{\frac{\pi}{2}}\cos^{\alpha}x\sin^{\beta}x\,\mathrm{d}x&\xlongequal{\begin{subarray}{c}t=\sin^{2}x\\\mathrm{d}t=2\sin x\cos x\,\mathrm{d}x\end{subarray}}\frac{1}{2}\int_{0}^{1}t^{\frac{\beta-1}{2}}(1-t)^{\frac{\alpha-1}{2}}\,\mathrm{d}t\\
    &=\frac{1}{2}\mathrm{B}\left(\frac{\beta+1}{2},\frac{\alpha+1}{2}\right)=\frac{1}{2}\frac{\Gamma\left(\frac{\beta+1}{2}\right)\Gamma\left(\frac{\alpha+1}{2}\right)}{\Gamma\left(\frac{\alpha+\beta}{2}+1\right)}.
  \end{aligned}
\]

(2) 因为
\[
  \begin{aligned}
    \Gamma\left(\frac{5}{2}\right)&=\frac{3}{2}\Gamma\left(\frac{3}{2}\right)=\frac{3}{2}\cdot\frac{1}{2}\Gamma\left(\frac{1}{2}\right)=\frac{3}{4}\sqrt{\pi},\\
    \Gamma\left(\frac{7}{2}\right)&=\frac{5}{2}\Gamma\left(\frac{5}{2}\right)=\frac{5}{2}\cdot\frac{3}{4}\sqrt{\pi}=\frac{15}{8}\sqrt{\pi},
  \end{aligned}
\]
所以
\[
  \int_{0}^{\frac{\pi}{2}}\cos^{6}x\sin^{4}x\,\mathrm{d}x=\frac{1}{2}\frac{\Gamma\left(\frac{7}{2}\right)\Gamma\left(\frac{5}{2}\right)}{\Gamma\left(\frac{10}{2}+1\right)}=\frac{1}{2}\frac{\frac{15}{8}\sqrt{\pi}\cdot\frac{3}{4}\sqrt{\pi}}{5!}=\frac{3\pi}{512}.
\]
\end{proof}

\begin{example}\label{ex:15.4.4}
证明：瑕积分
\[
  \int_{0}^{\frac{\pi}{2}}(\tan x)^{\alpha}\,\mathrm{d}x=\frac{\pi}{2\cos\frac{\alpha\pi}{2}},\quad |\alpha|<1.
\]
\end{example}
\begin{proof}
\[
  \begin{aligned}
    \int_{0}^{\frac{\pi}{2}}(\tan x)^{\alpha}\,\mathrm{d}x&=\int_{0}^{\frac{\pi}{2}}\sin^{\alpha}x\cos^{-\alpha}x\,\mathrm{d}x\xlongequal{\text{例~\ref{ex:15.4.3}(1)}}\frac{1}{2}\frac{\Gamma\left(\frac{\alpha+1}{2}\right)\Gamma\left(\frac{-\alpha+1}{2}\right)}{\Gamma\left(\frac{\alpha+(-\alpha)}{2}+1\right)}\\
    &=\frac{1}{2}\Gamma\left(\frac{1+\alpha}{2}\right)\Gamma\left(\frac{1-\alpha}{2}\right)\xlongequal{\text{余元公式}}\frac{1}{2}\frac{\pi}{\sin\frac{1+\alpha}{2}\pi}=\frac{\pi}{2\cos\frac{\alpha\pi}{2}}.
  \end{aligned}
\]
\end{proof}

\begin{example}\label{ex:15.4.5}
证明：
\begin{enumerate}
  \item $\forall\alpha\in\mathbb{R}$, 有 $\displaystyle\lim_{x\to+\infty}\frac{x^{\alpha}\Gamma(x)}{\Gamma(x+\alpha)}=1$;
  \item $\displaystyle\lim_{n\to+\infty}\sqrt{n}\int_{0}^{1}(1-x^{2})^{n}\,\mathrm{d}x=\frac{\sqrt{\pi}}{2}$.
\end{enumerate}
\end{example}
\begin{proof}
(1) 由 Stirling 公式,
\[
  \begin{aligned}
    \lim_{x\to+\infty}\frac{x^{\alpha}\Gamma(x)}{\Gamma(x+\alpha)}&=\lim_{x\to+\infty}\frac{x^{\alpha}\sqrt{2\pi(x-1)}\left(\frac{x-1}{\mathrm{e}}\right)^{x-1}\mathrm{e}^{\frac{\theta(x-1)}{12(x-1)}}}{\sqrt{2\pi(x+\alpha-1)}\left(\frac{x+\alpha-1}{\mathrm{e}}\right)^{x+\alpha-1}\mathrm{e}^{\frac{\theta(x+\alpha-1)}{12(x+\alpha-1)}}}\\
    &=\lim_{x\to+\infty}\sqrt{\frac{x-1}{x+\alpha-1}}\frac{1}{\left[\left(1+\frac{\alpha}{x-1}\right)^{\frac{x-1}{\alpha}}\right]^{\alpha}}\left(\frac{x}{x+\alpha-1}\right)^{\alpha}\cdot\mathrm{e}^{\alpha}\cdot\mathrm{e}^{\frac{\theta(x-1)}{12(x-1)}-\frac{\theta(x+\alpha-1)}{12(x+\alpha-1)}}\\
    &=1\cdot\frac{1}{\mathrm{e}^{\alpha}}\cdot 1\cdot\mathrm{e}^{\alpha}\cdot 1=1.
  \end{aligned}
\]

(2) $\displaystyle\int_{0}^{1}(1-x^{2})^{n}\,\mathrm{d}x\xlongequal{\begin{subarray}{c}t=x^{2}\\\mathrm{d}t=2x\,\mathrm{d}x\end{subarray}}\int_{0}^{1}(1-t)^{n}\cdot\frac{1}{2}t^{-\frac{1}{2}}\,\mathrm{d}t=\frac{1}{2}\mathrm{B}\left(\frac{1}{2},n+1\right)$
\[
  =\frac{1}{2}\frac{\Gamma\left(\frac{1}{2}\right)\Gamma(n+1)}{\Gamma\left(n+\frac{3}{2}\right)}=\frac{\sqrt{\pi}}{2}\frac{\Gamma(n+1)}{\Gamma\left(n+\frac{3}{2}\right)},
\]
\[
  \lim_{n\to+\infty}\sqrt{n}\int_{0}^{1}(1-x^{2})^{n}\,\mathrm{d}x=\frac{\sqrt{\pi}}{2}\lim_{n\to+\infty}\frac{\sqrt{n}\Gamma(n+1)}{\Gamma\left(n+\frac{3}{2}\right)}=\frac{\sqrt{\pi}}{2}\lim_{n\to+\infty}\frac{n^{\frac{3}{2}}\Gamma(n)}{\Gamma\left(n+\frac{3}{2}\right)}\xlongequal{(1)}\frac{\sqrt{\pi}}{2}\cdot 1=\frac{\sqrt{\pi}}{2}.
\]
或者
\[
  \begin{aligned}
    \lim_{n\to+\infty}\sqrt{n}\int_{0}^{1}(1-x^{2})^{n}\,\mathrm{d}x&=\frac{\sqrt{\pi}}{2}\lim_{n\to+\infty}\frac{\sqrt{n}\Gamma(n+1)}{\Gamma\left(n+\frac{3}{2}\right)}=\lim_{n\to+\infty}\frac{\sqrt{\pi}}{2}\cdot\frac{\sqrt{n}\cdot n!}{\frac{2n+1}{2}\cdots\frac{1}{2}\Gamma\left(\frac{1}{2}\right)}\\
    &=\lim_{n\to+\infty}\frac{\sqrt{n}\cdot(2n)!!}{(2n+1)!!}\\
    &\xlongequal{\text{\extref{例 6.4.5}(2)}}\lim_{n\to+\infty}\sqrt{\frac{n}{2n+1}}\sqrt{\left[\frac{(2n)!!}{(2n-1)!!}\right]^{2}\frac{1}{2n+1}}\\
    &=\frac{1}{\sqrt{2}}\sqrt{\frac{\pi}{2}}=\frac{\sqrt{\pi}}{2}.
  \end{aligned}
\]
\end{proof}

\begin{example}\label{ex:15.4.6}
证明：
\begin{enumerate}
  \item $\displaystyle I=\int_{0}^{1}\ln\Gamma(x)\,\mathrm{d}x=\ln\sqrt{2\pi}$;
  \item $\displaystyle J=\int_{0}^{1}\sin\pi x\ln\Gamma(x)\,\mathrm{d}x=\frac{1}{\pi}\left(1+\ln\frac{\pi}{2}\right)$.
\end{enumerate}
\end{example}
\begin{proof}
(1) 显然
\[
  I=\int_{0}^{1}\ln\Gamma(x)\,\mathrm{d}x\xlongequal{x=1-t}\int_{1}^{0}\ln\Gamma(1-t)(-\mathrm{d}t)=\int_{0}^{1}\ln\Gamma(1-t)\,\mathrm{d}t,
\]
则
\[
  \begin{aligned}
    2I&=\int_{0}^{1}\ln\Gamma(x)\,\mathrm{d}x+\int_{0}^{1}\ln\Gamma(1-x)\,\mathrm{d}x=\int_{0}^{1}\ln[\Gamma(x)\Gamma(1-x)]\,\mathrm{d}x\\
    &\xlongequal{\text{余元公式}}\int_{0}^{1}\ln\frac{\pi}{\sin\pi x}\,\mathrm{d}x=\ln\pi-\int_{0}^{1}\ln\sin\pi x\,\mathrm{d}x\xlongequal{\theta=\pi x}\ln\pi-\frac{1}{\pi}\int_{0}^{\pi}\ln\sin\theta\,\mathrm{d}\theta\\
    &=\ln\pi-\frac{2}{\pi}\int_{0}^{\frac{\pi}{2}}\ln\sin\theta\,\mathrm{d}\theta\xlongequal{\text{\extref{例 6.5.18}}}\ln\pi-\frac{2}{\pi}\cdot\left(-\frac{\pi}{2}\ln 2\right)=\ln(2\pi),
  \end{aligned}
\]
\[
  I=\frac{1}{2}\ln(2\pi)=\ln\sqrt{2\pi}.
\]

(2) 显然
\[
  J=\int_{0}^{1}\sin\pi x\ln\Gamma(x)\,\mathrm{d}x\xlongequal{x=1-t}\int_{1}^{0}\sin\pi(1-t)\ln\Gamma(1-t)(-\mathrm{d}t)=\int_{0}^{1}\sin\pi t\ln\Gamma(1-t)\,\mathrm{d}t,
\]
故
\[
  \begin{aligned}
    2J&=\int_{0}^{1}\sin\pi t[\ln\Gamma(t)+\ln\Gamma(1-t)]\,\mathrm{d}t=\int_{0}^{1}\sin\pi t\ln[\Gamma(t)\Gamma(1-t)]\,\mathrm{d}t\\
    &\xlongequal{\text{余元公式}}\int_{0}^{1}\sin\pi x[\ln\pi-\ln\sin\pi x]\,\mathrm{d}x=\ln\pi\int_{0}^{1}\sin\pi x\,\mathrm{d}x-\int_{0}^{1}\sin\pi x\ln\sin\pi x\,\mathrm{d}x\\
    &\xlongequal{t=\pi x}\frac{2}{\pi}\ln\pi-\frac{1}{\pi}\int_{0}^{\pi}\sin t\ln\sin t\,\mathrm{d}t=\frac{2}{\pi}\ln\pi-\frac{2}{\pi}\int_{0}^{\frac{\pi}{2}}\sin x\ln\sin x\,\mathrm{d}x\\
    &=\frac{2}{\pi}\ln\pi-\frac{4}{\pi}\int_{0}^{\frac{\pi}{2}}\sin\frac{x}{2}\cos\frac{x}{2}\ln\sin x\,\mathrm{d}x\\
    &=\frac{2}{\pi}\ln\pi-\frac{4}{\pi}\left(-\frac{1}{2}+\frac{1}{2}\ln 2\right)=\frac{2}{\pi}\ln\pi+\frac{2}{\pi}-\frac{2}{\pi}\ln 2.
  \end{aligned}
\]
所以
\[
  J=\int_{0}^{1}\sin\pi x\ln\Gamma(x)\,\mathrm{d}x=\frac{1}{\pi}\left(1+\ln\frac{\pi}{2}\right).
\]
剩下需证的是
\[
  \begin{aligned}
    \int_{0}^{\frac{\pi}{2}}\sin\frac{x}{2}\cos\frac{x}{2}\ln\sin x\,\mathrm{d}x&=\int_{0}^{\frac{\pi}{2}}\ln\sin x\,\mathrm{d}\sin^{2}\frac{x}{2}\\
    &\xlongequal{\text{分部积分}}\left.\sin^{2}\frac{x}{2}\ln\sin x\right|_{0}^{\frac{\pi}{2}}-\int_{0}^{\frac{\pi}{2}}\sin^{2}\frac{x}{2}\frac{\cos x}{\sin x}\,\mathrm{d}x\\
    &=-\frac{1}{2}\int_{0}^{\frac{\pi}{2}}\frac{\sin\frac{x}{2}\cos x}{\cos\frac{x}{2}}\,\mathrm{d}x=-\frac{1}{2}\int_{0}^{\frac{\pi}{2}}\frac{\sin\frac{x}{2}}{\cos\frac{x}{2}}\left(2\cos^{2}\frac{x}{2}-1\right)\mathrm{d}x\\
    &\xlongequal{t=\cos\frac{x}{2}}\int_{1}^{\frac{\sqrt{2}}{2}}\frac{2t^{2}-1}{t}\,\mathrm{d}t=\left.(t^{2}-\ln t)\right|_{1}^{\frac{\sqrt{2}}{2}}\\
    &=\left(\frac{1}{2}-\ln\frac{1}{\sqrt{2}}\right)-(1-0)=-\frac{1}{2}+\frac{1}{2}\ln 2.
  \end{aligned}
\]
\end{proof}

\begin{example}\label{ex:15.4.7}
证明：级数 $\displaystyle\sum_{n=1}^{\infty}\frac{1}{n\mathrm{C}_{2n}^{n}}=\frac{\pi}{3\sqrt{3}}$.
\end{example}
\begin{proof}
\[
  \begin{aligned}
    \sum_{n=1}^{\infty}\frac{1}{n\mathrm{C}_{2n}^{n}}&=\sum_{n=1}^{\infty}\frac{1}{n}\frac{n!n!}{(2n)!}=\sum_{n=1}^{\infty}\frac{(n-1)!n!}{(2n)!}\\
    &=\sum_{n=1}^{\infty}\frac{\Gamma(n)\Gamma(n+1)}{\Gamma(2n+1)}\xlongequal{\text{定理~\ref{thm:15.4.5}}}\sum_{n=1}^{\infty}\mathrm{B}(n,n+1)\\
    &=\sum_{n=1}^{\infty}\mathrm{B}(n+1,n)=\sum_{n=1}^{\infty}\int_{0}^{1}t^{n}(1-t)^{n-1}\,\mathrm{d}t\\
    &=\int_{0}^{1}\sum_{n=1}^{\infty}t^{n}(1-t)^{n-1}\,\mathrm{d}t=\int_{0}^{1}t\sum_{n=1}^{\infty}[t(1-t)]^{n-1}\,\mathrm{d}t\\
    &=\int_{0}^{1}\frac{t}{1-t(1-t)}\,\mathrm{d}t=\int_{0}^{1}\frac{t}{t^{2}-t+1}\,\mathrm{d}t\\
    &=\frac{1}{2}\int_{0}^{1}\frac{(2t-1)+1}{t^{2}-t+1}\,\mathrm{d}t=\frac{1}{2}\left[\int_{0}^{1}\frac{\mathrm{d}(t^{2}-t+1)}{t^{2}-t+1}+\int_{0}^{1}\frac{\mathrm{d}t}{t^{2}-t+1}\right]\\
    &=\frac{1}{2}\left[\left.\ln(t^{2}-t+1)\right|_{0}^{1}+\left.\frac{2}{\sqrt{3}}\arctan\frac{t-\frac{1}{2}}{\frac{\sqrt{3}}{2}}\right|_{0}^{1}\right]\\
    &=\frac{1}{\sqrt{3}}\left[\arctan\frac{1}{\sqrt{3}}-\left(-\arctan\frac{1}{\sqrt{3}}\right)\right]=\frac{\pi}{3\sqrt{3}}.
  \end{aligned}
\]
\end{proof}

\begin{example}[$\Gamma$ 函数的延拓与其图像]\label{ex:15.4.8}
因为 $\forall s>0$, $\Gamma(s)>0$ 与
\[
  \Gamma''(s)=\int_{0}^{+\infty}x^{s-1}(\ln x)^{2}\mathrm{e}^{-x}\,\mathrm{d}x>0,
\]
所以 $\Gamma(s)$ 的图形位于 $x$ 轴的上方, 且是向下凸的. 由于 $\Gamma(1)=\Gamma(2)=1$, 故 $\Gamma'(s)$ 在 $s\in(0,+\infty)$ 上存在唯一的极小值点 $s_{0}\in(1,2)$. 又由 $\Gamma''(s)>0$ 知 $\Gamma'(s)$ 严格单调增加. 根据 Fermat 定理知, $\Gamma'(s_{0})=0$. 由此推得
\[
  \Gamma'(s)
  \begin{cases}
    <0, & 0<s<s_{0},\ \Gamma(s)\text{ 严格单调减少},\\
    >0, & s>s_{0},\ \Gamma(s)\text{ 严格单调增加}.
\end{cases}
\]
从 $\Gamma(s)=\frac{s\Gamma(s)}{s}=\frac{\Gamma(s+1)}{s}\ (s>0)$ 及 $\lim_{s\to 0^{+}}\Gamma(s+1)=\Gamma(1)=1$, 有
\[
  \lim_{s\to 0^{+}}\Gamma(s)=\lim_{s\to 0^{+}}\frac{\Gamma(s+1)}{s}=+\infty.
\]
由 $\Gamma(n+1)=n!$ 及 $\Gamma(s)$ 在 $(s_{0},+\infty)$ 上严格单调增加可推得
\[
  \lim_{s\to+\infty}\Gamma(s)=+\infty.
\]
综上所述, $\Gamma$ 函数的图像如图~\ref{fig:15.4.1} 中 $s>0$ 部分所示.

为延拓 $\Gamma(s)$, 我们改写公式 $\Gamma(s+1)=s\Gamma(s)\ (s>0)$ 为
\[
  \Gamma(s)=\frac{\Gamma(s+1)}{s}.
\]
当 $s\in(-1,0)$ 时, 上式右边有意义, 于是可应用它定义左边函数 $\Gamma(s)$ 在 $(-1,0)$ 内的值, 并且可推得 $\Gamma(s)<0$. 用同样的方法, 利用 $\Gamma(s)$ 已在 $(-1,0)$ 内有定义这一事实, $\Gamma(s)=\frac{\Gamma(s+1)}{s}$ 又可定义 $\Gamma(s)$ 在 $(-2,-1)$ 内的值, 而且这时 $\Gamma(s)>0$. 依次下去可将 $\Gamma(s)$ 延拓到 $\mathbb{R}\setminus\{0,-1,-2,\cdots\}$, 其图像如图~\ref{fig:15.4.1} 所示.
\end{example}

\begin{figure}[htbp]
\centering
\begin{tikzpicture}[>=latex, scale=1.12]
  \draw[->] (-3.8, 0) -- (4.2, 0) node[below] {$x$};
  \draw[->] (0, -2.8) -- (0, 3.8) node[left] {$y$};
  
  \draw[thin] (-1, -2.5) -- (-1, 3.2);
  \draw[thin] (-2, -2.5) -- (-2, 3.2);
  \draw[thin] (-3, -2.5) -- (-3, 3.2);
  
  \draw (-3, 0) ++(0, 2pt) -- ++(0, -4pt) node[below left=-1pt] {$-3$};
  \draw (-2, 0) ++(0, 2pt) -- ++(0, -4pt) node[below left=-1pt] {$-2$};
  \draw (-1, 0) ++(0, 2pt) -- ++(0, -4pt) node[below left=-1pt] {$-1$};
  \node[below right] at (0, 0) {$0$};
  \draw (1.46, 0) ++(0, 2pt) -- ++(0, -4pt) node[below] {$s_0$};
  \draw (2.8, 0) ++(0, 2pt) -- ++(0, -4pt);
  
  % 各分支曲线
  \draw[thick] plot[smooth, tension=0.7] coordinates {(-2.9, -2.2) (-2.8, -1.5) (-2.5, -1.0) (-2.2, -1.5) (-2.1, -2.2)};
  \draw[thick] plot[smooth, tension=0.7] coordinates {(-1.9, 2.5) (-1.8, 1.8) (-1.5, 1.3) (-1.2, 1.8) (-1.1, 2.5)};
  \draw[thick] plot[smooth, tension=0.7] coordinates {(-0.9, -2.5) (-0.8, -1.8) (-0.5, -1.4) (-0.2, -1.8) (-0.1, -2.5)};
  \draw[thick] plot[smooth, tension=0.8] coordinates {(0.15, 3.2) (0.35, 1.8) (0.8, 1.1) (1.46, 0.885) (2.2, 1.1) (2.8, 1.8) (3.3, 3.2)};
  
  \draw[dashed] (1.46, 0) -- (1.46, 0.885);
\end{tikzpicture}
\caption{$\Gamma$ 函数的延拓与其图像}\label{fig:15.4.1}
\end{figure}

\section*{练习题 15.4}
\begin{enumerate}
  \item 证明：$\Gamma(s)$ 可表示为
  \begin{enumerate}
    \item $\displaystyle\Gamma(s)=2\int_{0}^{+\infty}x^{2s-1}\mathrm{e}^{-x^{2}}\,\mathrm{d}x,\quad s>0$;
    \item $\displaystyle\Gamma(s)=a^{s}\int_{0}^{+\infty}x^{s-1}\mathrm{e}^{-ax}\,\mathrm{d}x,\quad s>0, a>0$.
  \end{enumerate}
  \item 证明：$\mathrm{B}(p,q)$ 可表示为
  \begin{enumerate}
    \item $\displaystyle\mathrm{B}(p,q)=\int_{0}^{+\infty}\frac{t^{q-1}}{(1+t)^{p+q}}\,\mathrm{d}t,\quad p>0, q>0$;
    \item $\displaystyle\mathrm{B}(p,q)=2\int_{0}^{\frac{\pi}{2}}\sin^{2p-1}\theta\cos^{2q-1}\theta\,\mathrm{d}\theta,\quad p>0, q>0$.
  \end{enumerate}
  \item 利用 $\Gamma$ 函数计算下列积分：
  \begin{enumerate}
    \item $\displaystyle\int_{0}^{1}\sqrt{x-x^{2}}\,\mathrm{d}x$;
    \item $\displaystyle\int_{0}^{+\infty}\frac{\sqrt[4]{x}}{(1+x)^{2}}\,\mathrm{d}x$;
    \item $\displaystyle\int_{0}^{+\infty}\frac{\mathrm{d}x}{1+x^{3}}$;
    \item $\displaystyle\int_{0}^{\frac{\pi}{2}}\sin^{5}x\cos^{6}x\,\mathrm{d}x$.
  \end{enumerate}
  \item 计算极限
  \[
    \lim_{n\to+\infty}\sqrt{n}\int_{0}^{1}x^{\frac{3}{2}}(1-x^{5})^{n}\,\mathrm{d}x.
\]
  \item 证明：$\ln\mathrm{B}(p,q)$ 关于变量 $p$ 为 $(0,+\infty)$ 上的凸函数.
  \item 证明：$\displaystyle\mathrm{B}\left(\frac{1}{3},\frac{1}{2}\right)=\frac{\sqrt{3}}{2\pi}\frac{\left[\Gamma\left(\frac{1}{3}\right)\right]^{3}}{\sqrt[3]{2}}$.
  \item 作变量代换 $u=\sin^{2}\frac{x}{2}$, 证明：$\displaystyle\int_{0}^{\pi}\frac{\mathrm{d}x}{\sqrt{3-\cos x}}=\frac{1}{2\sqrt{2}}\mathrm{B}\left(\frac{1}{4},\frac{1}{2}\right)=\frac{1}{4\sqrt{\pi}}\Gamma^{2}\left(\frac{1}{4}\right)$.
\end{enumerate}

\section*{思考题 15.4}
\begin{enumerate}
  \item 讨论函数
  \[
    f(s,p,q)=\int_{0}^{+\infty}\frac{x^{s}}{(a+bx^{p})^{q}}\,\mathrm{d}x,\quad a>0, b>0, p>0
\]
  的定义域, 并用 $\Gamma$ 函数表示 $f(s,p,q)$.
  \item 证明：$\displaystyle\int_{0}^{1}\ln\Gamma(x)\,\mathrm{d}x=\ln\sqrt{2\pi}$.
  \item 证明：$\displaystyle\int_{0}^{1}\sin\pi x\ln\Gamma(x)\,\mathrm{d}x=\frac{1}{\pi}\left(\ln\frac{\pi}{2}+1\right)$.
  \item 证明：
  \[
    \int_{0}^{+\infty}\mathrm{e}^{-x^{4}}\,\mathrm{d}x\int_{0}^{+\infty}x^{2}\mathrm{e}^{-x^{4}}\,\mathrm{d}x=\frac{\pi}{8\sqrt{2}}.
\]
  \item 证明：$\displaystyle\Gamma\left(\frac{1}{n}\right)\Gamma\left(\frac{2}{n}\right)\cdots\Gamma\left(\frac{n-1}{n}\right)=\frac{(2\pi)^{\frac{n-1}{2}}}{\sqrt{n}}$.
  \item 设 $p,q,m>0$, 证明：
  \[
    \int_{0}^{1}x^{p-1}(1-x^{m})^{q-1}\,\mathrm{d}x=\frac{1}{m}\mathrm{B}\left(\frac{p}{m},q\right)=\frac{1}{m}\frac{\Gamma\left(\frac{p}{m}\right)\Gamma(q)}{\Gamma\left(\frac{p}{m}+q\right)},
\]
  并由此推得
  \[
    \int_{0}^{1}\frac{\mathrm{d}x}{\sqrt{1-x^{4}}}\int_{0}^{1}\frac{x^{2}\,\mathrm{d}x}{\sqrt{1-x^{4}}}=\frac{\pi}{4}.
\]
  \item 证明：
  \begin{enumerate}
    \item $\displaystyle\int_{0}^{+\infty}\mathrm{e}^{-x^{n}}\,\mathrm{d}x=\Gamma\left(\frac{1}{n}+1\right)$;
    \item $\displaystyle\lim_{n\to+\infty}\int_{0}^{+\infty}\mathrm{e}^{-x^{n}}\,\mathrm{d}x=1$.
  \end{enumerate}
  \item 设 $0\leqslant h<1$, 正整数 $n\geqslant 3$. 作变量代换 $t=hu$, 证明：
  \[
    \int_{0}^{h}(1-t^{2})^{\frac{n-3}{2}}\,\mathrm{d}t\geqslant h\int_{0}^{1}(1-u^{2})^{\frac{n-3}{2}}\,\mathrm{d}u=\frac{\sqrt{\pi}}{2}\frac{\Gamma\left(\frac{n-1}{2}\right)}{\Gamma\left(\frac{n}{2}\right)}h.
\]
  \item 设 $n>0$, 证明：$\displaystyle\int_{0}^{1}\frac{\mathrm{d}x}{\sqrt[n]{1-x^{n}}}=\frac{\pi}{n\sin\frac{\pi}{n}}$.
  \item 设 $n>m>0$, 证明：$\displaystyle\int_{0}^{+\infty}\frac{x^{m-1}}{1+x^{n}}\,\mathrm{d}x=\frac{\pi}{n\sin\frac{m\pi}{n}}$.
  \item 设 $0<p<1$, 应用余元公式证明：
  \[
    \int_{0}^{+\infty}\frac{x^{p-1}\ln x}{1+x}\,\mathrm{d}x=\int_{0}^{+\infty}\frac{\partial}{\partial p}\left(\frac{x^{p-1}}{1+x}\right)\mathrm{d}x=-\frac{\pi^{2}\cos p\pi}{\sin^{2}p\pi}.
\]
  \item 利用 Stirling 公式
  \[
    n!=\sqrt{2\pi n}n^{n}\mathrm{e}^{-n+\frac{\theta_{n}}{12n}},\quad 0<\theta_{n}<1
\]
  求下列极限：
  \begin{enumerate}
    \item $\displaystyle\lim_{n\to+\infty}\sqrt[n^{2}]{n!}$;
    \item $\displaystyle\lim_{n\to+\infty}\frac{n}{\sqrt[n]{n!}}$;
    \item $\displaystyle\lim_{n\to+\infty}\frac{n}{\sqrt[n]{(2n-1)!!}}$;
    \item $\displaystyle\lim_{n\to+\infty}\frac{\ln n!}{\ln n^{n}}$.
  \end{enumerate}
\end{enumerate}

% ==================================================================
%  徐森林、金亚东、薛春华《数学分析》（第三册）
%  复习题 15
%  源：原书印刷页 241-243（PDF 物理页 247-249）
%  底本：PaddleOCR-VL-1.6，见 chapters/_ch15_15_4_raw.md 末尾
%  约定：只用 exblock，不设 section（故不入目录）
%  全 13 题题面照录
% ==================================================================

\begin{exblock}{复习题 15}
\begin{enumerate}
  \item 证明: 对任何实数 $u$, 有
  \[
    \frac{1}{2\pi}\int_{0}^{2\pi}\mathrm{e}^{u\cos x}\cos(u\sin x)\,\mathrm{d}x=1.
\]

  \item 设 $D=\{(x,y)\mid 1\leqslant x^{2}+y^{2}\leqslant 4\}$, 计算积分
  \[
    \frac{1}{2\pi}\int_{0}^{+\infty}\left[\iint\limits_{D}\frac{\sin(t\sqrt{x^{2}+y^{2}})}{\sqrt{x^{2}+y^{2}}}\,\mathrm{d}x\mathrm{d}y\right]\mathrm{d}t.
\]

  \item 设函数 $f(x,u)$ 与 $\displaystyle\frac{\partial f(x,u)}{\partial u}$ 都在 $[a,+\infty)\times[\alpha,\beta]$ 上连续, 积分 $\displaystyle\int_{a}^{+\infty}\frac{\partial f(x,u)}{\partial u}\,\mathrm{d}x$ 在 $[\alpha,\beta]$ 上一致收敛, $\displaystyle\int_{a}^{+\infty}f(x,u)\,\mathrm{d}x$ 在某点 $u_{0}\in[\alpha,\beta]$ 收敛. 证明: 积分 $\displaystyle\int_{a}^{+\infty}f(x,u)\,\mathrm{d}x$ 在 $[\alpha,\beta]$ 上一致收敛于
  \[
    \varphi(u)=\int_{a}^{+\infty}f(x,u)\,\mathrm{d}x,
\]
  且
  \[
    \varphi'(u)=\frac{\mathrm{d}}{\mathrm{d}u}\int_{a}^{+\infty}f(x,u)\,\mathrm{d}x=\int_{a}^{+\infty}\frac{\partial f(x,u)}{\partial u}\,\mathrm{d}x.
\]

  \item 证明:
  \[
    \sum_{k=0}^{n}\mathrm{C}_{n}^{k}\frac{(-1)^{k}}{m+k+1}=\frac{n!m!}{(n+m+1)!}.
\]

  \item 设 $\displaystyle f(x)=\left(\int_{0}^{x}\mathrm{e}^{-t^{2}}\,\mathrm{d}t\right)^{2}$, $\displaystyle g(x)=\int_{0}^{1}\frac{\mathrm{e}^{-x^{2}(1+t^{2})}}{1+t^{2}}\,\mathrm{d}t$, 证明:
  \[
    f(x)+g(x)=\frac{\pi}{4}.
\]
  并由此推得 $\displaystyle\int_{0}^{+\infty}\mathrm{e}^{-x^{2}}\,\mathrm{d}x=\frac{\sqrt{\pi}}{2}$.

  \item 利用公式
  \[
    \Gamma(x)=\lim_{n\to+\infty}\frac{n^{x}n!}{x(x+1)\cdots(x+n)},
\]
  证明: 对任意实数 $x$ 有
  \[
    \sin x=x\prod_{n=1}^{\infty}\left(1-\frac{x^{2}}{n^{2}\pi^{2}}\right).
\]

  \item 按照下列步骤, 给出公式
  \[
    \mathrm{B}(p,q)=\frac{\Gamma(p)\Gamma(q)}{\Gamma(p+q)}
\]
  的另一个证明:

  (1)
  \[
    \Gamma(p)=2\int_{0}^{+\infty}u^{2p-1}\mathrm{e}^{-u^{2}}\,\mathrm{d}u.
\]

  (2)
  \[
    \Gamma(p)\Gamma(q)=\lim_{A\to+\infty}4\iint\limits_{[0,A]^{2}}u^{2p-1}v^{2q-1}\mathrm{e}^{-(u^{2}+v^{2})}\,\mathrm{d}u\mathrm{d}v.
\]

  (3)
  \[
    \begin{aligned}
      \lim_{A\to+\infty}\iint\limits_{D(A)}u^{2p-1}v^{2q-1}\mathrm{e}^{-(u^{2}+v^{2})}\,\mathrm{d}u\mathrm{d}v&=\frac{1}{4}\mathrm{B}(p,q)\Gamma(p+q),\\
      \lim_{A\to+\infty}\iint\limits_{D(\sqrt{2}A)}u^{2p-1}v^{2q-1}\mathrm{e}^{-(u^{2}+v^{2})}\,\mathrm{d}u\mathrm{d}v&=\frac{1}{4}\mathrm{B}(p,q)\Gamma(p+q),
    \end{aligned}
\]
  其中 $D(R)=\left\{(r,\theta)\mid 0\leqslant r\leqslant R, 0\leqslant\theta\leqslant\frac{\pi}{2}\right\}$.

  (4)
  \[
    \mathrm{B}(p,q)=\frac{\Gamma(p)\Gamma(q)}{\Gamma(p+q)}.
\]

  \item (1) 用两种不同方法证明:
  \[
    I(\alpha)=\int_{0}^{+\infty}\mathrm{e}^{-x^{2}}\cos(2\alpha x)\,\mathrm{d}x=\frac{\sqrt{\pi}}{2}\mathrm{e}^{-\alpha^{2}};
\]

  (2) 证明:
  \[
    J(\alpha)=\int_{0}^{+\infty}\mathrm{e}^{-x^{2}}\sin(2\alpha x)\,\mathrm{d}x=\mathrm{e}^{-\alpha^{2}}\int_{0}^{\alpha}\mathrm{e}^{-x^{2}}\,\mathrm{d}x.
\]

  \item (1) 证明:
  \[
    I(\alpha)=\int_{0}^{1}(1-x^{\alpha})^{-\frac{1}{\alpha}}\,\mathrm{d}x
\]
  当 $\alpha>1$ 时收敛; 当 $\alpha\leqslant 1$ 时发散.

  (2) 证明:
  \[
    I(\alpha)=\frac{\pi}{\alpha\sin\frac{\pi}{\alpha}},\quad\alpha>1.
\]

  \item 设 $\displaystyle f(x)=\int_{0}^{\frac{\pi}{2}}\ln(1-x^{2}\cos^{2}\theta)\,\mathrm{d}\theta$, $|x|<1$.

  (1) 用积分号下对 $x$ 求导, 然后积分得
  \[
    f(x)=\pi\ln\frac{1+\sqrt{1-x^{2}}}{2},\quad |x|<1;
\]

  (2) 对 $\displaystyle\ln(1-x^{2}\cos^{2}\theta)=-\sum_{n=1}^{\infty}\frac{x^{2n}}{n}\cos^{2n}\theta$ 以及幂级数的逐项积分证明:
  \[
    f(x)=-\frac{\pi}{2}\sum_{n=1}^{\infty}\frac{(2n-1)!!}{n\cdot(2n)!!}x^{2n};
\]

  (3) 根据 (1) 与 (2) 给出 $\ln(1+\sqrt{1-x^{2}})$ 的 Maclaurin 展开式, 以及 $f^{(k)}(0)$, $k=0,1,2,\cdots$.

  \item 用下列各方法证明:
  \[
    \int_{0}^{+\infty}\frac{\mathrm{e}^{-ax^{2}}-\mathrm{e}^{-bx^{2}}}{x^{2}}\,\mathrm{d}x=\sqrt{\pi}(\sqrt{b}-\sqrt{a}),\quad a>0,b>0.
\]

  (1) 应用积分号下求导;

  (2) 应用积分号下求积分;

  (3) 应用分部积分 $\displaystyle\int_{0}^{+\infty}(\mathrm{e}^{-ax^{2}}-\mathrm{e}^{-bx^{2}})\,\mathrm{d}\left(-\frac{1}{x}\right)$;

  (4) 化为二重积分
  \[
    \int_{0}^{+\infty}\frac{\mathrm{e}^{-ax^{2}}-\mathrm{e}^{-bx^{2}}}{x^{2}}\,\mathrm{d}x=\int_{0}^{+\infty}\frac{\mathrm{d}x}{x^{2}}\int_{ax^{2}}^{bx^{2}}\mathrm{e}^{-y}\,\mathrm{d}y=\iint\limits_{D}\frac{\mathrm{e}^{-y}}{x^{2}}\,\mathrm{d}x\mathrm{d}y,
\]
  其中 $D=\left\{(x,y)\mid 0\leqslant x<+\infty, ax^{2}\leqslant y\leqslant bx^{2}\right\}$;

  (5) 将
  \[
    \frac{\mathrm{e}^{-ax^{2}}-\mathrm{e}^{-bx^{2}}}{x^{2}}=\sum_{k=0}^{\infty}\frac{1}{(k+1)!}\mathrm{e}^{-bx^{2}}(b-a)^{k+1}x^{2k}
\]
  两边积分, 并由
  \[
    \int_{0}^{+\infty}\mathrm{e}^{-bx^{2}}\,\mathrm{d}x=\frac{1}{\sqrt{b}}\frac{\sqrt{\pi}}{2},\quad\int_{0}^{+\infty}\mathrm{e}^{-bx^{2}}x^{2k}\,\mathrm{d}x=\frac{\sqrt{\pi}}{2}\frac{1}{2}\frac{3}{2}\cdots\frac{2k-1}{2}\frac{1}{b^{\frac{2k+1}{2}}}
\]
  代入推得结论.

  \item 作变量代换 $\displaystyle t=\tan\frac{\varphi}{2}$ 与 $\displaystyle\sqrt{\frac{1-k}{1+k}}t=\tan\frac{\theta}{2}$, 证明:
  \[
    \begin{aligned}
      \int_{0}^{\pi}\left(\frac{\sin\varphi}{1+\cos\varphi}\right)^{\alpha-1}\frac{\mathrm{d}\varphi}{1+k\cos\varphi}&=\frac{1}{1+k}\left(\sqrt{\frac{1+k}{1-k}}\right)^{\alpha}\int_{0}^{\pi}\tan^{\alpha-1}\frac{\theta}{2}\,\mathrm{d}\theta\\
      &=\frac{1}{1+k}\left(\sqrt{\frac{1+k}{1-k}}\right)^{\alpha}\frac{\pi}{\sin\frac{\alpha}{2}\pi}.
    \end{aligned}
\]

  \item 设 $0<p,q<1$, 用下面两种方法证明:
  \[
    I(p,q)=\int_{0}^{+\infty}\frac{x^{p-1}-x^{q-1}}{(1+x)\ln x}\,\mathrm{d}x=\ln\left|\frac{\tan\frac{p\pi}{2}}{\tan\frac{q\pi}{2}}\right|.
\]

  (1) 利用 $\displaystyle\frac{\partial I}{\partial p}=\frac{\pi}{\sin p\pi}$ 及积分号下求导;

  (2) 由 $\displaystyle\left.\frac{x^{u-1}}{\ln x}\right|_{u=q}^{u=p}=\int_{q}^{p}x^{u-1}\,\mathrm{d}u$ 及积分号下求积分.
\end{enumerate}
\end{exblock}

